% Les amines — Chimie Tle TI — Cours signé OnBuch+
% Programme officiel MINESEC, Terminale TI. Compiler avec : tectonic L03-les-amines.tex
\newcommand{\DOCMATIERE}{Chimie}
\newcommand{\DOCNIVEAU}{Tle TI}
\newcommand{\DOCMODULE}{Module 1 · Chimie organique}
\newcommand{\DOCLECON}{03}
\newcommand{\DOCTITRE}{Les amines}
% OnBuch+ — préambule « cours signés » ENRICHI (compilé avec tectonic / XeLaTeX)
% Système visuel « Pop » d'OnBuch+ : fond crème (jamais blanc), bordures encre
% épaisses, ombres « dures » décalées, titres Archivo Black en capitales.
% Version enrichie pour les cours de chimie : chimie (mhchem, chemfig),
% graphes (pgfplots), boîtes pédagogiques supplémentaires (définition,
% propriété, méthode, attention, le savais-tu, exercice résolu, exercice,
% TP, bilan), page de garde refaite et signature OnBuch+ ancrée en bas.
%
% Macros attendues dans main.tex AVANT \input{preamble} :
%   \DOCMATIERE  \DOCNIVEAU  \DOCMODULE  \DOCLECON (numéro)  \DOCTITRE
\documentclass[11pt]{article}

\usepackage{fontspec}
\usepackage[french]{babel}
\usepackage[a4paper,margin=2cm,top=3.4cm,bottom=2.8cm,headheight=1.4cm,footskip=1.3cm]{geometry}
\usepackage{amsmath,amssymb,amsfonts}
\usepackage[table]{xcolor} % [table] : fournit aussi \rowcolors (alternance de lignes)
\usepackage{pagecolor}
\usepackage{tikz}
\usetikzlibrary{arrows.meta,positioning,calc,shapes.geometric,shapes.misc,shapes.arrows,decorations.pathmorphing,decorations.markings,patterns,shadows,mindmap,trees,fit,backgrounds,matrix,intersections}
\usepackage{pgfplots}
\pgfplotsset{compat=1.18}
\usepackage[version=4]{mhchem}
\usepackage{chemfig}
\usepackage{enumitem}
\usepackage{fancyhdr}
% [most] charge toutes les bibliothèques tcolorbox (listings, minted, raster,
% documentation...) et gonfle inutilement la mémoire TeX sur un document long
% (~300 boîtes) : on ne charge que ce qui est réellement utilisé.
\usepackage[skins,breakable]{tcolorbox}
\usepackage{graphicx}
\usepackage{colortbl}
\usepackage{booktabs}
\usepackage{tabularx}
\usepackage{array}
\usepackage{multicol}
\usepackage{siunitx}
% parse-numbers=false : accepte aussi \SI{2,5 \times 10^{-3}}{...}
\sisetup{locale=FR,output-decimal-marker={,},group-minimum-digits=4,parse-numbers=false}
\DeclareSIUnit\ohm{\ensuremath{\symup{\Omega}}} % U+2126 absent de la police
\usepackage{qrcode}
% \degree : commande fréquemment utilisée par les modèles pour un angle ou une
% température en dehors de \SI{}{} (non fournie par les paquets chargés).
\providecommand{\degree}{^{\circ}}
\usepackage{lastpage}
\usepackage{hyperref}
\hypersetup{hidelinks}

% ============================================================
% Palette « Pop » OnBuch+ — mêmes hex que la classe P (plus_ui.dart)
% ============================================================
\definecolor{poporange}{HTML}{F59321}
\definecolor{popdark}{HTML}{E07A0C}
\definecolor{popcream}{HTML}{FFF6E8}
\definecolor{popink}{HTML}{1C1714}
\definecolor{poppurple}{HTML}{7A5AE0}
\definecolor{popgreen}{HTML}{2E9E6A}
\definecolor{poppink}{HTML}{E0457B}
\definecolor{popblue}{HTML}{2F7DE1}
\definecolor{popgold}{HTML}{C98A12}
\definecolor{popmuted}{HTML}{8A7F76}
\definecolor{popline}{HTML}{EADFD2}
\definecolor{popgray}{HTML}{8A7F76}
\colorlet{popgrey}{popgray}
% Alias tolérants (le modèle nomme parfois les couleurs différemment)
\colorlet{popwhite}{white}
\colorlet{popblack}{popink}
\colorlet{popred}{poppink}
\colorlet{popyellow}{popgold}
% Teintes claires pour fonds de tableaux / zones
\colorlet{popblueL}{popblue!10!white}
\colorlet{poporangeL}{poporange!14!white}
\colorlet{popgreenL}{popgreen!12!white}
\colorlet{poppurpleL}{poppurple!12!white}
\colorlet{poppinkL}{poppink!12!white}
\colorlet{popgoldL}{popgold!14!white}

\pagecolor{popcream}

% ============================================================
% Typographie
% ============================================================
\defaultfontfeatures{Ligatures=TeX}
\setmainfont{PlusJakartaSans-Medium.ttf}[Path=../../../fonts/,BoldFont=PlusJakartaSans-Bold.ttf]
% Maths en sans-serif (Fira Math) avec chiffres et lettres droites de Plus
% Jakarta, pour que les formules se fondent dans le texte.
\usepackage[math-style=ISO,bold-style=ISO]{unicode-math}
\setmathfont{FiraMath-Regular.otf}[Path=../../../fonts/]
\setmathfont{PlusJakartaSans-Medium.ttf}[Path=../../../fonts/,range={up/num,up/{Latin,latin},\mathup}]
\usepackage{icomma} % virgule décimale sans espace en mode math
% Caractères absents de la police de texte (sinon carré vide) : rendus par
% leur équivalent mathématique, en texte comme en formule.
\usepackage{newunicodechar}
\newunicodechar{α}{\ensuremath{\alpha}}
\newunicodechar{Α}{\ensuremath{\alpha}}
\newunicodechar{β}{\ensuremath{\beta}}
\newunicodechar{δ}{\ensuremath{\delta}}
\newunicodechar{✓}{\popcheck{popgreen}}
\newfontfamily\popdisplay{ArchivoBlack-Regular.ttf}[Path=../../../fonts/]
\newfontfamily\poplogo{SpaceGrotesk-Bold.ttf}[Path=../../../fonts/]
\newfontfamily\popsemi{PlusJakartaSans-SemiBold.ttf}[Path=../../../fonts/]
\newfontfamily\popxbold{PlusJakartaSans-ExtraBold.ttf}[Path=../../../fonts/]
\newfontfamily\popxboldls{PlusJakartaSans-ExtraBold.ttf}[Path=../../../fonts/,LetterSpace=3.0]

\setlist[enumerate]{leftmargin=1.7em,itemsep=5pt,topsep=4pt}
\setlist[itemize]{leftmargin=1.7em,itemsep=4pt,topsep=4pt,label={\color{poporange}\textbullet}}
\setlength{\parindent}{0pt}
\setlength{\parskip}{5pt}
\renewcommand{\arraystretch}{1.25}

% chemfig : liaisons un peu plus épaisses, encre Pop
\setchemfig{atom sep=2em,bond style={line width=0.9pt,popink},cram width=3pt}

% pgfplots : style Pop par défaut
\pgfplotsset{
  every axis/.append style={axis line style={popink,line width=1pt},tick style={popink},
    grid style={popline,line width=0.5pt},label style={font=\small},tick label style={font=\footnotesize}},
  popcurve/.style={poporange,line width=1.8pt,no markers,smooth},
}

% ============================================================
% Pill / squiggle
% ============================================================
\newcommand{\poppill}[3]{%
  \tikz[baseline=-0.55ex]\node[draw=popink,line width=1.2pt,rounded corners=2.6pt,%
    fill=#2,inner xsep=6.5pt,inner ysep=3pt]{%
    \popxboldls\fontsize{7}{7}\selectfont\color{#3}\MakeUppercase{#1}};%
}
\newcommand{\popsquiggle}[1]{%
  \tikz[baseline=-0.4ex]{
    \draw[#1,line width=2.6pt,line cap=round]
      plot [smooth] coordinates {(0,0.09) (0.34,0.01) (0.68,0.21) (1.02,0.01) (1.36,0.10)};
  }%
}
% Mot-clé surligné (marqueur orange sous le texte)
\newcommand{\cle}[1]{{\popxbold\color{popdark}#1}}

% ============================================================
% En-tête / pied
% ============================================================
\pagestyle{fancy}
\fancyhf{}
\renewcommand{\headrulewidth}{0pt}
\renewcommand{\footrulewidth}{0pt}
\lhead{\raisebox{-0.15ex}{\poplogo\fontsize{13}{13}\selectfont\color{popink}OnBuch\color{poporange}+}}
\rhead{\poppill{\DOCMATIERE\ \textbullet\ \DOCNIVEAU\ \textbullet\ Leçon \DOCLECON}{white}{popink}}
\lfoot{\popxbold\fontsize{7.6}{7.6}\selectfont\color{popmuted}\MakeUppercase{Cours signé OnBuch+}\ \textbullet\ {\popsemi\DOCTITRE}}
\rfoot{\tikz[baseline=-0.6ex]\node[rounded corners=8.5pt,fill=popink,minimum height=17pt,minimum width=17pt,inner xsep=5pt,inner sep=0]{\popxbold\fontsize{8}{8}\selectfont\color{popcream}\thepage\,/\,\pageref*{LastPage}};}

% ============================================================
% Titres
% ============================================================
\newcommand{\chaptitre}[2]{%
  \begin{tcolorbox}[enhanced,colback=poporange,colframe=popink,boxrule=2.6pt,arc=11pt,
      drop shadow=popink,left=15pt,right=15pt,top=12pt,bottom=11pt,before skip=3pt,after skip=18pt]
    \poppill{#2}{popink}{popcream}\par\vspace{9pt}
    {\raggedright\hyphenpenalty=10000\exhyphenpenalty=10000\popdisplay\fontsize{22}{24}\selectfont\color{popcream}\MakeUppercase{#1}\par}
  \end{tcolorbox}
}
% Section numérotée : pastille orange avec le numéro + titre Archivo
\newcounter{popsec}
\newcommand{\coursec}[1]{%
  \stepcounter{popsec}\setcounter{popsub}{0}%
  \par\vspace{16pt}\noindent
  \begin{minipage}{\linewidth}
  \tikz[baseline=-0.7ex]\node[draw=popink,line width=1.6pt,fill=poporange,rounded corners=4pt,minimum size=22pt,inner sep=0]{\popdisplay\fontsize{12}{12}\selectfont\color{popcream}\thepopsec};%
  \hspace{8pt}{\popdisplay\fontsize{15}{17}\selectfont\color{popink}\MakeUppercase{#1}}\par
  \vspace{4pt}\noindent{\color{poporange}\rule{36pt}{3.6pt}}
  \end{minipage}\par\nopagebreak\vspace{8pt}
}
\newcounter{popsub}
\newcommand{\courssub}[1]{%
  \stepcounter{popsub}%
  \par\vspace{9pt}\noindent{\popxbold\fontsize{11.5}{13}\selectfont\color{popink}{\color{poporange}\thepopsec.\thepopsub}\hspace{6pt}#1}\par\nopagebreak\vspace{4pt}
}

% ============================================================
% Boîtes pédagogiques — même recette : fond blanc, bordure épaisse colorée,
% ombre dure encre, badge plein. Titre optionnel [#1].
% ============================================================
\tcbset{popbase/.style={breakable,enhanced,colback=white,boxrule=2.3pt,arc=9pt,
  shadow={3pt}{-3pt}{0pt}{popink},left=14pt,right=14pt,top=11pt,bottom=11pt,before skip=11pt,after skip=15pt,
  fontupper=\popsemi\fontsize{10.6}{14.5}\selectfont\color{popink},
  fontlower=\popsemi\fontsize{10.6}{14.5}\selectfont\color{popink}}}
% Coche dessinée (le glyphe ✓ n'existe pas dans les polices du document)
\newcommand{\popcheck}[1]{\tikz[baseline=-0.5ex]\draw[#1,line width=1.6pt,line cap=round,line join=round](-0.13,0.0)--(-0.04,-0.09)--(0.13,0.1);}
\newcommand{\popboxhead}[3]{\poppill{#1}{#2}{white}\ifx\relax#3\relax\else\hspace{6pt}{\popxbold\fontsize{10.6}{12}\selectfont\color{popink}#3}\fi\par\vspace{7pt}}

\newtcolorbox{aretenir}[1][]{popbase,colframe=popgreen,before upper={\popboxhead{À retenir}{popgreen}{#1}}}
\newtcolorbox{exemplebox}[1][]{popbase,colframe=poporange,before upper={\popboxhead{Exemple}{poporange}{#1}}}
\newtcolorbox{definition}[1][]{popbase,colframe=popblue,before upper={\popboxhead{Définition}{popblue}{#1}}}
\newtcolorbox{propriete}[1][]{popbase,colframe=poppurple,before upper={\popboxhead{Propriété}{poppurple}{#1}}}
\newtcolorbox{methode}[1][]{popbase,colframe=popgold,before upper={\popboxhead{Méthode}{popgold}{#1}}}
\newtcolorbox{attention}[1][]{popbase,colframe=poppink,before upper={\popboxhead{Attention}{poppink}{#1}}}
\newtcolorbox{experience}[1][]{popbase,colframe=popblue,colback=popblueL,before upper={\popboxhead{Expérience}{popblue}{#1}}}
% « Le savais-tu ? » : carte violette pleine (contexte, culture, Cameroun)
\newtcolorbox{savaistu}[1][]{popbase,colback=poppurple,colframe=popink,
  fontupper=\popsemi\fontsize{10.4}{14.2}\selectfont\color{white},
  before upper={\poppill{Le savais-tu\,?}{white}{poppurple}\ifx\relax#1\relax\else\hspace{6pt}{\popxbold\color{white}#1}\fi\par\vspace{7pt}}}
% Exercice résolu : énoncé en haut, \tcblower puis la solution sur fond crème
\newcounter{popexo}\newcounter{popexr}
\newtcolorbox{exoresolu}[1][]{popbase,colframe=poporange,colbacklower=popcream,
  segmentation style={popink,line width=1.2pt,dashed},
  before upper={\stepcounter{popexr}\popboxhead{Exercice résolu \thepopexr}{poporange}{#1}},
  before lower={\poppill{Solution}{popink}{popcream}\par\vspace{6pt}}}
% Exercice (non résolu en place) — la correction est dans la partie Corrigés
\newtcolorbox{exercice}[1][]{popbase,colframe=popink,
  before upper={\stepcounter{popexo}\popboxhead{Exercice \thepopexo}{popink}{#1}}}
% Corrigé
\newtcolorbox{corrige}[1][]{popbase,colframe=popgreen,colback=popgreenL,
  before upper={\popboxhead{Corrigé}{popgreen}{#1}}}
% Objectifs / prérequis / bilan
\newtcolorbox{objectifs}{popbase,colframe=popink,colback=poporangeL,
  before upper={\popboxhead{Objectifs de la leçon}{popink}{}}}
\newtcolorbox{prerequis}{popbase,colframe=popink,colback=popblueL,
  before upper={\popboxhead{Prérequis}{popblue}{}}}
\newtcolorbox{bilan}{popbase,colframe=popink,colback=white,boxrule=2.8pt,
  before upper={\popboxhead{Fiche bilan}{poporange}{L'essentiel à connaître pour l'examen}}}
% Figure encadrée avec légende
\newtcolorbox{popfigure}[1][]{enhanced,breakable=false,colback=white,colframe=popline,boxrule=1.4pt,arc=8pt,
  left=10pt,right=10pt,top=10pt,bottom=8pt,before skip=10pt,after skip=12pt,halign=center,
  lower separated=false,halign lower=center,
  fontlower=\popsemi\fontsize{9}{11.5}\selectfont\color{popmuted},
  #1}
\newcommand{\legende}[1]{\par\vspace{6pt}{\popsemi\fontsize{9}{11.5}\selectfont\color{popmuted}#1}}

% ============================================================
% Signature OnBuch+
% ============================================================
\newcommand{\poplogoinline}[1]{{\poplogo\fontsize{#1}{#1}\selectfont\color{popink}OnBuch\color{poporange}+}}
\newcommand{\signatureonbuch}{%
  \begin{tcolorbox}[enhanced,colback=popink,colframe=popink,boxrule=2.2pt,arc=10pt,
      drop shadow=poporange,left=14pt,right=14pt,top=10pt,bottom=10pt,before skip=0pt,after skip=0pt]
    \tikz[baseline=-0.7ex]\node[circle,fill=popgreen,draw=popcream,line width=1.3pt,minimum size=22pt,inner sep=0]{\popcheck{white}};%
    \hspace{9pt}\begin{minipage}[c]{0.62\linewidth}
      {\popxbold\fontsize{12}{13}\selectfont\color{popcream}Signé\ \textbullet\ {\poplogo OnBuch\color{poporange}+}}\par\vspace{2pt}
      {\raggedright\popsemi\fontsize{8}{10.5}\selectfont\color{popline}\MakeUppercase{Équipe pédagogique OnBuch+}\\\MakeUppercase{Conforme au programme officiel MINESEC}\par}
    \end{minipage}\hfill
    {\poplogo\fontsize{22}{22}\selectfont\color{popcream}OnBuch\color{poporange}+}
  \end{tcolorbox}%
}

% ============================================================
% Page de garde
% #1 titre · #2 matière · #3 niveau · #4 module · #5 n° leçon · #6 durée ·
% #7 description / objectifs (texte ou itemize)
% ============================================================
\newcommand{\pagegarde}[7]{%
  \thispagestyle{empty}
  \newgeometry{a4paper,margin=1.7cm,top=1.6cm,bottom=1.4cm}%
  \begin{tikzpicture}[remember picture,overlay]
    \fill[poporange!18] ([xshift=-3.2cm,yshift=-3.6cm]current page.north east) circle (4.6cm);
    \fill[poppurple!14] ([xshift=2.4cm,yshift=7.6cm]current page.south west) circle (3.8cm);
  \end{tikzpicture}%
  {\poplogo\fontsize{30}{30}\selectfont\color{popink}OnBuch\color{poporange}+}\hfill
  \raisebox{0.6ex}{\poppill{Cours signé \textbullet\ Premium}{popink}{popcream}}\par
  \vspace{1pt}\popsquiggle{poppurple}\par
  \vspace{8pt}
  \begin{tcolorbox}[enhanced,colback=poporange,colframe=popink,boxrule=2.8pt,arc=13pt,
      drop shadow=popink,left=18pt,right=18pt,top=12pt,bottom=13pt,after skip=10pt]
    \poppill{#2 \textbullet\ #3}{popink}{popcream}\hspace{5pt}\poppill{#4}{white}{popink}\par\vspace{8pt}
    {\popdisplay\fontsize{14}{14}\selectfont\color{popink}LEÇON #5}\par\vspace{4pt}
    {\raggedright\hyphenpenalty=10000\exhyphenpenalty=10000\popdisplay\fontsize{25}{27}\selectfont\color{popcream}\MakeUppercase{#1}\par}
    \vspace{8pt}
    {\popxbold\fontsize{9.5}{9.5}\selectfont\color{popink}\MakeUppercase{Durée conseillée : #6}}
  \end{tcolorbox}
  \begin{tcolorbox}[enhanced,colback=white,colframe=popink,boxrule=2.2pt,arc=10pt,
      drop shadow=popink,left=16pt,right=16pt,top=10pt,bottom=10pt,after skip=8pt]
    \poppill{Dans cette leçon}{poppurple}{white}\par\vspace{5pt}
    \popsemi\fontsize{9.3}{12.6}\selectfont\color{popink}#7
  \end{tcolorbox}
  \noindent\begin{minipage}[t]{0.487\linewidth}
    \begin{tcolorbox}[enhanced,colback=popgreen,colframe=popink,boxrule=2.2pt,arc=10pt,
        drop shadow=popink,left=11pt,right=11pt,top=10pt,bottom=10pt,height=3.1cm,valign=center]
      \tcbox[colback=white,colframe=popink,boxrule=1.3pt,arc=5pt,boxsep=0pt,nobeforeafter,
        left=3pt,right=3pt,top=3pt,bottom=3pt,valign=center]{\qrcode[height=1.9cm]{https://play.google.com/store/apps/details?id=com.luvvix.onbuch&hl=fr}}%
      \hspace{8pt}\begin{minipage}[c]{0.5\linewidth}
        {\popdisplay\fontsize{11}{12.5}\selectfont\color{white}\MakeUppercase{Télécharge OnBuch}}\par\vspace{4pt}
        {\raggedright\popsemi\fontsize{8.4}{11}\selectfont\color{white}Gratuit \textbullet\ Google Play\\Tuteur IA, annales, concours, résultats\par}
      \end{minipage}
    \end{tcolorbox}
  \end{minipage}\hfill
  \begin{minipage}[t]{0.487\linewidth}
    \begin{tcolorbox}[enhanced,colback=poppurple,colframe=popink,boxrule=2.2pt,arc=10pt,
        drop shadow=popink,left=11pt,right=11pt,top=10pt,bottom=10pt,height=3.1cm,valign=center]
      \tikz[baseline=-0.7ex]\node[circle,fill=white,draw=popink,line width=1.3pt,minimum size=1.6cm,inner sep=0]{\popdisplay\fontsize{24}{24}\selectfont\color{poppurple}+};%
      \hspace{8pt}\begin{minipage}[c]{0.6\linewidth}
        {\popdisplay\fontsize{11}{12.5}\selectfont\color{white}\MakeUppercase{Passe à OnBuch+}}\par\vspace{4pt}
        {\raggedright\popsemi\fontsize{8.4}{11}\selectfont\color{white}Tous les cours signés, Tuteur IA illimité, fiches PDF — onglet OnBuch+ de l'app\par}
      \end{minipage}
    \end{tcolorbox}
  \end{minipage}
  \vspace{8pt}
  \popcontenu
  \vfill
  \signatureonbuch
  \restoregeometry
  \newpage
}

% Rangée « Ce cours contient » (page de garde)
\newcommand{\poptile}[3]{% #1 couleur #2 gros texte #3 légende
  \begin{tcolorbox}[enhanced,colback=#1,colframe=popink,boxrule=2pt,arc=9pt,shadow={2.5pt}{-2.5pt}{0pt}{popink},
      left=6pt,right=6pt,top=9pt,bottom=9pt,halign=center,valign=center,height=1.75cm,width=\linewidth,nobeforeafter]
    {\popdisplay\fontsize{12.5}{13}\selectfont\color{white}\MakeUppercase{#2}}\par\vspace{3pt}
    {\centering\popsemi\fontsize{7.8}{9.5}\selectfont\color{white}#3\par}
  \end{tcolorbox}}
\newcommand{\popcontenu}{%
  \par\noindent{\popxboldls\fontsize{7.5}{7.5}\selectfont\color{popmuted}CE COURS SIGNÉ CONTIENT}\par\vspace{7pt}
  \noindent\begin{minipage}[t]{0.235\linewidth}\poptile{popblue}{Cours}{complet, illustré, pas à pas}\end{minipage}\hfill
  \begin{minipage}[t]{0.235\linewidth}\poptile{poporange}{Méthodes}{et exercices résolus}\end{minipage}\hfill
  \begin{minipage}[t]{0.235\linewidth}\poptile{poppink}{Exercices}{type examen + corrigés}\end{minipage}\hfill
  \begin{minipage}[t]{0.235\linewidth}\poptile{popgreen}{Fiche bilan}{l'essentiel à retenir}\end{minipage}\par}

% Sommaire
\newtcolorbox{sommaire}{enhanced,colback=white,colframe=popink,boxrule=2.2pt,arc=10pt,
  drop shadow=popink,left=18pt,right=18pt,top=13pt,bottom=13pt,before skip=6pt,after skip=20pt,
  before upper={\poppill{Sommaire}{poppurple}{white}\par\vspace{9pt}\popsemi\fontsize{10.6}{14.5}\selectfont\color{popink}}}

% Signature de fin de cours — poussée tout en bas de la dernière page
\newcommand{\findecours}{%
  \par\vfill\nopagebreak
  \begin{center}\popsquiggle{poporange}\end{center}\vspace{4pt}
  \signatureonbuch
}
% Glyphes absents de Fira Math : redéfinis à partir de glyphes disponibles
\AtBeginDocument{%
  \def\setminus{\mathbin{\backslash}}\let\smallsetminus\setminus
  \def\vdots{\mathord{\vbox{\baselineskip3.2pt\lineskiplimit0pt\kern4pt\hbox{.}\hbox{.}\hbox{.}}}}%
  \def\ddots{\mathinner{\mkern1mu\raise6.5pt\hbox{.}\mkern2mu\raise3.6pt\hbox{.}\mkern2mu\raise0.7pt\hbox{.}\mkern1mu}}%
  \def\triangle{\mathord{\scalebox{0.9}{$\symup{\Delta}$}}}%
  \def\nabla{\mathord{\raisebox{\depth}{\scalebox{1}[-1]{$\symup{\Delta}$}}}}%
  \def\checkmark{\popcheck{popdark}}%
  \def\star{\mathbin{\ast}}\def\bigstar{\mathord{\ast}}%
  \def\diamond{\mathbin{\scalebox{0.6}{$\Diamond$}}}%
  \def\bigcup{\mathop{\mathchoice{\raisebox{-0.3em}{\scalebox{1.7}{$\cup$}}}{\scalebox{1.25}{$\cup$}}{\cup}{\cup}}\displaylimits}%
  \def\bigcap{\mathop{\mathchoice{\raisebox{-0.3em}{\scalebox{1.7}{$\cap$}}}{\scalebox{1.25}{$\cap$}}{\cap}{\cap}}\displaylimits}%
  \def\lesseqqgtr{\mathrel{\lessgtr}}\def\bigoplus{\mathop{\oplus}\displaylimits}%
}
\providecommand{\ssi}{si et seulement si} % abréviation fréquente

%% FIGFIX-BEGIN v5 — correctifs globaux des figures (docs/onbuch-generation/tools/figfix)
\usepackage{adjustbox}\usepackage{etoolbox}\tcbuselibrary{raster}
% toute figure TikZ/pgfplots plus large que la ligne est réduite à la largeur disponible
\BeforeBeginEnvironment{tikzpicture}{\begin{adjustbox}{max width=\linewidth}}
\AfterEndEnvironment{tikzpicture}{\end{adjustbox}}
% légendes pgfplots lisibles par-dessus les courbes
\pgfplotsset{every axis/.append style={legend style={fill=white,fill opacity=0.92,text opacity=1,draw=black!15,inner sep=3pt}}}
% étiquettes d'arêtes (syntaxe edge["..."]) avec fond blanc
\tikzset{every edge quotes/.append style={fill=white,inner sep=1.2pt}}
%% FIGFIX-END
\begin{document}

\pagegarde{Les amines}{Chimie}{Tle TI}{Module 1 · Chimie organique}{03}{3 h}{Cette leçon présente les amines, dérivés de l'ammoniac : structure, classes, nomenclature, propriétés physiques et chimiques (basicité, nucléophilie). Elle aboutit aux réactions clés de synthèse pharmaceutique : réaction d'Hofmann et préparation des amides à partir des chlorures d'acyle.
\vspace{4pt}
\begin{multicols}{2}\small
\begin{itemize}
\item À la fin de cette leçon, je sais décrire la structure de la molécule d'ammoniac et en déduire celle des amines.
\item À la fin de cette leçon, je sais écrire la formule générale des amines et identifier leur classe (primaire, secondaire, tertiaire).
\item À la fin de cette leçon, je sais nommer une amine à partir de sa formule semi-développée et inversement.
\item À la fin de cette leçon, je sais expliquer l'état physique et la solubilité des amines par les liaisons hydrogène.
\item À la fin de cette leçon, je sais écrire l'équation-bilan de la réaction d'une amine avec l'eau et justifier son caractère basique.
\item À la fin de cette leçon, je sais mettre en évidence le caractère nucléophile des amines à partir du doublet non liant de l'azote.
\end{itemize}\end{multicols}}

\begin{sommaire}
\begin{multicols}{2}
\begin{enumerate}
\item De l'ammoniac aux amines : structure et classes
\item Nomenclature des amines
\item Propriétés physiques des amines
\item Caractère basique des amines
\item Caractère nucléophile et réaction d'Hofmann
\item Action sur les chlorures d'acyle : synthèse des amides
\item Travaux pratiques
\item Méthodes et exercices résolus
\item Exercices
\item Corrigés des exercices
\item Fiche bilan
\end{enumerate}
\end{multicols}
\end{sommaire}

\begin{objectifs}
\begin{itemize}
\item À la fin de cette leçon, je sais décrire la structure de la molécule d'ammoniac et en déduire celle des amines.
\item À la fin de cette leçon, je sais écrire la formule générale des amines et identifier leur classe (primaire, secondaire, tertiaire).
\item À la fin de cette leçon, je sais nommer une amine à partir de sa formule semi-développée et inversement.
\item À la fin de cette leçon, je sais expliquer l'état physique et la solubilité des amines par les liaisons hydrogène.
\item À la fin de cette leçon, je sais écrire l'équation-bilan de la réaction d'une amine avec l'eau et justifier son caractère basique.
\item À la fin de cette leçon, je sais mettre en évidence le caractère nucléophile des amines à partir du doublet non liant de l'azote.
\item À la fin de cette leçon, je sais écrire les équations-bilans des réactions d'Hofmann (amine + dérivé halogéné).
\item À la fin de cette leçon, je sais écrire l'équation-bilan de la réaction d'une amine sur un chlorure d'acyle et identifier l'amide formé.
\end{itemize}
\end{objectifs}

\begin{prerequis}
\begin{itemize}
\item Nomenclature des alcanes et des dérivés halogénés (Première)
\item Couples acide/base et réaction d'un acide avec l'eau (Première)
\item Liaison hydrogène et propriétés physiques des alcools (Première)
\item Groupes caractéristiques : alcool, acide carboxylique, chlorure d'acyle (début de Terminale)
\item Écriture et ajustement des équations-bilans en chimie organique
\end{itemize}
\end{prerequis}

\newpage

\coursec{De l'ammoniac aux amines : structure et classes}

\courssub{Situation d'accroche : de l'ammoniac des champs au sirop de la pharmacie}

Il est neuf heures du matin dans une pharmacie du quartier Nlongkak, à Yaoundé. Une jeune pharmacienne prépare un sirop antipaludéen à base de chloroquine pour une maman inquiète. À quelques rues de là, un client savoure un café arabica de l'Ouest, riche en caféine. Ces deux molécules, la chloroquine et la caféine, ont un point commun remarquable : elles contiennent un ou plusieurs atomes d'\cle{azote} directement liés à des chaînes carbonées. Ce sont des \cle{amines}.

Pourtant, l'ancêtre de toutes ces molécules est bien plus modeste : l'\cle{ammoniac} \ce{NH3}, ce gaz incolore à l'odeur piquante que l'on retrouve dans les engrais azotés distribués par la SODECOTON pour les cotoncultures du Nord, ou dans certains produits de nettoyage. Comment passe-t-on de cette petite molécule à trois atomes d'hydrogène aux molécules complexes qui soignent, stimulent ou embaument ? Pourquoi les amines sont-elles à la fois \emph{basiques} et capables d'\emph{attaquer} d'autres molécules ? Et comment les chimistes les utilisent-ils pour fabriquer des médicaments comme le paracétamol vendu dans toutes les pharmacies du pays ?

C'est tout l'objet de cette leçon. Nous commençons ici par le fondement de tout l'édifice : comprendre la structure de l'ammoniac, définir les amines et apprendre à les classer en trois grandes familles.

\courssub{Structure de la molécule d'ammoniac \ce{NH3}}

\textbf{Intuition d'abord.} Imagine une tente montée sur le sol : trois piquets plantés dans la terre et un mât central qui culmine. La molécule d'ammoniac ressemble exactement à cela : les trois atomes d'hydrogène forment la base triangulaire, et l'atome d'azote se dresse au sommet. Mais ce qui rend cette molécule vraiment spéciale, c'est un détail invisible : un \cle{doublet non liant}, une paire d'électrons « libres » portée par l'azote, comme un petit nuage de charge négative pointant vers le haut de la pyramide.

\textbf{Étude rigoureuse.} L'atome d'azote (numéro atomique $Z = 7$) possède la configuration électronique $1s^2\,2s^2\,2p^3$ : il a donc $5$ électrons sur sa couche externe. Dans \ce{NH3}, il met en commun trois de ces électrons avec trois atomes d'hydrogène pour former trois \cle{liaisons covalentes} simples \ce{N-H}. Il reste alors sur l'azote une paire d'électrons non partagée : le \cle{doublet non liant}.

L'azote est ainsi entouré de \textbf{quatre doublets} d'électrons : trois doublets liants (les liaisons \ce{N-H}) et un doublet non liant. D'après la théorie de répulsion des paires électroniques (méthode VSEPR, type $AX_3E_1$), ces quatre doublets se disposent aux sommets d'un tétraèdre. Mais comme l'un des sommets est occupé par le doublet non liant (invisible), la géométrie \emph{observable} de la molécule est \cle{pyramidale} à base triangulaire.

\begin{propriete}[Géométrie de l'ammoniac]
La molécule d'ammoniac \ce{NH3} est \cle{pyramidale} : l'atome d'azote occupe le sommet et les trois atomes d'hydrogène forment la base triangulaire. L'angle de liaison $\widehat{HNH}$ vaut environ \SI{107}{\degree}, légèrement inférieur à l'angle tétraédrique (\SI{109,5}{\degree}), car le doublet non liant, plus « volumineux », repousse davantage les liaisons \ce{N-H}. La molécule est \cle{polaire} et possède un doublet non liant disponible sur l'azote.
\end{propriete}

Ce doublet non liant est la clé de tout le chapitre : c'est lui qui rendra les amines \emph{basiques} (capables de capter un proton \ce{H+}) et \emph{nucléophiles} (capables d'attaquer les sites pauvres en électrons d'autres molécules). Retiens bien cette image : le doublet non liant est la « main libre » de l'azote.

\courssub{Passage de l'ammoniac aux amines}

\textbf{Intuition d'abord.} Reprenons notre tente. Que se passe-t-il si l'on remplace un piquet (un atome d'hydrogène) par un poteau plus gros, un groupe carboné comme le groupe méthyle \ce{-CH3} ? On obtient une nouvelle molécule, \ce{CH3-NH2}, la méthylamine. Remplaçons un deuxième piquet : \ce{(CH3)2NH}, la diméthylamine. Puis le troisième : \ce{(CH3)3N}, la triméthylamine. À chaque étape, l'azote garde son doublet non liant, sa « main libre ». C'est exactement ainsi que l'on construit la famille des amines.

\begin{definition}[Amine]
Une \cle{amine} est un composé organique qui dérive de l'ammoniac \ce{NH3} par remplacement d'\emph{un}, de \emph{deux} ou de \emph{trois} atomes d'hydrogène par des groupes carbonés (alkyles $R$ ou aryles comme le groupe phényle \ce{-C6H5}). L'atome d'azote conserve toujours son doublet non liant.
\end{definition}

La \cle{formule générale} d'une amine s'écrit $R_{(3-n)}NH_n$ avec $n \in \{0, 1, 2\}$ le nombre d'atomes d'hydrogène restant sur l'azote, c'est-à-dire :
\begin{itemize}
\item $R-NH_2$ si un seul hydrogène a été remplacé ;
\item $R_2NH$ (ou $R-NH-R'$) si deux hydrogènes ont été remplacés ;
\item $R_3N$ (ou $R-N(-R')-R''$) si les trois hydrogènes ont été remplacés.
\end{itemize}

\begin{popfigure}
\begin{center}
\begin{tikzpicture}[scale=0.95, every node/.style={font=\small}]
% Ammoniac pyramidal
\begin{scope}
  \node[draw=popblue, fill=popblue!20, circle, minimum size=6mm, font=\bfseries\color{popblue}] (N) at (0,0.6) {N};
  % H atoms at base of pyramid
  \node[draw=popdark, fill=popdark, circle, minimum size=4mm] (H1) at (-1.0,-0.6) {H};
  \node[draw=popdark, fill=popdark, circle, minimum size=4mm] (H2) at (1.0,-0.6) {H};
  \node[draw=popdark, fill=popdark, circle, minimum size=4mm] (H3) at (0,-1.3) {H};
  \draw[thick, popdark] (N) -- (H1);
  \draw[thick, popdark] (N) -- (H2);
  \draw[thick, popdark] (N) -- (H3);
  \draw[dashed, popmuted] (H1) -- (H2) -- (H3) -- cycle;
  % Lone pair representation
  \draw[popblue, thick, decorate, decoration={snake, amplitude=1pt, segment length=4pt}] (N) -- ++(90:0.7) node[above, popblue, font=\scriptsize] {Doublet non liant};
\end{scope}

% Substitution arrows
\draw[-{Stealth[length=3mm]}, very thick, poporange] (1.8,0.8) -- (3.2,1.6) node[midway, above, sloped, font=\scriptsize] {$-$\ce{H}, $+R$};
\draw[-{Stealth[length=3mm]}, very thick, popgreen]  (2.0,0.0) -- (3.2,0.0) node[midway, above, font=\scriptsize] {$-2$\ce{H}, $+2R$};
\draw[-{Stealth[length=3mm]}, very thick, poppurple] (1.8,-0.8) -- (3.2,-1.6) node[midway, below, sloped, font=\scriptsize] {$-3$\ce{H}, $+3R$};

% Three classes boxes
\node[draw=poporange, fill=poporangeL, rounded corners, minimum width=2.8cm, minimum height=1.0cm, align=center, font=\small\bfseries] at (5.8,1.7) {$R-NH_2$ \\ \textit{Primaire}};
\node[draw=popgreen, fill=popgreenL, rounded corners, minimum width=2.8cm, minimum height=1.0cm, align=center, font=\small\bfseries] at (5.8,0.0) {$R_2NH$ \\ \textit{Secondaire}};
\node[draw=poppurple, fill=poppurpleL, rounded corners, minimum width=2.8cm, minimum height=1.0cm, align=center, font=\small\bfseries] at (5.8,-1.7) {$R_3N$ \\ \textit{Tertiaire}};
\end{tikzpicture}
\end{center}
\legende{De l'ammoniac aux trois classes d'amines : chaque atome d'hydrogène remplacé par un groupe carboné $R$ fait monter la molécule d'un « étage » de substitution. Le doublet non liant (représenté en bleu) est conservé à chaque étape.}
\end{popfigure}

\courssub{Les trois classes d'amines}

\begin{definition}[Classe d'une amine]
La \cle{classe} d'une amine est déterminée par le \textbf{nombre de groupes carbonés directement liés à l'atome d'azote} :
\begin{itemize}
\item \cle{amine primaire} : un seul groupe carboné sur l'azote, formule générale $R-NH_2$ ; le groupe caractéristique $-NH_2$ est le groupe \cle{amino} ;
\item \cle{amine secondaire} : deux groupes carbonés sur l'azote, formule générale $R_2NH$ ou $R-NH-R'$ ; le groupe $-NH-$ est le groupe \cle{imino} ;
\item \cle{amine tertiaire} : trois groupes carbonés sur l'azote, formule générale $R_3N$ ; l'azote ne porte plus aucun hydrogène.
\end{itemize}
\end{definition}

\begin{exemplebox}[Les trois classes avec le groupe méthyle]
Avec le groupe méthyle \ce{-CH3}, on obtient trois amines différentes :
\begin{itemize}
\item \ce{CH3-NH2} (\chemfig{CH_3-NH_2}) : un seul groupe carboné sur l'azote $\Rightarrow$ amine \textbf{primaire} (méthanamine) ;
\item \ce{(CH3)2NH} (\chemfig{CH_3-NH-CH_3}) : deux groupes carbonés $\Rightarrow$ amine \textbf{secondaire} (N-méthylméthanamine) ;
\item \ce{(CH3)3N} (\chemfig{CH_3-N(-[2]CH_3)-CH_3}) : trois groupes carbonés $\Rightarrow$ amine \textbf{tertiaire} (N,N-diméthylméthanamine).
\end{itemize}
Remarque que la classe ne dépend \emph{pas} de la taille des groupes : \ce{(CH3-CH2-CH2)3N} est tertiaire exactement comme \ce{(CH3)3N}, car elle porte trois groupes carbonés sur l'azote.
\end{exemplebox}

\begin{center}
\begin{tabularx}{\linewidth}{|l|X|X|X|}
\hline
\rowcolor{popblueL} \textbf{Classe} & \textbf{Formule générale} & \textbf{Exemple (représentation)} & \textbf{H restants sur N} \\
\hline
Primaire & $R-NH_2$ & \chemfig{CH_3-CH_2-NH_2} (éthylamine) & 2 \\
\hline
Secondaire & $R_2NH$ & \chemfig{CH_3-NH-CH_3} (diméthylamine) & 1 \\
\hline
Tertiaire & $R_3N$ & \chemfig{CH_3-N(-[2]CH_3)-CH_3} (triméthylamine) & 0 \\
\hline
\end{tabularx}
\end{center}

\begin{attention}[Erreur fréquente : confondre la classe de l'amine et celle de l'alcool]
La classe d'une amine se détermine par le \textbf{nombre de groupes carbonés liés à l'azote}, et \emph{non} par la nature primaire, secondaire ou tertiaire du carbone qui porte le groupe amino. Piège classique du baccalauréat : la \emph{tert}-butylamine \ce{(CH3)3C-NH2} est une amine \textbf{primaire}, car l'azote n'est lié qu'à \emph{un seul} groupe carboné, même si ce groupe est un carbone tertiaire ! Pour un alcool, c'est le nombre de groupes carbonés portés par le \emph{carbone fonctionnel} qui compte ; pour une amine, c'est le nombre de groupes carbonés portés par l'\emph{azote}. Deux critères différents : ne les mélange jamais.
\end{attention}

\begin{exoresolu}[Identifier la classe d'une amine]
On donne les quatre amines suivantes : (a) \ce{CH3-CH2-NH2} ; (b) \ce{CH3-NH-CH2-CH3} ; (c) \ce{(CH3)2CH-NH2} ; (d) \ce{CH3-CH2-N(CH3)-CH2-CH3}. Déterminer la classe de chacune en justifiant.
\tcblower
\textbf{Méthode} : on repère l'atome d'azote, puis on compte \emph{uniquement} le nombre de groupes carbonés qui lui sont directement liés.
\begin{itemize}
\item (a) \ce{CH3-CH2-NH2} : l'azote est lié à \textbf{un} groupe carboné (éthyle) et à deux hydrogènes $\Rightarrow$ amine \textbf{primaire}.
\item (b) \ce{CH3-NH-CH2-CH3} : l'azote est lié à \textbf{deux} groupes carbonés (méthyle et éthyle) et à un hydrogène $\Rightarrow$ amine \textbf{secondaire}.
\item (c) \ce{(CH3)2CH-NH2} : attention au piège ! L'azote est lié à \textbf{un seul} groupe carboné (isopropyle), même si le carbone voisin est ramifié $\Rightarrow$ amine \textbf{primaire}.
\item (d) \ce{CH3-CH2-N(CH3)-CH2-CH3} : l'azote est lié à \textbf{trois} groupes carbonés (deux éthyles et un méthyle) et à aucun hydrogène $\Rightarrow$ amine \textbf{tertiaire}.
\end{itemize}
\textbf{Bilan} : le critère est le décompte des liaisons \ce{N-C}, rien d'autre. Les groupes carbonés peuvent être identiques ou différents, cela ne change pas la classe.
\end{exoresolu}

\courssub{Des amines dans la vie quotidienne}

Les amines ne sont pas des curiosités de laboratoire : elles sont partout autour de toi. L'\cle{éthylamine} \ce{CH3-CH2-NH2} intervient dans la synthèse de colorants et de produits pharmaceutiques. La \cle{triméthylamine} \ce{(CH3)3N} est célèbre pour une raison beaucoup moins noble : c'est elle qui se dégage du poisson avarié. La \cle{chloroquine}, antipaludéen historique distribué dans les centres de santé du Cameroun, possède dans sa structure un atome d'azote d'amine tertiaire et un atome d'azote d'amine secondaire. La \cle{caféine} du café de l'Ouest et l'\cle{adrénaline} sécrétée par ton organisme lors d'un effort sont également des molécules azotées de cette grande famille.

\begin{savaistu}[Pourquoi le poisson avarié sent-il mauvais ?]
Les poissons marins accumulent dans leurs tissus un composé azoté, l'oxyde de triméthylamine, qui les protège du sel. Après la mort du poisson, les bactéries transforment ce composé en \cle{triméthylamine} \ce{(CH3)3N}, une amine tertiaire volatile à l'odeur âcre et nauséabonde. C'est elle la responsable de « l'odeur de poisson pas frais » sur les étals du marché Mokolo ! De manière générale, les amines de petite masse molaire sont des composés volatils aux odeurs désagréables : la putrescine et la cadavérine, produites lors de la décomposition des protéines, portent des noms qui parlent d'eux-mêmes. Astuce de grand-mère confirmée par la chimie : arroser le poisson de jus de citron (acide citrique) transforme l'amine en ion ammonium non volatil, et l'odeur disparaît — une première illustration du caractère basique des amines que nous étudierons bientôt.
\end{savaistu}

\begin{aretenir}[L'essentiel de la section]
\begin{itemize}
\item L'ammoniac \ce{NH3} est une molécule \cle{pyramidale} ($\widehat{HNH} \approx \SI{107}{\degree}$) dont l'azote porte un \cle{doublet non liant} : c'est la source de toute la réactivité des amines.
\item Une \cle{amine} dérive de \ce{NH3} par remplacement de 1, 2 ou 3 atomes d'hydrogène par des groupes carbonés.
\item Les trois classes : \textbf{primaire} $R-NH_2$ (1 groupe carboné), \textbf{secondaire} $R_2NH$ (2 groupes), \textbf{tertiaire} $R_3N$ (3 groupes).
\item La classe se détermine en comptant les groupes carbonés \emph{liés à l'azote}, jamais en regardant la nature des carbones.
\end{itemize}
\end{aretenir}

Maintenant que tu sais reconnaître et classer une amine, il est temps d'apprendre à les \emph{nommer} correctement : c'est l'objet de la section suivante, consacrée à la nomenclature.

\coursec{Nomenclature des amines}

Dans la section précédente, nous avons vu comment l'on passe de la molécule d'ammoniac \ce{NH3} aux amines en remplaçant un, deux ou trois atomes d'hydrogène par des groupes alkyles ou aryles, ce qui définit les trois classes d'amines (primaires, secondaires, tertiaires). Très bien : nous savons désormais \emph{reconnaître} une amine et sa classe. Mais en chimie, savoir reconnaître ne suffit pas : il faut aussi savoir \emph{nommer}. Imagine un pharmacien de Douala qui reçoit une ordonnance : si deux molécules différentes portaient le même nom, ou si une même molécule avait deux noms, les conséquences seraient dramatiques. La nomenclature est la langue officielle du chimiste : elle attribue à chaque molécule un nom unique, construit selon des règles précises fixées par l'IUPAC. C'est cette langue que nous allons apprendre ici, pas à pas, avec des exemples de difficulté croissante.

\courssub{Pourquoi une nomenclature spécifique pour les amines ?}

\begin{definition}[Nommer une amine]
Nommer une amine, c'est construire un nom qui renseigne sur :
\begin{itemize}
\item la \cle{chaîne carbonée principale} liée à l'atome d'azote ;
\item la \cle{position} de l'atome d'azote sur cette chaîne (indice de position) ;
\item les \cle{autres groupes alkyles} éventuellement portés par l'azote (préfixe $N$- ou $N,N$-).
\end{itemize}
Le nom d'une amine se termine toujours par le suffixe \cle{« amine »}.
\end{definition}

L'idée directrice est simple : l'atome d'azote joue le rôle de « centre de gravité » de la molécule. On choisit la chaîne carbonée la plus longue parmi celles qui lui sont attachées, on la nomme comme un alcane, on remplace le « e » final par « amine », et l'on signale tout le reste autour de l'azote.

\courssub{Règle 1 : le nom systématique — alcane + suffixe « amine »}

\begin{propriete}[Règle fondamentale de nommage]
Pour nommer une amine :
\begin{enumerate}
\item on identifie la \textbf{plus longue chaîne carbonée} liée à l'atome d'azote : c'est la \cle{chaîne principale} ;
\item on la nomme comme l'alcane correspondant, en remplaçant le « e » final par le suffixe \textbf{« amine »} ;
\item on fait précéder le suffixe de l'\textbf{indice de position} de l'atome d'azote sur la chaîne, en numérotant la chaîne de façon à donner à l'azote \textbf{l'indice le plus petit possible} ;
\item on cite devant le nom, par ordre alphabétique, les éventuels substituants portés par la chaîne, avec leurs indices de position.
\end{enumerate}
\end{propriete}

\begin{exemplebox}[Nommer \ce{CH3-CH(NH2)-CH3}]
La molécule \chemfig{CH_3-CH(-[2]NH_2)-CH_3} possède une chaîne de trois atomes de carbone : l'alcane correspondant est le \textbf{propane}. L'atome d'azote est porté par le carbone central. En numérotant la chaîne pour donner à l'azote l'indice le plus petit, on obtient l'indice 2 (et non 1 !). Le nom est donc \textbf{propan-2-amine}. C'est une amine primaire (l'azote n'est lié qu'à un seul carbone).
\end{exemplebox}

\begin{attention}[Piège classique : la position de l'azote]
Beaucoup d'élèves nomment \ce{CH3-CH(NH2)-CH3} « propan-1-amine » par réflexe, car ils lisent la formule de gauche à droite. C'est \textbf{faux} : l'azote est sur le carbone n° 2. Rappelle-toi la règle d'or : on numérote la chaîne dans le sens qui donne à l'azote \textbf{l'indice le plus petit}. Ainsi :
\begin{itemize}
\item \ce{CH3-CH2-CH2-NH2} : propan-1-amine (azote en bout de chaîne) ;
\item \ce{CH3-CH(NH2)-CH3} : propan-2-amine (azote au milieu).
\end{itemize}
Ces deux molécules sont des \textbf{isomères de position} : des noms différents pour des propriétés différentes !
\end{attention}

\courssub{Règle 2 : le nom courant en « alkylamine »}

Pour les amines simples, un usage très répandu (et toléré par l'IUPAC) consiste à nommer le ou les groupes alkyles liés à l'azote, suivis du mot \textbf{« amine »} en un seul mot.

\begin{exemplebox}[Quelques noms courants]
\begin{itemize}
\item \ce{CH3-NH2} : \textbf{méthylamine} (nom systématique : méthanamine) ;
\item \ce{CH3-CH2-NH2} : \textbf{éthylamine} (nom systématique : éthanamine) ;
\item \ce{CH3-CH2-CH2-NH2} : \textbf{propylamine} (nom systématique : propan-1-amine).
\end{itemize}
Pour les amines secondaires et tertiaires symétriques, on utilise les préfixes \emph{di-} et \emph{tri-} : \ce{(CH3)2NH} est la \textbf{diméthylamine}, \ce{(C2H5)3N} est la \textbf{triéthylamine}.
\end{exemplebox}

\begin{attention}[Les deux systèmes cohabitent]
Au Baccalauréat TI, les deux nomenclatures sont acceptées, mais le nom \textbf{systématique} (éthanamine, propan-2-amine...) est celui qui montre ta maîtrise des règles. Le nom courant (éthylamine) est pratique pour les molécules simples, mais il devient ambigu dès que la molécule se complique : c'est là qu'intervient la règle du préfixe $N$-.
\end{attention}

\courssub{Cas des amines secondaires et tertiaires : le préfixe $N$-}

Voici le point délicat — et le plus riche en erreurs au Baccalauréat. Quand l'azote porte \textbf{plusieurs groupes alkyles différents}, comment les distinguer dans le nom ?

\begin{propriete}[Règle du préfixe $N$-]
Pour une amine secondaire ou tertiaire dont les groupes liés à l'azote sont différents :
\begin{enumerate}
\item on choisit comme \textbf{chaîne principale} la \textbf{plus longue chaîne carbonée} liée à l'azote ; elle donne le nom de base (alcane + « amine » + indice de position) ;
\item les \textbf{autres groupes alkyles}, directement liés à l'azote, sont cités en préfixes, précédés de la lettre \textbf{$N$-} en italique (ou $N,N$- s'il y en a deux) pour indiquer qu'ils sont portés par l'\textbf{azote} et non par la chaîne carbonée ;
\item les préfixes sont cités par \textbf{ordre alphabétique}.
\end{enumerate}
\end{propriete}

\begin{exemplebox}[Application de la règle du préfixe $N$-]
\begin{itemize}
\item \ce{CH3-NH-CH2-CH3} : la plus longue chaîne liée à N a 2 carbones (éthane) ; le groupe méthyle est porté par l'azote. Nom : \textbf{$N$-méthyléthanamine}.
\item \ce{CH3-N(-CH3)-CH3} : trois groupes méthyles identiques ; la chaîne principale a 1 carbone (méthane), les deux autres méthyles sont sur l'azote. Nom : \textbf{$N,N$-diméthylméthanamine} (ou triméthylamine).
\item \ce{CH3-CH2-N(-CH3)-CH2-CH3} : chaîne principale à 2 carbones (éthane), un méthyle et un éthyle sur l'azote. Nom : \textbf{$N$-éthyl-$N$-méthyléthanamine} (ordre alphabétique : éthyl avant méthyl).
\end{itemize}
\end{exemplebox}

\begin{attention}[Erreur fréquente n° 1 : oublier le préfixe $N$-]
Écrire « méthyléthanamine » au lieu de « \textbf{$N$-méthyléthanamine} » est une faute grave : sans le $N$-, le nom signifierait que le groupe méthyle est greffé sur la \textbf{chaîne carbonée}, ce qui correspond à une molécule complètement différente (\ce{CH3-CH(NH2)-CH3}, la propan-2-amine !). Le préfixe $N$- est le panneau indicateur qui dit : « ce groupe est accroché à l'azote, pas au carbone ». De même, « éthylméthylamine » est un nom courant toléré mais imprécis : préfère toujours $N$-méthyléthanamine.
\end{attention}

\begin{popfigure}
\begin{center}
\chemfig{CH_3-N(-[2]CH_3)(-[6]CH_3)}
\hspace{2cm}
\chemfig{CH_3-CH_2-N(-[2]CH_3)(-[6]CH_2-CH_3)}
\end{center}
\vspace{0,3cm}
\begin{center}
\textbf{À gauche :} $N,N$-diméthylméthanamine — chaîne principale : méthane (1 C), deux méthyles sur N.\par
\textbf{À droite :} $N$-éthyl-$N$-méthyléthanamine — chaîne principale : éthane (2 C), un méthyle et un éthyle sur N.
\end{center}
\legende{Deux amines tertiaires : la chaîne principale est la plus longue chaîne carbonée liée à l'azote ; les autres groupes sont signalés par le préfixe $N$-.}
\end{popfigure}

\courssub{Nomenclature des amines aromatiques}

Le programme mentionne explicitement les groupes \emph{aryles}. Les amines aromatiques (arylamines) dérivent de l'aniline (nom retenu IUPAC pour \ce{C6H5NH2}, nom systématique : benzenamine).

\begin{propriete}[Règles pour les amines aromatiques]
\begin{enumerate}
\item La molécule de référence est l'\textbf{aniline} (\ce{C6H5NH2}).
\item Les substituants sur le \textbf{cycle} sont indiqués par leur position (ortho, méta, para ou indices 2, 3, 4) : \textbf{2-méthylaniline} (ortho-toluidine), \textbf{4-méthylaniline} (para-toluidine).
\item Les substituants sur l'\textbf{azote} portent le préfixe \textbf{$N$-} : \ce{C6H5-NH-CH3} est la \textbf{$N$-méthylaniline} (ou $N$-méthylbenzenamine).
\item Pour les amines di/triaryles : \ce{(C6H5)2NH} est la \textbf{diphénylamine}, \ce{(C6H5)3N} est la \textbf{triphénylamine}.
\end{enumerate}
\end{propriete}

\begin{exemplebox}[Exemples d'amines aromatiques]
\begin{itemize}
\item \ce{C6H5-NH2} : \textbf{aniline} (benzenamine) — amine primaire aromatique.
\item \ce{CH3-C6H4-NH2} (para) : \textbf{4-méthylaniline} (para-toluidine).
\item \ce{C6H5-NH-CH2-CH3} : \textbf{$N$-éthylaniline} (ou $N$-éthylbenzenamine) — amine secondaire mixte.
\item \ce{C6H5-N(-CH3)-CH2-CH3} : \textbf{$N$-éthyl-$N$-méthylaniline} — amine tertiaire mixte.
\end{itemize}
\end{exemplebox}

\courssub{Choix de la chaîne principale et numérotation}

Résumons les deux décisions clés que tout élève doit savoir prendre face à une formule d'amine.

\textbf{Choix de la chaîne principale.} Quand l'azote porte plusieurs chaînes carbonées, on retient \textbf{la plus longue}. Exemple : dans \ce{CH3-CH2-CH2-NH-CH3}, l'azote porte un groupe propyle (3 C) et un groupe méthyle (1 C). La chaîne principale est le propane ; le nom est \textbf{$N$-méthylpropan-1-amine} (et non $N$-propylméthanamine).

\textbf{Numérotation.} On numérote la chaîne principale de l'extrémité la plus proche de l'atome d'azote, de sorte que l'azote reçoive \textbf{l'indice le plus petit}. Exemple : \ce{CH3-CH2-CH(NH2)-CH2-CH3} : l'azote est sur le carbone 3 quel que soit le sens de numérotation : c'est la \textbf{pentan-3-amine}. Mais dans \ce{CH3-CH(NH2)-CH2-CH3}, on numérote par la gauche : \textbf{butan-2-amine} (et non butan-3-amine).

\begin{methode}[Pour nommer une amine sans se tromper]
\begin{enumerate}
\item \textbf{Repérer l'atome d'azote} et compter le nombre de groupes carbonés qu'il porte (classe de l'amine : primaire, secondaire ou tertiaire).
\item \textbf{Identifier la plus longue chaîne carbonée} liée à l'azote : c'est la chaîne principale.
\item \textbf{Nommer cette chaîne} comme l'alcane correspondant, remplacer le « e » final par « amine » et donner à l'azote \textbf{l'indice de position le plus petit} (ex. : butan-2-amine).
\item \textbf{Citer les autres groupes} portés par l'azote avec le préfixe $N$- (ou $N,N$-), par ordre alphabétique, devant le nom de base.
\item \textbf{Vérifier} : le nom obtenu permet-il de redessiner la molécule sans ambiguïté ? Si oui, c'est gagné.
\end{enumerate}
\end{methode}

\begin{popfigure}
\begin{center}
\begin{tikzpicture}[
node distance=0.55cm,
etape/.style={rectangle, rounded corners=4pt, draw=popblue, fill=popblueL, text width=6.6cm, align=center, font=\small, inner sep=6pt},
decision/.style={rectangle, rounded corners=4pt, draw=poporange, fill=poporangeL, text width=6.6cm, align=center, font=\small, inner sep=6pt},
fleche/.style={-{Stealth[length=3mm]}, thick, popdark}]
\node[etape] (e1) {\textbf{Étape 1} — Repérer l'atome d'azote et compter les groupes carbonés qu'il porte (classe de l'amine).};
\node[etape, below=of e1] (e2) {\textbf{Étape 2} — Choisir la \textbf{plus longue chaîne carbonée} liée à N : chaîne principale.};
\node[etape, below=of e2] (e3) {\textbf{Étape 3} — Numéroter la chaîne pour donner à N \textbf{l'indice le plus petit}.};
\node[decision, below=of e3] (e4) {\textbf{Étape 4} — Nommer : alcane + « amine » + indice (ex. : propan-2-amine).};
\node[etape, below=of e4] (e5) {\textbf{Étape 5} — Ajouter devant, par ordre alphabétique, les groupes portés par N avec le préfixe \textbf{$N$-} ou \textbf{$N,N$-}.};
\node[decision, below=of e5] (e6) {\textbf{Vérification} — Le nom permet-il de redessiner la molécule sans ambiguïté ?};
\draw[fleche] (e1) -- (e2);
\draw[fleche] (e2) -- (e3);
\draw[fleche] (e3) -- (e4);
\draw[fleche] (e4) -- (e5);
\draw[fleche] (e5) -- (e6);
\end{tikzpicture}
\end{center}
\legende{Organigramme récapitulatif : les cinq étapes du nommage d'une amine, suivies de la vérification finale.}
\end{popfigure}

\courssub{Exercices d'application : de la plus simple à la plus complexe}

\begin{exoresolu}[Nommer des amines de complexité croissante]
Nomme les cinq amines suivantes et précise leur classe :
\begin{enumerate}
\item \ce{CH3-CH2-CH2-NH2}
\item \ce{CH3-CH(NH2)-CH2-CH3}
\item \ce{CH3-NH-CH2-CH2-CH3}
\item \ce{CH3-CH2-N(-CH3)-CH2-CH2-CH3}
\item \ce{(CH3)2CH-N(-CH3)-CH2-CH3}
\end{enumerate}
\tcblower
\textbf{Solution détaillée.}
\begin{enumerate}
\item \ce{CH3-CH2-CH2-NH2} : une seule chaîne de 3 carbones, azote en bout de chaîne. Nom : \textbf{propan-1-amine} (propylamine). Amine \textbf{primaire}.
\item \ce{CH3-CH(NH2)-CH2-CH3} : chaîne de 4 carbones ; on numérote par la gauche pour donner à N l'indice 2. Nom : \textbf{butan-2-amine}. Amine \textbf{primaire}. (Piège évité : pas butan-3-amine !)
\item \ce{CH3-NH-CH2-CH2-CH3} : amine secondaire ; la plus longue chaîne liée à N a 3 carbones (propane), le méthyle est sur l'azote. Nom : \textbf{$N$-méthylpropan-1-amine}. Amine \textbf{secondaire}.
\item \ce{CH3-CH2-N(-CH3)-CH2-CH2-CH3} : amine tertiaire ; la plus longue chaîne a 3 carbones (propane) ; sur l'azote : un éthyle et un méthyle, cités par ordre alphabétique. Nom : \textbf{$N$-éthyl-$N$-méthylpropan-1-amine}. Amine \textbf{tertiaire}.
\item \ce{(CH3)2CH-N(-CH3)-CH2-CH3} : l'azote porte un groupe éthyle (2 C), un méthyle (1 C) et un groupe isopropyle \ce{-CH(CH3)2} (3 C). La plus longue chaîne liée à l'azote est le groupe isopropyle (3 atomes de carbone) : la chaîne principale est donc le \textbf{propane}. L'azote est porté par le carbone 2 de cette chaîne (carbone secondaire) : le nom de base est \textbf{propan-2-amine}. Sur l'azote, on a un méthyle et un éthyle (ordre alphabétique : éthyl, méthyl) $\rightarrow$ \textbf{$N$-éthyl-$N$-méthyl}. Sur la chaîne, le carbone 2 porte un méthyle $\rightarrow$ \textbf{2-méthyl}. Nom complet : \textbf{$N$-éthyl-$N$,2-diméthylpropan-2-amine}. Amine \textbf{tertiaire}. (Le méthyle sur le carbone 2 porte un simple indice, celui sur l'azote porte le préfixe $N$-.)
\end{enumerate}
\end{exoresolu}

\begin{exercice}[À toi de jouer]
\begin{enumerate}
\item Nomme les amines suivantes et donne leur classe : \textbf{a)} \ce{CH3-CH2-NH-CH3} ; \textbf{b)} \ce{CH3-CH2-CH(NH2)-CH3} ; \textbf{c)} \ce{(CH3-CH2)3N} ; \textbf{d)} \ce{CH3-N(-CH3)-CH2-CH2-CH2-CH3}.
\item Écris la formule semi-développée de : \textbf{a)} la pentan-2-amine ; \textbf{b)} la $N,N$-diéthyléthanamine ; \textbf{c)} la $N$-méthylbutan-2-amine.
\item Un élève nomme \ce{CH3-NH-CH2-CH3} « éthylméthanamine ». Explique pourquoi ce nom est incorrect et propose le nom exact.
\end{enumerate}
\end{exercice}

\begin{corrige}[Exercice 1]
\begin{enumerate}
\item \textbf{a)} $N$-méthyléthanamine, secondaire ; \textbf{b)} butan-2-amine, primaire ; \textbf{c)} $N,N$-diéthyléthanamine (triéthylamine), tertiaire ; \textbf{d)} $N,N$-diméthylbutan-1-amine, tertiaire.
\item \textbf{a)} \ce{CH3-CH(NH2)-CH2-CH2-CH3} ; \textbf{b)} \ce{CH3-CH2-N(-CH2-CH3)-CH2-CH3} ; \textbf{c)} \ce{CH3-NH-CH(CH3)-CH2-CH3}.
\item « éthylméthanamine » placerait le groupe éthyle sur la chaîne méthane, ce qui est impossible (le méthane n'a qu'un carbone). Le nom correct est \textbf{$N$-méthyléthanamine} : la chaîne principale est l'éthane (2 C) et le méthyle est porté par l'azote.
\end{enumerate}
\end{corrige}

\begin{aretenir}[L'essentiel de la nomenclature des amines]
\begin{itemize}
\item Chaîne principale = \textbf{plus longue chaîne carbonée liée à l'azote}.
\item Nom de base = alcane correspondant dont le « e » final est remplacé par \textbf{« amine »}, précédé de l'\textbf{indice de position de N} (le plus petit possible) : propan-2-amine.
\item Groupes supplémentaires portés par l'azote : préfixe \textbf{$N$-} ou \textbf{$N,N$-}, cités par ordre alphabétique : $N$-méthyléthanamine, $N,N$-diméthylméthanamine.
\item Amines aromatiques : base \textbf{aniline} (benzenamine) ; substituants sur cycle (indices 2,3,4), sur azote ($N$-).
\item Noms courants tolérés pour les cas simples : éthylamine, diméthylamine, triéthylamine.
\item Vérification finale : le nom doit permettre de redessiner la molécule \textbf{sans aucune ambiguïté}.
\end{itemize}
\end{aretenir}

Tu sais désormais donner un nom rigoureux à n'importe quelle amine, qu'elle soit primaire, secondaire ou tertiaire, aliphatique ou aromatique. Cette compétence te servira immédiatement dans la suite de la leçon : quand nous étudierons les propriétés physiques des amines (état, solubilité, liaisons hydrogène), nous comparerons des molécules précises — et il faudra pouvoir les identifier sans hésiter. En avant !

\coursec{Propriétés physiques des amines}

Dans la section précédente, tu as appris à nommer les amines et à reconnaître leur classe (primaire, secondaire ou tertiaire). Cette classification n'est pas qu'un exercice de vocabulaire : elle a des conséquences \emph{directes et mesurables} sur le comportement physique des amines. Pourquoi la méthylamine est-elle un gaz alors que l'aniline est un liquide ? Pourquoi l'éthylamine bout-elle à \SI{17}{\celsius} alors que son isomère tertiaire, la triméthylamine, bout à \SI{3}{\celsius} seulement ? Pourquoi la méthylamine se mélange-t-elle à l'eau en toutes proportions alors que la phénylamine n'y est que très peu soluble ? Toutes ces questions trouvent leur réponse dans une notion capitale que tu as déjà rencontrée avec les alcools : la \cle{liaison hydrogène}. Cette section va te montrer comment la structure moléculaire d'une amine commande son état physique, sa température d'ébullition et sa solubilité.

\courssub{État physique des amines à température ambiante}

\textbf{Intuition.} Une molécule est d'autant plus difficile à « arracher » à ses voisines (pour passer à l'état gazeux) qu'elle est lourde et qu'elle est fortement attirée par ses voisines. On s'attend donc à ce que l'état physique d'une amine dépende de sa \cle{masse molaire} et de la force des interactions entre ses molécules.

\begin{propriete}[État physique des amines]
À température ambiante (environ \SI{25}{\celsius}) :
\begin{itemize}
\item les amines de très faible masse molaire, comme la \cle{méthylamine} \ce{CH3NH2} (T$_{\text{éb}} = \SI{-6}{\celsius}$), la diméthylamine (T$_{\text{éb}} = \SI{7}{\celsius}$) ou la triméthylamine (T$_{\text{éb}} = \SI{3}{\celsius}$), sont \textbf{gazeuses} ;
\item les amines de masse molaire intermédiaire, comme l'\cle{éthylamine} \ce{C2H5NH2} (T$_{\text{éb}} = \SI{17}{\celsius}$), la propan-1-amine (T$_{\text{éb}} = \SI{49}{\celsius}$) ou l'aniline \ce{C6H5NH2} (T$_{\text{éb}} = \SI{184}{\celsius}$), sont \textbf{liquides} ;
\item les amines de grande masse molaire (longues chaînes carbonées, généralement plus de 12 atomes de carbone) sont \textbf{solides}.
\end{itemize}
\end{propriete}

On retrouve la tendance générale observée pour toutes les familles organiques : \textbf{quand la masse molaire augmente, la température d'ébullition augmente}, et l'on passe progressivement de l'état gazeux à l'état liquide puis solide. Mais cette tendance seule n'explique pas tout : il faut maintenant comparer des amines \emph{de masse voisine} pour découvrir le rôle décisif des liaisons hydrogène.

\courssub{Les liaisons hydrogène intermoléculaires : la clé de tout}

\textbf{Intuition.} L'atome d'azote est plus électronégatif que l'hydrogène. Dans une liaison \ce{N-H}, les électrons de la liaison sont attirés vers l'azote : l'hydrogène porte alors une charge partielle positive $\delta^+$ et l'azote une charge partielle négative $\delta^-$, sans oublier le \cle{doublet non liant} de l'azote. L'hydrogène $\delta^+$ d'une molécule peut donc être attiré par le doublet de l'azote $\delta^-$ d'une molécule voisine : c'est la liaison hydrogène, notée \ce{N-H\bond{...}N}.

\begin{definition}[Liaison hydrogène intermoléculaire dans les amines]
Une \cle{liaison hydrogène} intermoléculaire est une interaction attractive qui s'établit entre un atome d'hydrogène porté par un atome très électronégatif (donneur, ici \ce{N-H}) et le doublet non liant d'un atome très électronégatif d'une molécule voisine (accepteur, ici \ce{N}).
\begin{itemize}
\item Les \textbf{amines primaires} ($R\text{-}NH_2$) possèdent \emph{deux} hydrogènes sur l'azote : elles forment des liaisons hydrogène ($N\text{-}H\cdots N$) nombreuses entre leurs molécules.
\item Les \textbf{amines secondaires} ($R\text{-}NH\text{-}R'$) possèdent \emph{un} hydrogène sur l'azote : elles forment aussi des liaisons hydrogène intermoléculaires, mais en nombre moindre.
\item Les \textbf{amines tertiaires} ($R\text{-}N(R')\text{-}R''$) ne possèdent \emph{aucun} hydrogène sur l'azote : elles \textbf{ne peuvent pas} former de liaisons hydrogène entre leurs propres molécules.
\end{itemize}
\end{definition}

\begin{popfigure}
\begin{center}
\begin{tikzpicture}[scale=1.0, every node/.style={font=\small}]
% --- Ethylamine - ethylamine ---
\node[font=\bfseries\small, popblue] at (2.2,3.4) {Entre deux molécules d'éthylamine};
% molecule 1
\node[popdark] (N1) at (0.6,2.2) {N};
\node[popdark, left=0.28cm of N1] (C1) {\ce{CH3-CH2}};
\node[popdark, above right=0.16cm and 0.18cm of N1] (H1a) {H};
\node[popdark, below right=0.16cm and 0.18cm of N1] (H1b) {H};
\draw[popline, thick] (C1) -- (N1);
\draw[popline, thick] (N1) -- (H1a);
\draw[popline, thick] (N1) -- (H1b);
% molecule 2
\node[popdark] (N2) at (4.0,2.2) {N};
\node[popdark, right=0.28cm of N2] (C2) {\ce{CH2-CH3}};
\node[popdark, above left=0.16cm and 0.18cm of N2] (H2a) {H};
\node[popdark, below left=0.16cm and 0.18cm of N2] (H2b) {H};
\draw[popline, thick] (N2) -- (C2);
\draw[popline, thick] (N2) -- (H2a);
\draw[popline, thick] (N2) -- (H2b);
% H bond
\draw[poporange, very thick, dashed] (H1a.east) -- (N2.west)
  node[midway, above, poporange, font=\footnotesize] {liaison H};
% lone pair N2
\fill[popgreen] ($(N2.west)+(-0.12,0.10)$) circle (1.2pt);
\fill[popgreen] ($(N2.west)+(-0.12,-0.10)$) circle (1.2pt);
\node[popmuted, font=\footnotesize] at (2.3,1.35) {\ce{N-H\bond{...}N} : doublet de l'azote accepteur};

% --- Ethylamine - water ---
\node[font=\bfseries\small, poppurple] at (8.6,3.4) {Entre éthylamine et eau};
\node[popdark] (N3) at (7.0,2.2) {N};
\node[popdark, left=0.28cm of N3] (C3) {\ce{CH3-CH2}};
\node[popdark, above right=0.16cm and 0.18cm of N3] (H3a) {H};
\node[popdark, below right=0.16cm and 0.18cm of N3] (H3b) {H};
\draw[popline, thick] (C3) -- (N3);
\draw[popline, thick] (N3) -- (H3a);
\draw[popline, thick] (N3) -- (H3b);
% water
\node[popdark] (O4) at (10.4,2.2) {O};
\node[popdark, above right=0.16cm and 0.18cm of O4] (H4a) {H};
\node[popdark, below right=0.16cm and 0.18cm of O4] (H4b) {H};
\draw[popline, thick] (O4) -- (H4a);
\draw[popline, thick] (O4) -- (H4b);
\draw[poporange, very thick, dashed] (H3a.east) -- (O4.west)
  node[midway, above, poporange, font=\footnotesize] {liaison H};
\fill[popgreen] ($(O4.west)+(-0.12,0.10)$) circle (1.2pt);
\fill[popgreen] ($(O4.west)+(-0.12,-0.10)$) circle (1.2pt);
\node[popmuted, font=\footnotesize] at (8.7,1.35) {\ce{N-H\bond{...}O} : responsable de la solubilité};
\end{tikzpicture}
\end{center}
\legende{Liaisons hydrogène (pointillés orange) : à gauche, entre deux molécules d'éthylamine \ce{N-H\bond{...}N} ; à droite, entre l'éthylamine et une molécule d'eau \ce{N-H\bond{...}O}. Les points verts représentent le doublet non liant accepteur.}
\end{popfigure}

\begin{attention}[Piège classique sur les amines tertiaires]
\textbf{Erreur fréquente au Baccalauréat :} affirmer que les amines tertiaires forment des liaisons hydrogène entre elles. C'est \textbf{faux} : une amine tertiaire comme la triméthylamine \ce{N(CH3)3} ne porte \emph{aucun} atome d'hydrogène sur son azote. Or, pour \emph{donner} une liaison hydrogène, il faut un hydrogène lié à un atome très électronégatif. Les molécules d'amine tertiaire ne sont donc liées entre elles que par des interactions de Van der Waals, beaucoup plus faibles. Retiens : \textbf{pas de \ce{N-H}, pas de liaison hydrogène intermoléculaire}.
\end{attention}

\courssub{Conséquence sur les températures d'ébullition}

\textbf{Intuition.} Faire bouillir un liquide, c'est fournir assez d'énergie pour séparer les molécules les unes des autres. Si les molécules sont « scotchées » par des liaisons hydrogène, il faut chauffer davantage : la température d'ébullition est plus élevée.

\begin{propriete}[Comparaison des températures d'ébullition]
À nombre d'atomes de carbone comparable :
\begin{enumerate}
\item une \textbf{amine primaire} bout plus haut qu'une \textbf{amine secondaire}, elle-même plus haute qu'une \textbf{amine tertiaire} isomère : T$_{\text{éb}}$(primaire) $>$ T$_{\text{éb}}$(secondaire) $>$ T$_{\text{éb}}$(tertiaire), car le nombre de liaisons hydrogène intermoléculaires diminue ;
\item une \textbf{amine} bout à une température \textbf{inférieure} à celle de l'\textbf{alcool} de masse molaire voisine, car l'azote est moins électronégatif que l'oxygène : la liaison \ce{N-H} est moins polarisée que la liaison \ce{O-H}, donc les liaisons hydrogène \ce{N-H\bond{...}N} sont plus faibles que les liaisons \ce{O-H\bond{...}O} ;
\item une amine bout nettement plus haut que l'\textbf{alcane} de masse voisine, qui ne forme aucune liaison hydrogène.
\end{enumerate}
\end{propriete}

\begin{exemplebox}[L'effet spectaculaire des liaisons hydrogène]
Comparons deux amines \emph{isomères} (même formule brute \ce{C3H9N}, même masse molaire \SI{59}{g.mol^{-1}}) :
\begin{itemize}
\item la \textbf{propan-1-amine} \ce{CH3CH2CH2NH2} (amine primaire) : T$_{\text{éb}} = \SI{49}{\celsius}$ ;
\item la \textbf{triméthylamine} \ce{N(CH3)3} (amine tertiaire) : T$_{\text{éb}} = \SI{3}{\celsius}$.
\end{itemize}
Bien que ces deux amines aient \emph{exactement la même masse molaire}, la propan-1-amine bout \SI{46}{\celsius} \emph{au-dessus} de la triméthylamine. La seule explication : la propan-1-amine, primaire, établit des liaisons hydrogène \ce{N-H\bond{...}N} entre ses molécules, alors que la triméthylamine, tertiaire, n'en établit aucune.
Autre comparaison éloquente : l'éthanol \ce{C2H5OH} ($M = \SI{46}{g.mol^{-1}}$, masse quasi identique à celle de l'éthylamine \SI{45}{g.mol^{-1}}) bout à \SI{78}{\celsius}, soit \SI{61}{\celsius} de plus que l'éthylamine (\SI{17}{\celsius}), car les liaisons hydrogène \ce{O-H\bond{...}O} sont plus fortes que les liaisons \ce{N-H\bond{...}N}.
\end{exemplebox}

\begin{popfigure}
\begin{center}
\begin{tikzpicture}
\begin{axis}[
  width=0.95\linewidth, height=7cm,
  xlabel={Nombre d'atomes de carbone},
  ylabel={T$_{\text{éb}}$ (\si{\celsius})},
  xmin=0.5, xmax=6.5, ymin=-20, ymax=120,
  xtick={1,2,3,4,5,6},
  grid=major, grid style={popline!40},
  legend style={at={(0.03,0.97)}, anchor=north west, font=\footnotesize, draw=popline},
]
% Amines primaires : methylamine -6.5, ethylamine 17, propylamine 49, butylamine 78, pentylamine 104
\addplot[mark=*, mark size=1.8pt, popblue, very thick] coordinates {(1,-6.5) (2,16.6) (3,48.7) (4,77.8) (5,104)};
\addlegendentry{Amines primaires \ce{R-NH2}}
% Amines tertiaires : trimethylamine 3, triethylamine 89 (n=6)
\addplot[mark=square*, mark size=1.8pt, poporange, very thick] coordinates {(3,2.9) (6,89)};
\addlegendentry{Amines tertiaires \ce{R3N}}
% Alcanes : methane -162 (hors cadre) -> ethane -89, propane -42, butane -0.5, pentane 36, hexane 69
\addplot[mark=triangle*, mark size=2.2pt, popgreen, very thick] coordinates {(2,-88.6) (3,-42.1) (4,-0.5) (5,36.1) (6,68.7)};
\addlegendentry{Alcanes \ce{C_{n}H_{2n+2}}}
\end{axis}
\end{tikzpicture}
\end{center}
\legende{Températures d'ébullition en fonction du nombre d'atomes de carbone. Les amines primaires (bleu) bouillent nettement plus haut que les alcanes (vert) grâce aux liaisons hydrogène ; la triméthylamine (orange, $n=3$) bout presque aussi bas que le propane, faute de liaisons hydrogène.}
\end{popfigure}

La courbe résume tout : à $n = 3$, l'amine primaire (propan-1-amine, \SI{49}{\celsius}) bout \SI{46}{\celsius} au-dessus de l'amine tertiaire (triméthylamine, \SI{3}{\celsius}) et \SI{91}{\celsius} au-dessus du propane (\SI{-42}{\celsius}). Les liaisons hydrogène sont bien la grandeur qui fait la différence.

\courssub{Solubilité des amines dans l'eau}

\textbf{Intuition.} L'eau est un solvant qui « aime » les liaisons hydrogène : elle en forme elle-même en permanence. Une molécule organique sera soluble dans l'eau si elle peut « s'insérer » dans ce réseau de liaisons hydrogène, c'est-à-dire si elle peut en donner ou en accepter avec l'eau.

\begin{propriete}[Solubilité des amines dans l'eau]
\begin{itemize}
\item Les amines à \textbf{courte chaîne carbonée} (moins de 6 atomes de carbone environ) sont \textbf{solubles} dans l'eau, certaines même en toutes proportions (méthylamine, éthylamine). Cette solubilité s'explique par la formation de liaisons hydrogène avec les molécules d'eau : \ce{N-H\bond{...}O} (l'amine donne) et \ce{O-H\bond{...}N} (l'eau donne, l'azote accepte grâce à son doublet).
\item La solubilité \textbf{décroît rapidement} lorsque la chaîne carbonée s'allonge : la partie hydrocarbonée \ce{R} est \cle{hydrophobe} (elle ne forme pas de liaisons hydrogène et « repousse » l'eau). Quand cette partie devient prépondérante, l'amine devient pratiquement insoluble : l'aniline \ce{C6H5NH2} n'est soluble qu'à environ \SI{36}{g.L^{-1}} (à \SI{20}{\celsius}), et les amines à longue chaîne sont insolubles.
\end{itemize}
\end{propriete}

C'est le même compromis que pour les alcools : le groupe \ce{-NH2} est la « tête » \cle{hydrophile} (qui aime l'eau), la chaîne carbonée est la « queue » \cle{hydrophobe} (qui fuit l'eau). Tant que la tête domine, l'amine est soluble ; quand la queue prend le dessus, elle ne l'est plus.

\begin{savaistu}[Les amines tertiaires sont quand même un peu solubles !]
On pourrait croire que les amines tertiaires, incapables de former des liaisons hydrogène entre elles, sont totalement insolubles dans l'eau. Il n'en est rien : la triméthylamine est très soluble dans l'eau ! En effet, si elle ne peut pas \emph{donner} de liaison hydrogène (pas de \ce{N-H}), elle peut en \emph{accepter} grâce au \textbf{doublet non liant de son azote}, qui attire les hydrogènes $\delta^+$ des molécules d'eau : \ce{O-H\bond{...}N(CH3)3}. Cette solubilité reste toutefois plus faible que celle des amines primaires de masse comparable, qui peuvent à la fois donner et accepter des liaisons hydrogène. Ce doublet non liant, tu le retrouveras dans la section suivante : c'est lui qui confère aux amines leur caractère basique et nucléophile.
\end{savaistu}

\begin{methode}[Expliquer une propriété physique d'une amine]
Face à toute question du type « expliquer la température d'ébullition » ou « justifier la solubilité », suis toujours la même démarche en trois étapes :
\begin{enumerate}
\item \textbf{Identifier la classe} de l'amine (primaire, secondaire ou tertiaire) et compter les atomes d'hydrogène portés par l'azote.
\item \textbf{Conclure sur les liaisons hydrogène} : deux \ce{N-H} (primaire) $\Rightarrow$ liaisons hydrogène nombreuses ; un seul \ce{N-H} (secondaire) $\Rightarrow$ liaisons hydrogène moins nombreuses ; aucun \ce{N-H} (tertiaire) $\Rightarrow$ aucune liaison hydrogène intermoléculaire (mais l'azote peut en accepter avec l'eau).
\item \textbf{Examiner la longueur de la chaîne carbonée} : plus elle est longue, plus la masse molaire augmente (T$_{\text{éb}}$ plus élevée) et plus la partie hydrophobe domine (solubilité dans l'eau plus faible).
\end{enumerate}
Rédige ensuite ta réponse en reliant explicitement la structure à la propriété : « l'amine X est primaire, donc ses molécules sont associées par des liaisons hydrogène \ce{N-H\bond{...}N}, donc sa température d'ébullition est élevée ».
\end{methode}

\begin{exoresolu}[Comparaison de trois composés de masses voisines]
\textbf{Énoncé.} On considère trois composés de masses molaires voisines : le propane \ce{C3H8} ($M = \SI{44}{g.mol^{-1}}$, T$_{\text{éb}} = \SI{-42}{\celsius}$), l'éthylamine \ce{C2H5NH2} ($M = \SI{45}{g.mol^{-1}}$, T$_{\text{éb}} = \SI{17}{\celsius}$) et l'éthanol \ce{C2H5OH} ($M = \SI{46}{g.mol^{-1}}$, T$_{\text{éb}} = \SI{78}{\celsius}$).
\begin{enumerate}
\item Classer ces composés par température d'ébullition croissante.
\item Justifier ce classement en invoquant les interactions intermoléculaires.
\item Prévoir, en le justifiant, lequel de ces trois composés est le plus soluble dans l'eau.
\end{enumerate}
\tcblower
\textbf{Solution détaillée.}
\begin{enumerate}
\item T$_{\text{éb}}$(propane) $<$ T$_{\text{éb}}$(éthylamine) $<$ T$_{\text{éb}}$(éthanol), soit $\SI{-42}{\celsius} < \SI{17}{\celsius} < \SI{78}{\celsius}$.
\item Le propane est un alcane : molécule apolaire, il n'établit \textbf{aucune liaison hydrogène} ; seules de faibles interactions de Van der Waals relient ses molécules, d'où sa température d'ébullition très basse. L'éthylamine est une \textbf{amine primaire} \ce{C2H5NH2} : ses deux hydrogènes portés par l'azote permettent des liaisons hydrogène \ce{N-H\bond{...}N} entre molécules ; il faut donc davantage d'énergie pour les séparer, et la température d'ébullition monte à \SI{17}{\celsius}. L'éthanol forme aussi des liaisons hydrogène \ce{O-H\bond{...}O}, mais l'oxygène est \textbf{plus électronégatif} que l'azote ($\chi_O = 3{,}5 > \chi_N = 3{,}0$) : la liaison \ce{O-H} est plus polarisée que \ce{N-H}, les liaisons hydrogène sont plus fortes, et l'éthanol bout à \SI{78}{\celsius}.
\item L'éthanol est soluble dans l'eau \emph{en toutes proportions}, car il donne et accepte des liaisons hydrogène fortes avec l'eau. L'éthylamine est également très soluble (liaisons hydrogène \ce{N-H\bond{...}O} et \ce{O-H\bond{...}N}), mais un peu moins que l'éthanol. Le propane, apolaire et hydrophobe, est pratiquement insoluble. Le plus soluble est donc \textbf{l'éthanol}.
\end{enumerate}
\end{exoresolu}

\begin{aretenir}[L'essentiel sur les propriétés physiques des amines]
\begin{itemize}
\item L'état physique dépend de la masse molaire : amines légères gazeuses (méthylamine, diméthylamine, triméthylamine), puis liquides (éthylamine, propylamines, aniline), puis solides.
\item Les amines \textbf{primaires et secondaires} forment des liaisons hydrogène intermoléculaires \ce{N-H\bond{...}N} ; les amines \textbf{tertiaires} n'en forment pas entre elles (pas de \ce{N-H}).
\item Conséquence : T$_{\text{éb}}$(primaire) $>$ T$_{\text{éb}}$(secondaire) $>$ T$_{\text{éb}}$(tertiaire) à isomérie près ; et T$_{\text{éb}}$(amine) $<$ T$_{\text{éb}}$(alcool) de masse voisine car N est moins électronégatif que O.
\item Les amines à moins de 6 atomes de carbone sont solubles dans l'eau grâce aux liaisons hydrogène avec l'eau ; la solubilité chute quand la chaîne hydrophobe s'allonge.
\item Une amine tertiaire peut néanmoins \emph{accepter} des liaisons hydrogène de l'eau grâce au doublet de l'azote : elle garde une solubilité partielle.
\end{itemize}
\end{aretenir}

Tu maîtrises désormais le lien entre la structure des amines et leurs propriétés physiques. Mais le doublet non liant de l'azote ne sert pas qu'à accepter des liaisons hydrogène : il peut aussi capter un proton \ce{H+}. C'est précisément ce qui fait des amines des \textbf{bases} — et c'est l'objet de la prochaine section, où tu découvriras pourquoi une solution d'éthylamine dans l'eau rougit la phénolphtaléine.

\coursec{Caractère basique des amines}

Dans la section précédente, nous avons vu que les propriétés physiques des amines (état, solubilité dans l'eau, températures d'ébullition) s'expliquent par la présence du \cle{doublet non liant} porté par l'atome d'azote et par les liaisons hydrogène qu'il permet de former. Ce même doublet non liant est aussi à l'origine des propriétés \emph{chimiques} les plus importantes des amines : leur \cle{caractère basique} et leur \cle{caractère nucléophile}. C'est le caractère basique que nous allons maintenant étudier en détail. Comprends bien cette section : c'est elle qui explique pourquoi tant de médicaments vendus dans les pharmacies de Douala ou de Yaoundé se présentent sous forme de « chlorhydrates » !

\courssub{Origine de la basicité : le doublet non liant de l'azote}

\textbf{L'intuition d'abord.} Rappelle-toi la structure de la molécule d'ammoniac \ce{NH3} : l'atome d'azote y est lié à trois atomes d'hydrogène et porte un \cle{doublet non liant} (deux électrons non engagés dans une liaison). Ce doublet est une région de l'espace riche en électrons, donc chargée négativement. Or un proton \ce{H+} est une entité chargée positivement, dépourvue d'électrons. L'attraction entre les deux est donc naturelle : le doublet de l'azote peut « capturer » un proton en formant avec lui une \cle{liaison covalente de coordination} (liaison dative) : l'azote fournit les deux électrons de la liaison.

Dans une amine, l'atome d'azote conserve ce doublet non liant, qu'il soit lié à un, deux ou trois groupes alkyles. Toute amine possède donc, comme l'ammoniac, la capacité de capter un proton.

\begin{definition}[Une amine est une base de Brønsted]
Selon la théorie de Brønsted, une \cle{base} est une espèce chimique capable de \textbf{capter un proton} \ce{H+}. Une amine est une \cle{base de Brønsted} car le \textbf{doublet non liant} de son atome d'azote peut accepter un proton \ce{H+} pour former une liaison covalente de coordination $\ce{N} \rightarrow \ce{H}$. L'espèce chargée obtenue est un ion \cle{alkylammonium} (ou dialkylammonium, trialkylammonium selon la classe de l'amine).
\end{definition}

\begin{popfigure}
\begin{center}
\begin{tikzpicture}[scale=1.0, every node/.style={font=\small}]
% Amine avec doublet
\node at (0,0) {\chemfig{R-[:-30]N(-[2]H)(-[:-150]H)}};
\node[popblue, font=\large] (doublet) at (0.55,0.75) {\textbf{:}};
\node[popblue, below=0.05cm of doublet, font=\scriptsize] {doublet non liant};
% Flèche courbe du doublet vers H+
\node[popink] (Hplus) at (3.4,1.1) {\ce{H+}};
\draw[-{Stealth[length=3mm]}, popblue, very thick] (doublet.east) .. controls (1.8,1.5) and (2.6,1.5) .. (Hplus.west);
% Flèche de réaction
\draw[-{Stealth[length=3mm]}, popdark, very thick] (4.1,0.3) -- (5.3,0.3);
% Produit : ion ammonium
\node at (7.3,0) {\ce{R-NH3+}};
\node[popink, font=\small] at (7.3,-1.15) {ion alkylammonium};
\end{tikzpicture}
\end{center}
\legende{Capture d'un proton \ce{H+} par le doublet non liant de l'azote d'une amine primaire \ce{R-NH2} : la flèche courbe représente le déplacement du doublet vers le proton pour former la liaison $\ce{N} \rightarrow \ce{H}$.}
\end{popfigure}

\begin{attention}[Ne confonds pas les flèches]
La \textbf{flèche courbe} (utilisée dans les mécanismes) décrit le déplacement d'un \emph{doublet d'électrons} : elle part toujours d'un doublet (ou d'une liaison) et arrive sur un atome. La \textbf{flèche droite} $\ce{->}$ ou la double flèche $\ce{<=>}$ sépare les réactifs des produits dans une équation-bilan. Ne les mélange jamais dans une copie.
\end{attention}

\courssub{Réaction d'une amine avec l'eau : un équilibre acido-basique}

\textbf{L'intuition.} Lorsqu'on dissout une amine dans l'eau, deux espèces capables de capter ou de céder un proton se trouvent en présence : l'amine (qui veut capter \ce{H+}) et l'eau (qui peut céder \ce{H+}, car l'eau est un \cle{ampholyte}). Une partie des molécules d'amine arrache donc un proton aux molécules d'eau : il se forme un ion alkylammonium et un ion hydroxyde \ce{HO-}. La présence des ions \ce{HO-} rend la solution \textbf{basique} : c'est pourquoi une solution aqueuse d'éthylamine a un pH supérieur à 7.

Prenons l'exemple de l'\textbf{éthylamine} \ce{C2H5NH2} :

\[ \ce{C2H5NH2 + H2O <=> C2H5NH3+ + HO-} \]

Commentons cette équation :
\begin{itemize}
\item l'éthylamine capte un proton : c'est la \textbf{base} du couple \ce{C2H5NH3+}/\ce{C2H5NH2} ;
\item l'eau cède un proton : elle joue ici le rôle d'\textbf{acide} du couple \ce{H2O}/\ce{HO-} ;
\item l'ion \ce{C2H5NH3+} formé est l'\cle{ion éthylammonium} ;
\item la double flèche $\ce{<=>}$ indique que la réaction est \textbf{limitée} : c'est un \cle{équilibre chimique}.
\end{itemize}

\begin{methode}[Écrire l'équation de la réaction d'une amine avec l'eau]
Pour écrire correctement l'équation-bilan de la réaction d'une amine avec l'eau, procède en quatre étapes :
\begin{enumerate}
\item \textbf{Identifie} la classe de l'amine et écris sa formule semi-développée (\ce{RNH2}, \ce{R2NH} ou \ce{R3N}).
\item \textbf{Fais capter} un proton \ce{H+} par le doublet de l'azote : l'amine gagne un H et devient un ion ammonium chargé positivement : \ce{RNH3+}, \ce{R2NH2+} ou \ce{R3NH+}.
\item \textbf{Écris} le sous-produit : l'eau, ayant perdu un proton, devient l'ion hydroxyde \ce{HO-}.
\item \textbf{Place la double flèche} $\ce{<=>}$ : la réaction est un équilibre, jamais une réaction totale. Vérifie enfin la conservation des atomes et des charges.
\end{enumerate}
Exemple avec la diéthylamine (amine secondaire) :
\[ \ce{(C2H5)2NH + H2O <=> (C2H5)2NH2+ + HO-} \]
L'ion formé est l'ion \textbf{diéthylammonium}.
\end{methode}

\begin{popfigure}
\begin{center}
\begin{tikzpicture}[scale=1.0, every node/.style={font=\small}]
\node[popblue] (a) at (0,0) {\ce{C2H5NH2}};
\node[popmuted] at (0,-0.75) {base 1};
\node[popink] (b) at (2.6,0) {$+$ \ \ce{H2O}};
\node[popmuted] at (2.6,-0.75) {acide 2};
\draw[{Stealth[length=2.5mm]}-{Stealth[length=2.5mm]}, popdark, very thick] (4.0,0) -- (5.4,0);
\node[popblue] (c) at (7.2,0) {\ce{C2H5NH3+}};
\node[popmuted] at (7.2,-0.75) {acide 1};
\node[popink] (d) at (9.6,0) {$+$ \ \ce{HO-}};
\node[popmuted] at (9.6,-0.75) {base 2};
% Couples
\draw[popgreen, thick, rounded corners] (-1.3,-1.15) rectangle (1.3,0.55);
\draw[popgreen, thick, rounded corners] (5.9,-1.15) rectangle (8.5,0.55);
\draw[popgreen, thick, -{Stealth[length=2mm]}] (0,-1.35) .. controls (2,-2.1) and (5,-2.1) .. (7.2,-1.35);
\node[popgreen, font=\footnotesize] at (3.6,-2.25) {couple 1 : \ce{C2H5NH3+}/\ce{C2H5NH2}};
\draw[poporange, thick, -{Stealth[length=2mm]}] (2.6,-1.35) .. controls (4.5,-2.9) and (8,-2.9) .. (9.6,-1.35);
\node[poporange, font=\footnotesize] at (6.1,-3.05) {couple 2 : \ce{H2O}/\ce{HO-}};
\end{tikzpicture}
\end{center}
\legende{Équilibre entre l'éthylamine et l'eau : les deux couples acide/base mis en jeu. Le proton est transféré de l'acide 2 (\ce{H2O}) vers la base 1 (\ce{C2H5NH2}).}
\end{popfigure}

\textbf{Les amines sont des bases faibles.} La réaction d'une amine avec l'eau est \emph{limitée} : seule une petite fraction des molécules d'amine est transformée en ions alkylammonium. À l'équilibre, la solution contient majoritairement l'amine sous sa forme moléculaire \ce{RNH2}. On dit que les amines sont des \cle{bases faibles}. Par exemple, pour l'éthylamine, la constante d'acidité du couple \ce{C2H5NH3+}/\ce{C2H5NH2} vaut $K_a \approx 1{,}6 \times 10^{-11}$, soit $\mathrm{p}K_a \approx 10{,}8$ : le couple est celui d'une base faible (et d'un acide faible).

\begin{attention}[Erreur fréquente : la flèche simple]
L'erreur la plus fréquente au baccalauréat est d'écrire la réaction d'une amine avec l'eau avec une \textbf{flèche simple} $\ce{->}$, comme si la réaction était totale. C'est faux : les amines sont des \textbf{bases faibles}, leur réaction avec l'eau est un \textbf{équilibre} et doit être écrite avec la double flèche $\ce{<=>}$. Une solution d'amine à \SI{0,10}{mol.L^{-1}} ne contient donc \emph{pas} \SI{0,10}{mol.L^{-1}} d'ions \ce{HO-} : son pH est inférieur à 13 (il vaut environ 11,9 pour l'éthylamine).
\end{attention}

\courssub{Comparaison avec l'ammoniac : l'effet donneur des groupes alkyles}

Les amines sont-elles plus ou moins basiques que l'ammoniac ? L'expérience (et les valeurs de $\mathrm{p}K_a$) montre que, en solution aqueuse, les amines primaires et secondaires sont \textbf{plus basiques} que l'ammoniac.

\begin{propriete}[Effet inductif donneur des groupes alkyles]
Un groupe alkyle (comme \ce{CH3} ou \ce{C2H5}) est \cle{donneur d'électrons} par \textbf{effet inductif} : il « pousse » légèrement de la densité électronique vers l'atome d'azote. Le doublet non liant de l'azote devient alors \textbf{plus disponible} et l'ion alkylammonium formé est \textbf{stabilisé} par cette répartition de la charge positive. Conséquence : une amine capte plus facilement un proton que l'ammoniac.
\end{propriete}

\begin{center}
\begin{tabularx}{\linewidth}{|l|X|c|}
\hline
\rowcolor{popblueL} \textbf{Base} & \textbf{Couple acide/base} & $\mathrm{p}K_a$ du couple \\
\hline
Ammoniac \ce{NH3} & \ce{NH4+}/\ce{NH3} & 9,2 \\
\hline
Méthylamine \ce{CH3NH2} & \ce{CH3NH3+}/\ce{CH3NH2} & 10,7 \\
\hline
Éthylamine \ce{C2H5NH2} & \ce{C2H5NH3+}/\ce{C2H5NH2} & 10,8 \\
\hline
Diéthylamine \ce{(C2H5)2NH} & \ce{(C2H5)2NH2+}/\ce{(C2H5)2NH} & 11,0 \\
\hline
\end{tabularx}
\end{center}

\begin{aretenir}[À retenir]
Plus le $\mathrm{p}K_a$ du couple acide/base est \textbf{grand}, plus la \textbf{base} est \textbf{forte}. Les amines aliphatiques ($\mathrm{p}K_a \approx 10$ à $11$) sont donc des bases \emph{plus fortes} que l'ammoniac ($\mathrm{p}K_a = 9{,}2$), tout en restant des \emph{bases faibles} au sens où leur réaction avec l'eau reste un équilibre très limité.
\end{aretenir}

\courssub{Réaction des amines avec les acides : formation de sels d'ammonium}

Contrairement à la réaction avec l'eau, la réaction d'une amine avec un \textbf{acide fort} est \textbf{totale} : l'amine capte quantitativement le proton cédé par l'acide. Il se forme un \cle{sel d'alkylammonium}, composé ionique généralement \textbf{solide} à température ambiante et \textbf{très soluble dans l'eau}.

Avec l'acide chlorhydrique, l'éthylamine donne le \textbf{chlorure d'éthylammonium} :

\[ \ce{C2H5NH2 + HCl -> C2H5NH3+ + Cl-} \quad \text{soit} \quad \ce{C2H5NH2 + HCl -> C2H5NH3Cl} \]

Commentons cette équation :
\begin{itemize}
\item la réaction est \textbf{totale} (flèche simple $\ce{->}$), \textbf{rapide} et \textbf{exothermique} ;
\item elle se produit à température ambiante, sans catalyseur ;
\item le produit est un solide ionique blanc, le chlorure d'éthylammonium \ce{C2H5NH3Cl}, analogue au chlorure d'ammonium \ce{NH4Cl}.
\end{itemize}

De manière générale, en milieu aqueux acide :
\[ \ce{R-NH2 + H3O+ -> R-NH3+ + H2O} \]

Cette réaction est une véritable \textbf{réaction de dosage} possible : elle permet de déterminer la concentration d'une solution d'amine par titrage avec un acide fort, en suivant le pH.

\begin{exoresolu}[pH d'une solution d'éthylamine]
\textbf{Énoncé.} On dissout $n = \SI{4,5e-3}{mol}$ d'éthylamine \ce{C2H5NH2} dans de l'eau distillée de façon à obtenir $V = \SI{500}{mL}$ de solution. Le pH mesuré à \SI{25}{\celsius} vaut 11,9.
\begin{enumerate}
\item Écrire l'équation-bilan de la réaction de l'éthylamine avec l'eau.
\item Calculer la concentration $C$ de la solution en éthylamine apportée.
\item Calculer la concentration en ions hydroxyde \ce{HO-} dans la solution.
\item Montrer que la réaction est limitée et calculer son taux d'avancement final $\tau$.
\end{enumerate}
\tcblower
\textbf{Solution détaillée.}
\begin{enumerate}
\item L'éthylamine est une base de Brønsted ; avec l'eau :
\[ \ce{C2H5NH2 + H2O <=> C2H5NH3+ + HO-} \]
\item Concentration apportée :
\[ C = \frac{n}{V} = \frac{\num{4,5e-3}}{\num{0,500}} = \SI{9,0e-3}{mol.L^{-1}} \]
\item À \SI{25}{\celsius}, $\mathrm{p}K_e = 14$, donc :
\[ [\ce{H3O+}] = 10^{-\mathrm{pH}} = 10^{-11{,}9} = \SI{1,26e-12}{mol.L^{-1}} \]
\[ [\ce{HO-}] = \frac{K_e}{[\ce{H3O+}]} = \frac{10^{-14}}{\num{1,26e-12}} = \SI{7,9e-3}{mol.L^{-1}} \]
\item Si la réaction était totale, on aurait $[\ce{HO-}] = C = \SI{9,0e-3}{mol.L^{-1}}$. D'après le tableau d'avancement, à l'équilibre : $[\ce{HO-}] = x_f/V$ et $C = n/V$. Le taux d'avancement final est :
\[ \tau = \frac{x_f}{x_{max}} = \frac{[\ce{HO-}]}{C} = \frac{\num{7,9e-3}}{\num{9,0e-3}} \approx 0{,}88 \]
Ici $\tau \approx 88\,\% < 100\,\%$ : la réaction est bien \textbf{limitée}, l'éthylamine est une \textbf{base faible}. (Remarque : à cette concentration très diluée, l'avancement est élevé ; pour une solution plus concentrée, $\tau$ serait nettement plus faible, ce qui confirme le caractère d'équilibre régi par la loi d'Ostwald.)
\end{enumerate}
\end{exoresolu}

\courssub{Application pharmaceutique : les sels d'ammonium des médicaments}

\begin{savaistu}[Pourquoi tant de médicaments sont-ils des « chlorhydrates » ?]
De très nombreux médicaments contiennent une fonction amine : la \textbf{chloroquine} (antipaludique longtemps utilisé au Cameroun contre le paludisme), l'\textbf{amoxicilline}, la \textbf{quinine} extraite du quinquina, la \textbf{prométhazine}, etc. Or les amines sont souvent des molécules \emph{peu solubles dans l'eau} et parfois instables à l'air. Les laboratoires pharmaceutiques les transforment donc en \cle{sels d'ammonium} — le plus souvent des \textbf{chlorhydrates} — par réaction avec l'acide chlorhydrique :
\[ \ce{R-NH2 + HCl -> R-NH3+ Cl-} \]
Ces sels ioniques sont \textbf{beaucoup plus solubles dans l'eau}, donc mieux absorbés par l'organisme (le milieu sanguin et digestif est aqueux), plus stables et plus faciles à conserver sous nos climats chauds et humides. C'est ainsi que la chloroquine était vendue sous forme de \emph{chlorhydrate} ou de \emph{phosphate de chloroquine}. La prochaine fois que tu lis « chlorhydrate de... » sur une boîte de médicament, tu sauras qu'il s'agit d'un sel d'alkylammonium !
\end{savaistu}

\begin{exemplebox}[Solubilisation d'un principe actif]
Un sirop contre la toux contient un principe actif aminé, noté \ce{R-NH2}, peu soluble dans l'eau ($s \approx \SI{2}{g.L^{-1}}$). Pour le rendre soluble, le fabricant ajoute une quantité stœchiométrique d'acide chlorhydrique : la réaction totale
\[ \ce{R-NH2 + HCl -> R-NH3+ Cl-} \]
transforme intégralement l'amine en chlorhydrate, dont la solubilité dépasse \SI{200}{g.L^{-1}}. Le principe actif passe alors entièrement en solution : le sirop est limpide et le médicament est absorbable par l'organisme. Dans l'estomac, le milieu acide maintient le principe actif sous forme ionique \ce{R-NH3+}.
\end{exemplebox}

\begin{aretenir}[L'essentiel de la section]
\begin{itemize}
\item Une amine est une \textbf{base de Brønsted} : le doublet non liant de l'azote capte un proton \ce{H+}.
\item Avec l'eau, la réaction est un \textbf{équilibre} (base faible) : $\ce{RNH2 + H2O <=> RNH3+ + HO-}$ ; l'ion formé est un ion \textbf{alkylammonium}.
\item Les amines sont \textbf{plus basiques que l'ammoniac} à cause de l'effet inductif donneur des groupes alkyles ($\mathrm{p}K_a \approx 10$ à $11$ contre $9{,}2$).
\item Avec un acide fort, la réaction est \textbf{totale} et produit un \textbf{sel d'ammonium} soluble dans l'eau : $\ce{RNH2 + HCl -> RNH3+ Cl-}$.
\item Application majeure : les médicaments aminés sont formulés en \textbf{chlorhydrates} pour améliorer leur solubilité et leur absorption.
\end{itemize}
\end{aretenir}

Maintenant que tu maîtrises le caractère basique des amines — c'est-à-dire leur capacité à capter un proton — nous allons exploiter une autre conséquence de la présence du doublet non liant : leur \textbf{caractère nucléophile}, qui leur permet d'attaquer des atomes de carbone porteurs de charges partielles positives. C'est la clé de la réaction d'Hofmann et de la synthèse des amides, objets des deux prochaines sections.

\coursec{Caractère nucléophile et réaction d'Hofmann}

Dans la section précédente, nous avons vu que le doublet non liant de l'atome d'azote confère aux amines un \cle{caractère basique} : elles sont capables de capter un proton \ce{H+}. Mais ce même doublet peut faire bien plus que fixer un simple proton. Il peut attaquer un atome de carbone porteur d'une charge partielle positive, et c'est là toute la richesse de la chimie des amines : elles deviennent de véritables outils de construction moléculaire. C'est ce deuxième visage, le \cle{caractère nucléophile}, que nous allons maintenant explorer. C'est lui qui permet, en cascade, de passer de l'ammoniac aux amines primaires, secondaires, tertiaires, puis aux sels d'ammonium quaternaires : c'est la célèbre \cle{réaction d'Hofmann}.

\courssub{Le caractère nucléophile des amines}

\textbf{Intuition : un doublet en quête d'un site pauvre en électrons.}\par
Imagine un donneur généreux qui cherche à qui offrir son bien. Le doublet non liant de l'azote est ce donneur : il est disponible, riche en électrons, et il est naturellement attiré par les sites de la matière qui, eux, manquent d'électrons. Un proton \ce{H+} est un tel site (c'est la basicité), mais un atome de carbone rendu déficitaire en électrons par un atome d'halogène très électronégatif en est un autre. C'est cette seconde possibilité qui nous intéresse ici.

\begin{definition}[Nucléophile et électrophile]
Un \cle{nucléophile} (du grec \emph{nucleus}, noyau, et \emph{philos}, qui aime) est une espèce chimique \textbf{riche en électrons} (doublet non liant ou doublet d'une liaison multiple), attirée par les sites \textbf{déficitaires en électrons} (charges positives ou charges partielles positives $\delta^+$), auxquels elle peut céder un doublet pour former une nouvelle liaison covalente.
Un \cle{électrophile} est, à l'inverse, une espèce \textbf{pauvre en électrons}, attirée par les sites riches en électrons.
\end{definition}

\begin{propriete}[Caractère nucléophile des amines]
Grâce au \textbf{doublet non liant} porté par l'atome d'azote, les amines (et l'ammoniac \ce{NH3}) sont des \textbf{nucléophiles}. Elles ne sont \textbf{pas électrophiles} : l'azote, plus électronégatif que le carbone et l'hydrogène, porte une charge partielle négative et ne peut pas accepter de doublet supplémentaire.
L'atome ciblé par l'attaque nucléophile est typiquement un \textbf{carbone porteur d'un halogène} dans un dérivé halogéné \ce{R-X} (X = Cl, Br, I), car la liaison \ce{C-X} est polarisée : le carbone porte une charge partielle $\delta^+$.
\end{propriete}

\begin{attention}[Basicité et nucléophilie : deux rôles, un même doublet]
Ne confonds pas les deux rôles du doublet de l'azote :
\begin{itemize}
\item \textbf{Base} (au sens de Brønsted) : l'amine capte un \textbf{proton} \ce{H+}. C'est une question d'équilibre acide-base.
\item \textbf{Nucléophile} : l'amine attaque un \textbf{atome de carbone} déficitaire en électrons. C'est une question de formation d'une nouvelle liaison \ce{C-N}.
\end{itemize}
Une même amine joue les deux rôles, mais l'écriture de l'équation-bilan dépend du partenaire : un acide donne un transfert de proton, un dérivé halogéné donne une substitution nucléophile.
\end{attention}

\courssub{Substitution nucléophile d'une amine sur un dérivé halogéné}

\textbf{Le mécanisme, pas à pas.}\par
Considérons un dérivé halogéné, par exemple le bromométhane \ce{CH3Br}. La liaison \ce{C-Br} est polarisée car le brome est plus électronégatif que le carbone : le carbone porte une charge partielle $\delta^+$, le brome une charge partielle $\delta^-$. L'amine, riche en électrons, attaque ce carbone $\delta^+$ avec son doublet non liant. Simultanément, la liaison \ce{C-Br} se rompt : le brome part avec le doublet de la liaison (ion bromure \ce{Br-}). Un intermédiaire ionique (un ion alkylammonium) se forme, puis un proton est arraché à l'azote par une base présente dans le milieu (souvent une autre molécule d'amine ou d'ammoniac). Au bilan, l'halogène est parti sous forme d'acide halogénhydrique \ce{HX}, et l'azote s'est substitué à l'halogène : c'est une \cle{substitution nucléophile}.

\begin{popfigure}
\begin{center}
\begin{tikzpicture}[remember picture, baseline=(amine.base)]
\node[inner sep=0] (amine) {\chemfig{H_3C-\charge{45:3pt=$\scriptstyle\oplus$}{N}(-[2]H)(-[6]CH_2CH_3)-H}};
\node[inner sep=0, right=1.5cm of amine] (halo) {\chemfig{@{C}C(-[2]H)(-[6]H)(-[0]@{Br}Br)(-[4]H)}};
\draw[-{Stealth[length=3mm]}, popink, thick, shorten <=2pt, shorten >=2pt] 
  (amine.east) .. controls +(0:1.2) and +(180:1.2) .. (halo.west);
\draw[-{Stealth[length=2mm]}, poporange, thick, shorten <=1pt] 
  (halo.west) .. controls +(180:0.5) and +(90:0.5) .. ($(halo.west)+(-0.4,-0.4)$) node[below left, font=\scriptsize, poporange] {\ce{Br-}};
\end{tikzpicture}
\end{center}
\vspace{0.5em}
\begin{center}
\setchemfig{atom sep=2.2em}
\chemfig{CH_3-\charge{45:3pt=$\scriptstyle\oplus$}{N}(-[2]H)(-[6]H)(-[0]CH_2CH_3)-H}
\quad \ce{+ Br-} \quad \ce{->} \quad
\chemfig{CH_3-NH(-[2]H)-[0]CH_2CH_3}
\quad \ce{+ HBr}
\end{center}
\legende{Mécanisme simplifié de la substitution nucléophile de l'éthylamine sur le bromométhane : le doublet de l'azote (flèche rose) attaque le carbone $\delta^+$ ; la liaison \ce{C-Br} se rompt (flèche orange), le brome part sous forme \ce{Br-}, puis la perte d'un proton \ce{H+} donne la $N$-méthyléthanamine et \ce{HBr}.}
\end{popfigure}

\begin{methode}[Écrire l'équation-bilan de la réaction d'une amine sur un dérivé halogéné]
\begin{enumerate}
\item \textbf{Repérer le carbone porteur de l'halogène} dans \ce{R-X} : c'est lui, chargé $\delta^+$, qui subit l'attaque.
\item \textbf{Faire attaquer l'azote} par son doublet non liant : le groupe alkyle \ce{R} se greffe sur l'azote, en remplacement d'un hydrogène de l'amine (ou de l'ammoniac).
\item \textbf{Faire partir l'halogène} : il emporte le doublet de la liaison \ce{C-X}, puis capte le proton \ce{H+} libéré par l'azote pour former \ce{HX}.
\item \textbf{Vérifier la conservation des atomes} : à gauche comme à droite, mêmes nombres de C, H, N, X. En particulier, l'amine de départ perd exactement \textbf{un} hydrogène sur l'azote, qui se retrouve dans \ce{HX}.
\item \textbf{Préciser les conditions} : réaction lente, souvent réalisée à chaud, en solution alcoolique, sous pression si l'ammoniac est gazeux.
\end{enumerate}
\end{methode}

\begin{exemplebox}[Éthylamine sur le bromométhane]
L'éthylamine \ce{CH3CH2NH2} réagit avec le bromométhane \ce{CH3Br}. Le groupe méthyle se greffe sur l'azote, le brome part avec un proton :
\[ \ce{CH3CH2NH2 + CH3Br -> CH3CH2NHCH3 + HBr} \]
Le produit est la \textbf{$N$-méthyléthanamine} \ce{CH3CH2-NH-CH3}, une amine \textbf{secondaire}. Vérification : à gauche 3 C, 10 H, 1 N, 1 Br ; à droite 3 C, 10 H, 1 N, 1 Br. L'équation est équilibrée.
\end{exemplebox}

\courssub{La réaction d'Hofmann : une alkylation en cascade}

\begin{definition}[Réaction d'Hofmann]
La \cle{réaction d'Hofmann} est l'\textbf{alkylation d'une amine (ou de l'ammoniac) par un dérivé halogéné} \ce{R-X}, conduisant à une amine de \textbf{classe supérieure} : l'ammoniac donne une amine primaire, qui donne une amine secondaire, qui donne une amine tertiaire, qui donne enfin un \textbf{sel d'ammonium quaternaire}.
\end{definition}

\textbf{Pourquoi une cascade ?}\par
Voici le point capital, et le plus déroutant pour un élève : le produit de la réaction est \emph{lui-même nucléophile}. Quand l'ammoniac \ce{NH3} attaque le bromoéthane, il se forme l'éthanamine \ce{CH3CH2NH2}, une amine primaire qui possède \emph{toujours} un doublet non liant sur l'azote. Elle peut donc attaquer une deuxième molécule de bromoéthane, puis une troisième. La réaction ne s'arrête que lorsque l'azote a perdu tous ses hydrogènes et porte quatre groupes alkyles : on obtient alors un \textbf{ion ammonium quaternaire}, qui n'a plus de doublet disponible et ne peut plus réagir.

\begin{popfigure}
\begin{center}
\begin{tikzpicture}[
  node distance=0.55cm and 0.9cm,
  box/.style={draw=popblue, thick, rounded corners=3pt, fill=popblueL, align=center, minimum height=0.9cm, inner sep=4pt, font=\small},
  arr/.style={-{Stealth[length=3mm]}, thick, poporange},
  lbl/.style={font=\scriptsize\itshape, text=popdark}
]
\node[box] (nh3) {Ammoniac\\ \ce{NH3}};
\node[box, right=of nh3] (prim) {Amine primaire\\ \ce{R-NH2}};
\node[box, right=of prim] (sec) {Amine secondaire\\ \ce{R2NH}};
\node[box, right=of sec] (tert) {Amine tertiaire\\ \ce{R3N}};
\node[box, right=of tert] (quat) {Ammonium quaternaire\\ \ce{[R4N]+ X-}};
\draw[arr] (nh3) -- node[lbl, above]{\ce{+ R-X}} node[lbl, below]{\ce{- HX}} (prim);
\draw[arr] (prim) -- node[lbl, above]{\ce{+ R-X}} node[lbl, below]{\ce{- HX}} (sec);
\draw[arr] (sec) -- node[lbl, above]{\ce{+ R-X}} node[lbl, below]{\ce{- HX}} (tert);
\draw[arr] (tert) -- node[lbl, above]{\ce{+ R-X}} (quat);
\end{tikzpicture}
\end{center}
\legende{Cascade de la réaction d'Hofmann : chaque étape greffe un groupe alkyle \ce{R} sur l'azote. À la dernière étape, l'azote n'a plus d'hydrogène à perdre : il n'y a pas de \ce{HX} formé, mais un sel d'ammonium quaternaire.}
\end{popfigure}

\textbf{Les équations-bilans successives avec le bromoéthane.}\par
Prenons l'exemple complet de l'ammoniac et du bromoéthane \ce{CH3CH2Br}, chauffés en solution alcoolique.

\textbf{Étape 1 — formation de l'amine primaire (éthanamine) :}
\[ \ce{NH3 + CH3CH2Br -> CH3CH2NH2 + HBr} \]

\textbf{Étape 2 — formation de l'amine secondaire (diéthylamine) :}
\[ \ce{CH3CH2NH2 + CH3CH2Br -> (CH3CH2)2NH + HBr} \]

\textbf{Étape 3 — formation de l'amine tertiaire (triéthylamine) :}
\[ \ce{(CH3CH2)2NH + CH3CH2Br -> (CH3CH2)3N + HBr} \]

\textbf{Étape 4 — formation du sel d'ammonium quaternaire (bromure de tétraéthylammonium) :}
\[ \ce{(CH3CH2)3N + CH3CH2Br -> [(CH3CH2)4N]+ Br-} \]

Remarque bien la particularité de la dernière étape : l'azote de l'amine tertiaire n'a \textbf{plus d'hydrogène}. Il attaque avec son doublet, forme une quatrième liaison \ce{C-N}, devient chargé positivement, et le bromure \ce{Br-} assure la neutralité électrique du sel. Il n'y a donc \textbf{pas de \ce{HBr} formé} à cette étape.

\begin{attention}[Erreur fréquente : la réaction ne s'arrête pas à l'amine primaire]
Le piège classique du baccalauréat : on te demande le produit de l'action de l'ammoniac sur un dérivé halogéné, et tu écris uniquement l'amine primaire en la présentant comme le produit unique. C'est faux : l'amine primaire formée \textbf{reste nucléophile} et réagit à nouveau sur le dérivé halogéné. En pratique, on obtient un \textbf{mélange} d'amine primaire, secondaire, tertiaire et de sel quaternaire. En exercice, précise toujours la classe de l'amine obtenue et envisage la poursuite de la réaction si le dérivé halogéné est en excès.
\end{attention}

\textbf{Limite de la réaction et contrôle des proportions.}\par
Cette polyalkylation est à la fois une faiblesse et une richesse. Faiblesse, car si l'on mélange \ce{NH3} et \ce{R-X} en proportions stœchiométriques, on obtient un mélange de quatre produits difficiles à séparer. Richesse, car en jouant sur les \textbf{proportions relatives des réactifs}, on oriente le résultat :
\begin{itemize}
\item avec un \textbf{large excès d'ammoniac} (par exemple 10 mol de \ce{NH3} pour 1 mol de \ce{R-X}), chaque molécule de dérivé halogéné a statistiquement toutes les chances de rencontrer \ce{NH3} plutôt qu'une amine déjà formée : on obtient \textbf{majoritairement l'amine primaire} ;
\item avec un \textbf{excès de dérivé halogéné}, la cascade va jusqu'au bout : on obtient majoritairement le \textbf{sel d'ammonium quaternaire}.
\end{itemize}

\courssub{Exemple chiffré et application industrielle}

\begin{exoresolu}[Orientation de la réaction d'Hofmann]
\textbf{Énoncé.} On fait réagir \SI{0,10}{mol} de bromoéthane \ce{CH3CH2Br} avec \SI{1,0}{mol} d'ammoniac en solution alcoolique, à chaud.
\begin{enumerate}
\item Écrire l'équation-bilan de la première étape de la réaction.
\item Quel produit obtient-on majoritairement ? Justifier.
\item Calculer la masse maximale d'éthanamine ($M = \SI{45}{g.mol^{-1}}$) que l'on peut espérer recueillir.
\item Que se passerait-il si l'on utilisait \SI{0,50}{mol} de bromoéthane pour \SI{0,10}{mol} d'ammoniac ?
\end{enumerate}
\tcblower
\textbf{Solution détaillée.}
\begin{enumerate}
\item Première étape (substitution nucléophile de \ce{NH3} sur le carbone $\delta^+$ du bromoéthane) :
\[ \ce{NH3 + CH3CH2Br -> CH3CH2NH2 + HBr} \]
\item L'ammoniac est en large excès : $n(\ce{NH3})/n(\ce{CH3CH2Br}) = 1{,}0/0{,}10 = 10$. Chaque molécule de bromoéthane a dix fois plus de chances de rencontrer une molécule d'ammoniac qu'une molécule d'éthanamine déjà formée. La cascade s'arrête donc essentiellement à la première étape : on obtient \textbf{majoritairement l'éthanamine} \ce{CH3CH2NH2} (amine primaire).
\item Le bromoéthane est le réactif limitant : $n_{\max}(\ce{CH3CH2NH2}) = \SI{0,10}{mol}$.
\[ m_{\max} = n \times M = 0{,}10 \times 45 = \SI{4,5}{g} \]
\item Avec \SI{0,50}{mol} de bromoéthane pour \SI{0,10}{mol} d'ammoniac, le dérivé halogéné est en excès (5 mol pour 1 mol de \ce{NH3}). Les amines formées réagissent encore et encore : la cascade se poursuit jusqu'au \textbf{bromure de tétraéthylammonium} \ce{[(CH3CH2)4N]+ Br-}, produit majoritaire. On remarque qu'il faut théoriquement 4 mol de \ce{R-X} par mole de \ce{NH3} pour quaterneriser tout l'azote ; ici les 5 mol disponibles suffisent largement.
\end{enumerate}
\end{exoresolu}

\begin{savaistu}[Les amines dans l'industrie et le médicament]
L'alkylation des amines par des dérivés halogénés est une voie industrielle classique pour fabriquer des amines de synthèse : solvants, agents de surface, et surtout \textbf{intermédiaires pharmaceutiques}. De nombreux médicaments (antihistaminiques, anesthésiques locaux comme la lidocaïne) contiennent une amine secondaire ou tertiaire construite par ce type de réaction. Les \textbf{sels d'ammonium quaternaires}, produits ultimes de la cascade d'Hofmann, sont eux-mêmes très utiles : ce sont des \textbf{antiseptiques} (le chlorure de benzalkonium, présent dans de nombreux désinfectants vendus en pharmacie à Douala ou Yaoundé) et des agents adoucissants textiles. La réaction porte le nom du chimiste allemand \emph{August Wilhelm von Hofmann}, qui l'étudia systématiquement au XIX\textsuperscript{e} siècle.
\end{savaistu}

\begin{exercice}[À toi de jouer]
On chauffe en solution alcoolique la méthanamine \ce{CH3NH2} avec un excès d'iodométhane \ce{CH3I}.
\begin{enumerate}
\item Écrire les équations-bilans des trois étapes conduisant au sel d'ammonium quaternaire final, en nommant chaque produit.
\item Préciser la classe de chaque amine intermédiaire.
\item Pourquoi n'y a-t-il pas de \ce{HI} formé à la dernière étape ?
\end{enumerate}
\end{exercice}

\begin{corrige}[Exercice : méthanamine et iodométhane en excès]
\begin{enumerate}
\item Étape 1 : \ce{CH3NH2 + CH3I -> (CH3)2NH + HI} : diméthylamine.
Étape 2 : \ce{(CH3)2NH + CH3I -> (CH3)3N + HI} : triméthylamine.
Étape 3 : \ce{(CH3)3N + CH3I -> [(CH3)4N]+ I-} : iodure de tétraméthylammonium.
\item La méthanamine de départ est primaire, la diméthylamine est secondaire, la triméthylamine est tertiaire.
\item À la dernière étape, l'azote de la triméthylamine ne porte plus aucun atome d'hydrogène : il ne peut donc pas libérer de proton \ce{H+}. L'ion iodure \ce{I-} reste simplement associé à l'ion ammonium quaternaire chargé positivement pour former le sel.
\end{enumerate}
\end{corrige}

\begin{aretenir}[L'essentiel de la section]
\begin{itemize}
\item Le doublet non liant de l'azote fait des amines des \textbf{nucléophiles} : elles attaquent les carbones $\delta^+$ des dérivés halogénés \ce{R-X}.
\item La \textbf{réaction d'Hofmann} est l'alkylation d'une amine (ou de \ce{NH3}) par \ce{R-X} : à chaque étape, un groupe alkyle remplace un hydrogène de l'azote et \ce{HX} se forme.
\item La réaction se poursuit \textbf{en cascade} : amine primaire $\rightarrow$ secondaire $\rightarrow$ tertiaire $\rightarrow$ \textbf{sel d'ammonium quaternaire} \ce{[R4N]+ X-} (sans \ce{HX} à la dernière étape).
\item Un \textbf{excès d'ammoniac} favorise l'amine primaire ; un \textbf{excès de dérivé halogéné} conduit au sel quaternaire.
\end{itemize}
\end{aretenir}

\coursec{Action sur les chlorures d'acyle : synthèse des amides}

Dans la section précédente, nous avons exploité le caractère \cle{nucléophile} des amines dans la réaction d'Hofmann : le doublet libre de l'azote attaque un carbone porteur d'un halogène. Mais le carbone d'un dérivé halogéné n'est pas le seul site électrophile que l'azote d'une amine peut attaquer. Il existe un autre électrophile, beaucoup plus réactif encore : le carbone du groupe \cle{carbonyle} d'un \cle{chlorure d'acyle}. Cette réaction, rapide et totale, est la voie royale vers une famille de composés d'une importance capitale : les \cle{amides}, présents dans les protéines de notre organisme et dans de très nombreux médicaments, dont le fameux paracétamol que l'on trouve dans toutes les pharmacies de Yaoundé à Maroua.

\courssub{Le chlorure d'acyle : un électrophile puissant}

\textbf{Structure du chlorure d'acyle}\par

Un \cle{chlorure d'acyle} est un composé de formule générale $R-COCl$, dans lequel un groupe acyle $R-CO-$ est lié à un atome de chlore. On le nomme en remplaçant la terminaison « oïque » de l'acide carboxylique correspondant par « yle » : \ce{CH3-COCl} est le \textbf{chlorure d'éthanoyle} (dérivé de l'acide éthanoïque), \ce{C6H5-COCl} le chlorure de benzoyle.

\begin{center}
\chemfig{CH_3-C(=[2]O)(-[6]Cl)}
\hspace{2cm}
\chemfig{C_6H_5-C(=[2]O)(-[6]Cl)}
\end{center}

\textbf{Pourquoi le carbone du carbonyle est-il électrophile ?}\par

Deux effets se conjuguent pour appauvrir ce carbone en électrons :
\begin{itemize}
\item l'atome d'\textbf{oxygène}, très électronégatif, attire vers lui les électrons de la double liaison $C=O$ ;
\item l'atome de \textbf{chlore}, également électronégatif, attire les électrons de la liaison $C-Cl$.
\end{itemize}
Le carbone du carbonyle porte donc une charge partielle $\delta^+$ importante : c'est un site \cle{électrophile} remarquable, encore plus réactif que celui d'un dérivé halogéné. De plus, le chlore constitue un excellent \emph{groupe partant} : la liaison $C-Cl$ se rompt facilement.

\begin{attention}[Un électrophile plus réactif que le dérivé halogéné]
Ne confonds pas les deux situations. Dans la réaction d'Hofmann, l'amine attaque le carbone d'un alkyle $R-X$ : la réaction est lente et conduit à des mélanges d'amines. Avec un chlorure d'acyle, l'attaque a lieu sur le carbone du carbonyle, bien plus électrophile : la réaction est \textbf{rapide et totale}, et elle s'arrête proprement au stade de l'amide.
\end{attention}

\courssub{La réaction : équation-bilan et mécanisme simplifié}

\textbf{Intuition}\par

Imagine une rencontre entre deux partenaires complémentaires : le chlorure d'acyle apporte un groupe acyle $R-CO-$ et un chlore « encombrant » ; l'amine apporte son azote nucléophile et, sur cet azote, au moins un hydrogène « encombrant ». La réaction consiste simplement à \emph{éliminer ensemble} le chlore et l'hydrogène (sous forme de \ce{HCl}), puis à \emph{relier} le carbonyle à l'azote. Le nouveau pont $-CO-NH-$ créé entre les deux partenaires est la \cle{liaison amide}.

\textbf{Équation-bilan générale (amine primaire)}\par

\[ \ce{\text{R}-COCl + \text{R'}-NH2 -> \text{R}-CO-NH-\text{R'} + HCl} \]

\begin{definition}[Amide]
Un \cle{amide} est un composé organique de formule générale $R-CO-NR'R''$, où $R'$, $R''$ sont des atomes d'hydrogène ou des groupes alkyles ou aryles. On l'obtient par action d'une amine (ou de l'ammoniac) sur un \textbf{chlorure d'acyle} (ou sur un \textbf{anhydride d'acide}). Le groupe caractéristique $-CO-N-$ est appelé \textbf{groupe amide} ; la liaison $C-N$ qui le constitue est la \textbf{liaison amide}, dite \emph{liaison peptidique} lorsqu'elle relie deux acides aminés dans une protéine.
\end{definition}

\begin{methode}[Écrire l'équation-bilan en trois étapes]
\begin{enumerate}
\item \textbf{Identifier les partenaires} : repérer le groupe acyle $R-CO-$ du chlorure d'acyle et le reste $R'$ porté par l'azote de l'amine.
\item \textbf{Éliminer HCl} : le chlore du chlorure d'acyle s'élimine avec \emph{un} hydrogène fixé sur l'azote, sous forme de chlorure d'hydrogène \ce{HCl}.
\item \textbf{Relier} : écrire la nouvelle liaison entre le carbone du carbonyle et l'azote : $R-CO-NH-R'$ (amine primaire) ou $R-CO-NR'R''$ (amine secondaire). Vérifier que l'équation est équilibrée en atomes.
\end{enumerate}
\end{methode}

\begin{exemplebox}[Chlorure d'éthanoyle + méthylamine]
\[ \ce{CH3-COCl + CH3-NH2 -> CH3-CO-NH-CH3 + HCl} \]
Le produit est le \textbf{N-méthyléthanamide} : le préfixe « N- » indique que le groupe méthyle est porté par l'\textbf{azote}, et non par la chaîne carbonée. La réaction est vive, exothermique, et se déroule à température ordinaire, sans catalyseur.
\end{exemplebox}

\begin{popfigure}
\begin{center}
\setchemfig{atom sep=2.2em}
\chemfig{CH_3-C(=[2]O)(-[6]Cl)}
\hspace{0.8em}+\hspace{0.8em}
\chemfig{CH_3-NH_2}
\hspace{0.8em}\chemfig{[,0.6]->}\hspace{0.8em}
\chemfig{CH_3-C(=[2]O)(-[6]@{N}N(-[6]H)-[0]CH_3)}
\hspace{0.8em}+\hspace{0.8em}
\chemfig{HCl}
\chemmove{\draw[poporange,line width=1.2pt,rounded corners=4pt]
([shift={(-0.25,-0.35)}]N.south west) rectangle ([shift={(0.25,0.55)}]N.north east);}
\end{center}
\legende{Réaction du chlorure d'éthanoyle sur la méthylamine. L'atome d'azote (encadré en orange) s'est inséré entre le carbonyle et le groupe méthyle : la liaison $C-N$ créée est la \textbf{liaison amide}. Le chlore et un hydrogène de l'azote ont été éliminés sous forme de \ce{HCl}.}
\end{popfigure}

\textbf{Caractéristiques de la réaction}\par

\begin{propriete}[Une réaction rapide et totale]
L'action d'une amine sur un chlorure d'acyle est :
\begin{itemize}
\item \textbf{rapide} : elle est souvent vive, exothermique, dès la température ordinaire ;
\item \textbf{totale} : elle n'est pas limitée par un équilibre, contrairement à l'action d'une amine sur un \textbf{acide carboxylique}, qui ne donne qu'un \emph{sel d'ammonium} $R-COO^- \, R'NH_3^+$ (réaction acido-basique réversible), lequel ne se déshydrate en amide que par chauffage intense.
\end{itemize}
C'est pourquoi, en synthèse industrielle, on préfère toujours passer par le chlorure d'acyle (ou l'anhydride) plutôt que par l'acide carboxylique lui-même.
\end{propriete}

\begin{attention}[Le piège du sel d'ammonium]
Si tu écris $\mathrm{R{-}COOH + R'NH_2 \longrightarrow R{-}CO{-}NH{-}R' + H_2O}$ à température ordinaire, c'est \textbf{faux} : l'acide carboxylique réagit avec l'amine par transfert de proton (réaction acide-base) pour donner un \emph{carboxylate d'alkylammonium}. L'amide ne s'obtient à partir de l'acide qu'en chauffant fortement ce sel. Avec le chlorure d'acyle, en revanche, l'amide se forme directement et totalement.
\end{attention}

\courssub{Influence de la classe de l'amine}

La condition essentielle pour former un amide est la présence d'\textbf{au moins un atome d'hydrogène sur l'azote}, hydrogène qui sera éliminé avec le chlore sous forme de \ce{HCl}. Examinons les trois classes.

\textbf{Amine primaire $R'-NH_2$ : amide N-substitué}\par

L'azote possède deux hydrogènes ; un seul est éliminé. On obtient un amide \textbf{N-substitué} $R-CO-NH-R'$ (un seul groupe alkyle sur l'azote) :
\[ \ce{CH3-COCl + C2H5-NH2 -> CH3-CO-NH-C2H5 + HCl} \]
Le produit est le \textbf{N-éthyléthanamide}.

\textbf{Amine secondaire $R'-NH-R''$ : amide N,N-disubstitué}\par

L'azote ne possède plus qu'\emph{un} hydrogène, mais cela suffit. On obtient un amide \textbf{N,N-disubstitué} $R-CO-NR'R''$ :
\[ \ce{CH3-COCl + (CH3)2NH -> CH3-CO-N(CH3)2 + HCl} \]
Le produit est le \textbf{N,N-diméthyléthanamide}. Remarque que l'azote de l'amide formé ne porte \emph{plus aucun hydrogène} : la réaction ne peut pas aller plus loin, elle s'arrête proprement là.

\textbf{Amine tertiaire $NR'R''R'''$ : pas d'amide}\par

L'azote ne porte \textbf{aucun hydrogène} : l'élimination de \ce{HCl} est impossible. Il ne se forme \textbf{pas d'amide}. On observe au mieux un intermédiaire ionique instable (un chlorure d'acylammonium $[R-CO-NR'_3]^+ Cl^-$) qui se décompose.

\begin{attention}[Erreur fréquente au Bac]
Faire réagir une amine tertiaire comme la N,N-diméthylméthanamine \ce{N(CH3)3} avec un chlorure d'acyle \emph{en espérant un amide} est une erreur classique : \textbf{sans hydrogène sur l'azote, l'amide ne peut pas se former}. Avant d'écrire l'équation, compte toujours les hydrogènes portés par l'azote : il en faut au moins un. Astuce : ammoniac \ce{NH3} donne un amide non substitué $R-CO-NH_2$ ; amine primaire $\to$ amide N-substitué ; amine secondaire $\to$ amide N,N-disubstitué ; amine tertiaire $\to$ rien.
\end{attention}

\begin{center}
\begin{tabularx}{\linewidth}{|l|X|X|}
\hline
\rowcolor{popblueL}
\textbf{Réactif azoté} & \textbf{H sur l'azote} & \textbf{Produit avec $R-COCl$} \\
\hline
Ammoniac \ce{NH3} & 3 & Amide non substitué $R-CO-NH_2$ \\
\hline
Amine primaire $R'NH_2$ & 2 & Amide N-substitué $R-CO-NH-R'$ \\
\hline
Amine secondaire $R'NHR''$ & 1 & Amide N,N-disubstitué $R-CO-NR'R''$ \\
\hline
Amine tertiaire $NR'R''R'''$ & 0 & \textbf{Pas d'amide} \\
\hline
\end{tabularx}
\end{center}

\courssub{La liaison amide : des protéines aux médicaments}

La liaison amide $-CO-NH-$ possède une propriété remarquable : par \emph{mésomérie}, le doublet libre de l'azote se délocalise vers le carbonyle, ce qui confère à la liaison $C-N$ un \emph{caractère partiel de double liaison}. Le groupe amide est donc \textbf{plan} et \textbf{rigide}, et il est chimiquement très stable. C'est précisément cette stabilité qui en fait le ciment du monde vivant.

\begin{savaistu}[La liaison peptidique, fil de la vie]
Dans les \textbf{protéines} — les enzymes qui digèrent ton ndolé, l'hémoglobine qui transporte ton oxygène, les anticorps qui te défendent —, les acides aminés sont enchaînés les uns aux autres par des liaisons amides, appelées alors \textbf{liaisons peptidiques}. Une protéine est un \emph{polyamide} naturel. Le nylon, lui, est un polyamide de synthèse fabriqué par le même type de chimie ! Et côté médicaments, la fonction amide est partout : paracétamol, pénicillines, lidocaïne (anesthésique local)...
\end{savaistu}

\courssub{Application : la synthèse du paracétamol}

Le \cle{paracétamol} (ou acétaminophène) est l'antalgique et antipyrétique le plus consommé au monde. Sa molécule, le \emph{N-(4-hydroxyphényl)éthanamide}, possède une fonction \textbf{amide} et une fonction \textbf{phénol}.

Industriellement, on le prépare en faisant réagir le \textbf{para-aminophénol} (qui contient une amine primaire $-NH_2$ en position \emph{para} d'un groupe $-OH$ sur le cycle benzénique) avec un dérivé acétylant : l'anhydride éthanoïque en pratique industrielle, mais le \textbf{chlorure d'éthanoyle} convient parfaitement et illustre notre réaction :

\[ \ce{HO-C6H4-NH2 + CH3-COCl -> HO-C6H4-NH-CO-CH3 + HCl} \]

Seul le groupe $-NH_2$ réagit (l'azote est bien plus nucléophile que l'oxygène du phénol) : la fonction phénol $-OH$ est conservée dans le produit. La réaction est rapide et totale à température ordinaire.

\begin{popfigure}
\begin{center}
\setchemfig{atom sep=2em}
\chemfig{CH_3-C(=[2]O)(-[:-30]@{Na}N(-[:-90]H)-[:30]**6(---(-OH)---))}
\chemmove{\draw[popgreen,line width=1.3pt,rounded corners=5pt]
([shift={(-1.6,-1.5)}]Na.south west) rectangle ([shift={(1.7,1.6)}]Na.north east);}
\end{center}
\legende{Molécule de paracétamol. Le groupe \textbf{amide} $-CO-NH-$ (encadré en vert) relie le groupe acétyle \ce{CH3-CO}- au cycle aromatique ; le groupe $-OH$ en position \emph{para} est la fonction phénol. C'est le groupe amide qui est formé lors de l'action du chlorure d'éthanoyle (ou de l'anhydride éthanoïque) sur le para-aminophénol.}
\end{popfigure}

\begin{savaistu}[Le paracétamol au Cameroun]
Vendu sous de nombreux noms commerciaux dans toutes les pharmacies et les kiosques à médicaments du pays, le paracétamol est le médicament le plus consommé au Cameroun, contre la fièvre et les douleurs, notamment lors des crises de paludisme. Sa fonction amide, obtenue industriellement par le type de réaction étudié dans cette leçon, lui confère sa stabilité : il se conserve longtemps, même sous les climats chauds et humides de Douala. Un bel exemple de la chimie organique au service de la santé publique !
\end{savaistu}

\courssub{Complément Bac : action des amines sur les anhydrides d'acide}

Les \cle{anhydrides d'acide}, de formule générale $R-CO-O-CO-R'$, possèdent eux aussi un carbone de carbonyle très électrophile. Ils réagissent avec les amines de la même manière que les chlorures d'acyle, mais le groupe éliminé n'est plus \ce{HCl} : c'est un \textbf{acide carboxylique}.

\[ \ce{(CH3CO)2O + \text{R'}-NH2 -> CH3-CO-NH-\text{R'} + CH3COOH} \]

Avec l'anhydride éthanoïque et une amine primaire, on obtient l'amide N-substitué et l'acide éthanoïque. La réaction est également \textbf{rapide et totale}.

\begin{center}
\begin{tabularx}{\linewidth}{|l|X|X|X|}
\hline
\rowcolor{popblueL}
\textbf{Réactif acylant} & \textbf{Formule} & \textbf{Sous-produit} & \textbf{Remarque} \\
\hline
Chlorure d'acyle & $R-COCl$ & \ce{HCl} & Très réactif, mais corrosif et fumant \\
\hline
Anhydride d'acide & $(R-CO)_2O$ & $R-COOH$ & Moins agressif ; préféré en industrie (paracétamol) \\
\hline
Acide carboxylique & $R-COOH$ & \ce{H2O} & Sel d'ammonium d'abord ; amide seulement à chaud \\
\hline
\end{tabularx}
\end{center}

\begin{exoresolu}[Synthèse d'un amide : équation et quantités de matière]
On fait réagir le chlorure d'éthanoyle \ce{CH3-COCl} avec l'éthanamine \ce{C2H5-NH2}.
\begin{enumerate}
\item Écrire l'équation-bilan de la réaction et nommer le produit organique obtenu.
\item On utilise $n_1 = \SI{0,20}{mol}$ de chlorure d'éthanoyle et $n_2 = \SI{0,15}{mol}$ d'éthanamine. Déterminer le réactif limitant et la quantité d'amide formé.
\item Calculer la masse d'amide obtenue. Données : $M(C)=\SI{12}{g.mol^{-1}}$, $M(H)=\SI{1}{g.mol^{-1}}$, $M(O)=\SI{16}{g.mol^{-1}}$, $M(N)=\SI{14}{g.mol^{-1}}$.
\end{enumerate}
\tcblower
\textbf{1. Équation-bilan.} L'éthanamine est une amine \emph{primaire} : son azote porte deux hydrogènes, l'amide peut donc se former. Le chlore s'élimine avec un hydrogène de l'azote sous forme de \ce{HCl} :
\[ \ce{CH3-COCl + C2H5-NH2 -> CH3-CO-NH-C2H5 + HCl} \]
Le produit est le \textbf{N-éthyléthanamide}. La réaction est rapide et totale à température ordinaire.

\textbf{2. Réactif limitant.} L'équation montre que les réactifs réagissent \emph{mole à mole} (stœchiométrie $1:1$). Or $n_2 = \SI{0,15}{mol} < n_1 = \SI{0,20}{mol}$ : l'\textbf{éthanamine est le réactif limitant}. La réaction étant totale :
\[ n_{\text{amide}} = n_2 = \SI{0,15}{mol} \]
et il reste $0{,}20 - 0{,}15 = \SI{0,05}{mol}$ de chlorure d'éthanoyle en excès.

\textbf{3. Masse d'amide.} La formule brute du N-éthyléthanamide \ce{CH3-CO-NH-C2H5} est \ce{C4H9NO} :
\[ M = 4 \times 12 + 9 \times 1 + 14 + 16 = \SI{87}{g.mol^{-1}} \]
\[ m = n \times M = 0{,}15 \times 87 = \SI{13,05}{g} \approx \SI{13,1}{g} \]
On obtient donc environ $\SI{13,1}{g}$ de N-éthyléthanamide.
\end{exoresolu}

\begin{exercice}[À toi de jouer]
\begin{enumerate}
\item Écrire l'équation-bilan de l'action du chlorure d'éthanoyle sur la diéthylamine \ce{(C2H5)2NH} et nommer l'amide obtenu.
\item Le chlorure d'éthanoyle peut-il réagir avec la N,N-diméthylméthanamine \ce{N(CH3)3} pour donner un amide ? Justifier.
\item Écrire l'équation-bilan de la préparation industrielle du paracétamol à partir du para-aminophénol et de l'anhydride éthanoïque \ce{(CH3-CO)2O}.
\end{enumerate}
\end{exercice}

\begin{corrige}[Exercice n]
\textbf{1.} La diéthylamine est une amine \emph{secondaire} (un seul H sur N) :
\[ \ce{CH3-COCl + (C2H5)2NH -> CH3-CO-N(C2H5)2 + HCl} \]
Le produit est le \textbf{N,N-diéthyléthanamide} (amide N,N-disubstitué).

\textbf{2.} \textbf{Non.} La N,N-diméthylméthanamine est une amine \emph{tertiaire} : son azote ne porte \textbf{aucun atome d'hydrogène}. Or la formation de l'amide exige l'élimination de \ce{HCl}, donc d'un hydrogène fixé sur l'azote. Aucun amide ne peut se former.

\textbf{3.} L'anhydride éthanoïque acétyle le groupe $-NH_2$ du para-aminophénol ; le sous-produit est l'acide éthanoïque :
\[ \ce{HO-C6H4-NH2 + (CH3-CO)2O -> HO-C6H4-NH-CO-CH3 + CH3-COOH} \]
La réaction est rapide et totale ; la fonction phénol $-OH$ n'est pas touchée.
\end{corrige}

\begin{aretenir}[L'essentiel de la section]
\begin{itemize}
\item Un \textbf{chlorure d'acyle} $R-COCl$ possède un carbone de carbonyle très \textbf{électrophile} et un chlore excellent groupe partant.
\item Avec une amine, il donne un \textbf{amide} par élimination de \ce{HCl} : $\ce{{R-COCl} + {R'NH2} -> {R-CO-NH-R'} + HCl}$. La réaction est \textbf{rapide et totale} (contrairement à l'action sur un acide carboxylique, qui s'arrête au sel d'ammonium).
\item Il faut \textbf{au moins un H sur l'azote} : amine primaire $\to$ amide N-substitué ; amine secondaire $\to$ amide N,N-disubstitué ; amine tertiaire $\to$ \textbf{pas d'amide}.
\item La \textbf{liaison amide} (peptidique) est le ciment des protéines et se retrouve dans de nombreux médicaments, comme le \textbf{paracétamol}, synthétisé par acylation du para-aminophénol.
\item Les \textbf{anhydrides d'acide} donnent la même réaction, avec un acide carboxylique comme sous-produit au lieu de \ce{HCl}.
\end{itemize}
\end{aretenir}

\newpage

\coursec{Travaux pratiques}

\courssub{Objectifs du TP}

À l'issue de cette séance de travaux pratiques, tu seras capable de :

\begin{itemize}
\item mettre en évidence expérimentalement le \cle{caractère basique} d'une amine en solution aqueuse à l'aide d'un indicateur coloré et d'une mesure de pH ;
\item écrire et interpréter l'équation-bilan de la réaction de l'eau sur une amine (équilibre de Brønsted) ;
\item réaliser la \cle{neutralisation} d'une amine par un acide fort et écrire l'équation-bilan de formation d'un \cle{sel d'alkylammonium} ;
\item isoler par évaporation douce les cristaux de chlorure d'éthylammonium \ce{C2H5NH3+ Cl-} ;
\item respecter les règles de sécurité liées à la manipulation des amines et des acides.
\end{itemize}

Ce TP illustre concrètement la situation de la leçon : la plupart des médicaments azotés (antihistaminiques, anesthésiants locaux) sont commercialisés non pas sous forme d'amine libre, mais sous forme de \emph{sel d'ammonium}, plus stable, moins volatil et beaucoup plus soluble dans l'eau. Tu vas reproduire ici, à petite échelle, cette transformation essentielle.

\courssub{Matériel et produits}

\begin{center}
\rowcolors{2}{popblueL}{white}
\begin{tabularx}{\linewidth}{|l|X|}
\hline
\rowcolor{popblue!40} \textbf{Verrerie et matériel} & \textbf{Quantité} \\
\hline
Béchers de \SI{100}{mL} & 2 \\
\hline
Tubes à essais avec portoir & 3 \\
\hline
Pipettes compte-gouttes (Pasteur) & 3 \\
\hline
Verre de montre (ou coupelle propre) & 1 \\
\hline
Agitateur en verre & 1 \\
\hline
Papier pH (bandelettes) et échelle colorimétrique & 1 carnet \\
\hline
Pissette d'eau distillée & 1 \\
\hline
Plaque chauffante ou bec Bunsen (chauffage doux) & 1 \\
\hline
Gants, lunettes de protection, blouse & 1 paire par binôme \\
\hline
\end{tabularx}
\end{center}

\begin{center}
\rowcolors{2}{popgreenL}{white}
\begin{tabularx}{\linewidth}{|l|X|}
\hline
\rowcolor{popgreen!40} \textbf{Produits chimiques} & \textbf{Concentration / quantité} \\
\hline
Éthylamine \ce{C2H5NH2} en solution aqueuse & environ \SI{1}{mol.L^{-1}}, quelques gouttes \\
\hline
Solution d'acide chlorhydrique diluée & \SI{0,5}{mol.L^{-1}}, environ \SI{10}{mL} \\
\hline
Phénolphtaléine (solution alcoolique) & 2 gouttes \\
\hline
Eau distillée & environ \SI{50}{mL} \\
\hline
\end{tabularx}
\end{center}

\begin{attention}[Alternatives avec les ressources locales]
Si l'éthylamine pure n'est pas disponible dans ton laboratoire, plusieurs solutions existent :
\begin{itemize}
\item utiliser une solution commerciale d'éthylamine à \SI{70}{\%} dans l'eau (flacon courant dans les stocks de lycée) en la diluant au dixième ;
\item utiliser de la méthylamine ou de la butylamine si seules celles-ci sont en stock : le principe du TP reste identique ;
\item à défaut de phénolphtaléine, le papier pH seul suffit, ou encore quelques gouttes de jus de \emph{hibiscus} (bissap) fraîchement préparé : il vire du rouge au vert-jaune en milieu basique, c'est un excellent indicateur coloré naturel ;
\item à défaut de verre de montre, une coupelle en porcelaine ou même une assiette en verre propre convient pour l'évaporation.
\end{itemize}
\end{attention}

\courssub{Sécurité}

\begin{attention}[Consignes de sécurité impératives]
\textbf{Pictogrammes présents sur les flacons (décrits en mots) :}
\begin{itemize}
\item \emph{Flamme} (inflammable) : l'éthylamine pure est très volatile et inflammable. Aucune flamme nue à proximité du flacon ; le chauffage de l'évaporation ne concerne que la solution aqueuse finale, quasi débarrassée d'amine libre.
\item \emph{Silhouette corrodée} (corrosif) : l'acide chlorhydrique et les amines concentrées provoquent des brûlures. En cas de contact avec la peau, rincer abondamment à l'eau pendant plusieurs minutes et prévenir l'enseignant.
\item \emph{Point d'exclamation} (nocif, irritant) : les vapeurs d'amine irritent les voies respiratoires et les yeux.
\end{itemize}
\textbf{Gestes à adopter :}
\begin{itemize}
\item port obligatoire de la blouse, des gants et des lunettes pendant toute la séance ;
\item manipulation de l'amine sous hotte ou à proximité d'une fenêtre ouverte, en éloignant le visage des récipients ; ne jamais sentir directement le flacon ;
\item verser toujours l'acide dans la solution d'amine (et jamais l'inverse), goutte à goutte, en agitant ;
\item garder les flacons rebouchés dès que le prélèvement est terminé ;
\item ne rien jeter à l'évier sans l'accord de l'enseignant : les solutions neutralisées (sels d'ammonium dilués) peuvent être évacuées avec beaucoup d'eau.
\end{itemize}
\end{attention}

\courssub{Schéma du dispositif}

\begin{popfigure}
\begin{center}
\begin{tikzpicture}[scale=0.95, every node/.style={font=\small}]
% Bécher de la solution d'amine
\draw[thick, popblue] (0,0) -- (0,2.4) (2.4,0) -- (2.4,2.4);
\draw[thick, popblue] (0,0) -- (2.4,0);
\draw[thick, popblue] (2.4,2.4) -- (2.7,2.7);
% liquide rose
\fill[poppink!45] (0.06,0.06) rectangle (2.34,1.5);
\draw[poppink, thick] (0.06,1.5) -- (2.34,1.5);
\node[popdark] at (1.2,0.75) {solution rose};
\node[below, popdark] at (1.2,-0.15) {amine + eau + phénolphtaléine};
% pipette compte-gouttes avec HCl
\draw[thick, poporange] (1.2,4.6) -- (1.2,3.1);
\draw[thick, poporange] (1.05,3.6) -- (1.35,3.6);
\draw[thick, poporange] (1.05,3.6) -- (1.13,3.1) (1.35,3.6) -- (1.27,3.1);
\draw[thick, poporange] (1.13,3.1) -- (1.27,3.1);
\fill[poporange!60] (1.14,3.55) rectangle (1.26,3.15);
\fill[poporangeL] (1.1,4.6) circle (0.16);
\node[right, poporange!80!black] at (1.45,4.2) {\ce{HCl} dilué, goutte à goutte};
% goutte qui tombe
\fill[poporange!70] (1.2,2.6) circle (0.06);
\fill[poporange!70] (1.2,2.1) circle (0.06);
% flèche évolution
\draw[-{Stealth[length=3mm]}, very thick, popmuted] (3.4,1.2) -- (4.6,1.2);
% verre de montre avec cristaux
\draw[thick, popgreen!70!black] (5.4,0.55) .. controls (5.9,0.05) and (7.5,0.05) .. (8.0,0.55);
\draw[thick, popgreen!70!black] (5.4,0.55) -- (5.25,0.75) (8.0,0.55) -- (8.15,0.75);
\foreach \x/\y in {5.9/0.42,6.3/0.35,6.7/0.45,7.1/0.33,7.5/0.42,6.1/0.5,6.9/0.5}{
  \draw[popgreen!60!black, fill=popgreenL] (\x,\y) -- ++(0.09,0.09) -- ++(0.09,-0.09) -- ++(-0.09,-0.09) -- cycle;}
\node[below, popdark] at (6.7,-0.15) {cristaux de \ce{C2H5NH3+ Cl-} après évaporation};
% papier pH
\draw[thick, poppurple] (5.6,1.9) rectangle (7.8,2.5);
\fill[poppink!60] (5.7,2.0) rectangle (6.4,2.4);
\node[right, popdark] at (6.5,2.2) {papier pH : pH $\approx 11$};
\end{tikzpicture}
\end{center}
\legende{Dispositif du TP : mise en évidence de la basicité de l'éthylamine (solution rose avec la phénolphtaléine, pH $> 7$), neutralisation goutte à goutte par l'acide chlorhydrique, puis évaporation douce sur verre de montre pour obtenir les cristaux de chlorure d'éthylammonium.}
\end{popfigure}

\courssub{Protocole}

\begin{methode}[Déroulement de la séance]
\begin{enumerate}
\item \textbf{Préparation de la solution d'amine.} Introduire environ \SI{40}{mL} d'eau distillée dans un bécher de \SI{100}{mL}. Ajouter, sous hotte ou près de la fenêtre, \textbf{5 gouttes} de solution d'éthylamine à l'aide d'une pipette compte-gouttes. Homogénéiser avec l'agitateur en verre.
\item \textbf{Test à la phénolphtaléine.} Verser environ \SI{5}{mL} de cette solution dans un tube à essais. Ajouter 2 gouttes de phénolphtaléine. Observer et noter la couleur.
\item \textbf{Mesure du pH.} Déposer une goutte de la solution du bécher sur une bandelette de papier pH. Comparer avec l'échelle colorimétrique et noter la valeur du pH.
\item \textbf{Neutralisation par l'acide.} Ajouter au contenu du tube à essais la solution d'acide chlorhydrique \textbf{goutte à goutte}, en agitant après chaque goutte, jusqu'à \emph{décoloration complète} de la phénolphtaléine. Compter le nombre de gouttes d'acide nécessaires. Mesurer à nouveau le pH de la solution décolorée.
\item \textbf{Obtention du sel.} Verser le contenu du tube décoloré dans le bécher, puis transférer environ \SI{10}{mL} de la solution sur le verre de montre. Chauffer \emph{très doucement} (petite flamme ou plaque tiède) jusqu'à évaporation presque complète. Laisser refroidir.
\item \textbf{Observation des cristaux.} Observer à l'œil nu (ou à la loupe) le dépôt solide blanc formé sur le verre de montre : ce sont les cristaux de chlorure d'éthylammonium. Noter leur aspect.
\item \textbf{Nettoyage.} Rincer toute la verrerie à grande eau, ranger le poste de travail, se laver les mains.
\end{enumerate}
\end{methode}

\courssub{Observations et résultats attendus}

\begin{center}
\rowcolors{2}{popgoldL}{white}
\begin{tabularx}{\linewidth}{|X|X|X|}
\hline
\rowcolor{popgold!50} \textbf{Étape} & \textbf{Observation attendue} & \textbf{Mesure type} \\
\hline
Solution d'amine + phénolphtaléine & coloration rose franche & --- \\
\hline
Mesure du pH initial & teinte de papier pH correspondant à un milieu basique & pH $\approx 11$ \\
\hline
Ajout de \ce{HCl} goutte à goutte & décoloration progressive du rose & 8 à 12 gouttes \\
\hline
pH après décoloration & milieu quasi neutre à légèrement acide & pH $\approx 6$ \\
\hline
Évaporation sur verre de montre & dépôt de cristaux blancs, solubles dans l'eau & masse $\approx \SI{0,1}{g}$ \\
\hline
\end{tabularx}
\end{center}

\begin{attention}[Pièges d'observation]
\begin{itemize}
\item Si la solution ne rosit pas avec la phénolphtaléine, la solution d'amine est probablement trop diluée ou trop ancienne (les amines s'altèrent à l'air) : ajouter quelques gouttes d'amine.
\item La décoloration doit être obtenue \emph{goutte à goutte} : un excès brutal d'acide donne un pH très inférieur à 7 et fausse l'interprétation du point de neutralisation.
\item Un chauffage trop vif lors de l'évaporation fait projeter le sel et brunir les cristaux : toujours chauffer doucement.
\end{itemize}
\end{attention}

\courssub{Exploitation}

\begin{exoresolu}[Questions guidées d'exploitation]
\textbf{Question 1.} Pourquoi la phénolphtaléine vire-t-elle au rose dans la solution d'éthylamine ? Écrire l'équation-bilan de la réaction de l'eau sur l'éthylamine.

\textbf{Question 2.} Le pH mesuré vaut environ 11, et non 14 comme pour une solution de soude de même concentration formelle (environ \SI{1}{mol.L^{-1}}). Qu'en déduis-tu sur la force de la base éthylamine ?

\textbf{Question 3.} Écrire l'équation-bilan de la réaction qui se produit lors de l'ajout d'acide chlorhydrique. Comment nomme-t-on le produit formé ?

\textbf{Question 4.} Pourquoi le sel obtenu est-il solide et très soluble dans l'eau, alors que l'éthylamine est un liquide volatil à l'odeur désagréable ?

\textbf{Question 5.} Citer une application industrielle ou pharmaceutique de la transformation d'une amine en sel d'ammonium.
\tcblower
\textbf{Réponse 1.} La teinte rose de la phénolphtaléine indique un milieu \cle{basique} (pH $> 8{,}2$). L'éthylamine possède un doublet non liant sur l'atome d'azote : elle capte un proton de l'eau selon l'équilibre :
\[ \ce{C2H5NH2 + H2O <=> C2H5NH3+ + OH-} \]
Les ions hydroxyde \ce{OH-} formés sont responsables de la basicité de la solution.

\textbf{Réponse 2.} Une solution de soude (base forte) de concentration formelle voisine (\SI{1}{mol.L^{-1}}) aurait un pH proche de 14. Le pH de 11, plus faible, montre que la réaction avec l'eau est \emph{partielle} (équilibre limité) : l'éthylamine est une \cle{base faible}. C'est pourquoi on écrit la réaction avec une double flèche \ce{<=>}.

\textbf{Réponse 3.} L'acide chlorhydrique apporte des ions oxonium \ce{H3O+} qui neutralisent l'amine :
\[ \ce{C2H5NH2 + H3O+ + Cl- -> C2H5NH3+ + Cl- + H2O} \]
soit, en écriture simplifiée : \ce{C2H5NH2 + HCl -> C2H5NH3+ Cl-}. Le produit est le \cle{chlorure d'éthylammonium} (nom systématique : chlorure d'éthanaminium), un \emph{sel d'alkylammonium}. La décoloration de la phénolphtaléine marque la disparition des ions \ce{OH-} en excès.

\textbf{Réponse 4.} Le sel \ce{C2H5NH3+ Cl-} est un composé \emph{ionique} : les attractions électrostatiques entre ions \ce{C2H5NH3+} et \ce{Cl-} le rendent solide à température ambiante, non volatil (donc sans odeur) et très soluble dans l'eau, solvant polaire. L'amine moléculaire, elle, est volatile et peu soluble dès que la chaîne carbonée s'allonge.

\textbf{Réponse 5.} En pharmacie, de nombreux principes actifs azotés (chlorphéniramine, lidocaïne, quinine) sont formulés sous forme de chlorhydrates ou de citrates : le sel est stable à l'air, inodore, se conserve longtemps et se dissout dans le sang ou dans l'eau d'injection. C'est exactement la transformation réalisée dans ce TP.
\end{exoresolu}

\courssub{Conclusion}

Cette séance t'a permis de vérifier expérimentalement deux propriétés fondamentales des amines. D'une part, l'éthylamine se comporte comme une \cle{base de Brønsted} : grâce au doublet libre de l'atome d'azote, elle capte un proton de l'eau et rend la solution basique (virage de la phénolphtaléine au rose, pH $\approx 11$), selon l'équilibre \ce{C2H5NH2 + H2O <=> C2H5NH3+ + OH-}. D'autre part, l'action d'un acide fort sur l'amine conduit quantitativement à un \cle{sel d'alkylammonium}, ici le chlorure d'éthylammonium \ce{C2H5NH3+ Cl-}, isolé sous forme de cristaux blancs solubles dans l'eau. Cette réaction de neutralisation, simple et rapide, est à la base de la formulation de très nombreux médicaments : transformer une amine en sel, c'est la rendre stable, inodore et soluble, donc administrable. Tu disposes désormais des gestes expérimentaux et des équations-bilans indispensables pour aborder sereinement la suite de la leçon, notamment la réaction d'Hofmann et la synthèse des amides.

\newpage

\coursec{Méthodes et exercices résolus}

\courssub{Fiche méthode 1 : Écrire une formule semi-développée et identifier la classe d'une amine}

\begin{methode}[Reconnaître la classe d'une amine]
\begin{enumerate}
\item \textbf{Repérer l'atome d'azote} : toute amine dérive de l'ammoniac \ce{NH3} par remplacement d'un ou plusieurs atomes d'hydrogène par des groupes alkyle (ou aryle) R.
\item \textbf{Compter les groupes carbonés} liés directement à l'azote :
\begin{itemize}
\item 1 groupe R (et 2 H) : amine \cle{primaire} $R-\ce{NH2}$ ;
\item 2 groupes R (et 1 H) : amine \cle{secondaire} $R-\ce{NH}-R'$ ;
\item 3 groupes R (et 0 H) : amine \cle{tertiaire} $N(R)(R')(R'')$, notée $NR_3$.
\end{itemize}
\item \textbf{Attention au piège} : la classe dépend du nombre de groupes portés par \emph{l'azote}, pas de la nature des chaînes carbonées. \ce{(CH3)3C-NH2} est une amine \emph{primaire} même si le carbone porteur est tertiaire !
\item \textbf{Vérifier la valence} : l'azote forme toujours 3 liaisons et conserve un doublet non liant ; le compte des hydrogènes doit boucler ($N$ porte $3 - n$ hydrogènes pour une amine de classe $n$).
\end{enumerate}
\end{methode}

\begin{exoresolu}[Classes d'amines — type Bac]
On donne les composés suivants :
\textbf{(a)} \ce{CH3-CH2-NH2} ; \textbf{(b)} \ce{CH3-NH-CH(CH3)2} ; \textbf{(c)} \ce{N(CH3)3} ; \textbf{(d)} \ce{(CH3)3C-NH2} ; \textbf{(e)} \ce{CH3-CH2-CH2-NH-CH3}.
Pour chacun, écrire la formule semi-développée complète et préciser la classe de l'amine.
\tcblower
\textbf{Raisonnement} : on compte, pour chaque molécule, le nombre de groupes alkyle directement liés à l'atome d'azote.

\begin{center}
\begin{tabularx}{\linewidth}{|l|X|l|}
\hline
\rowcolor{popblueL} \textbf{Composé} & \textbf{Groupes sur N} & \textbf{Classe} \\
\hline
(a) \chemfig{CH_3-CH_2-NH_2} & 1 groupe (éthyle), 2 H & primaire \\
\hline
(b) \chemfig{CH_3-NH-CH(-[2]CH_3)(-[6]CH_3)} & 2 groupes (méthyle, isopropyle), 1 H & secondaire \\
\hline
(c) \chemfig{N(-[2]CH_3)(-[6]CH_3)-CH_3} & 3 groupes méthyle, 0 H & tertiaire \\
\hline
(d) \chemfig{(CH_3)_3C-NH_2} & 1 groupe (tert-butyle), 2 H & primaire \\
\hline
(e) \chemfig{CH_3-CH_2-CH_2-NH-CH_3} & 2 groupes (propyle, méthyle), 1 H & secondaire \\
\hline
\end{tabularx}
\end{center}

\textbf{Point clé pour (d)} : le groupe tert-butyle est ramifié, mais l'azote ne porte qu'\emph{un seul} groupe carboné : l'amine est primaire. Le piège classique consiste à confondre la classe de l'amine avec la classe du carbone.

\textbf{Conclusion} : (a) et (d) sont des amines primaires, (b) et (e) des amines secondaires, (c) une amine tertiaire.
\end{exoresolu}

\courssub{Fiche méthode 2 : Nommer une amine}

\begin{methode}[Nomenclature des amines]
\begin{enumerate}
\item \textbf{Amines primaires} \ce{R-NH2} : nommer la chaîne carbonée la plus longue contenant le carbone lié à N, remplacer le « e » final de l'alcane par le suffixe \cle{-amine}, et indiquer l'indice de position du groupe \ce{-NH2} : \emph{alcan-n-amine}. Exemple : \ce{CH3-CH(NH2)-CH3} = propan-2-amine.
\item \textbf{Amines secondaires et tertiaires} : choisir le groupe alkyle le plus long comme chaîne principale (suffixe \cle{-amine}) ; les autres groupes sont cités en préfixe précédés de la lettre \cle{N} (ou \cle{N,N} s'il y en a deux identiques) pour montrer qu'ils sont portés par l'azote. Exemple : \ce{CH3-NH-CH2-CH3} = N-méthyléthanamine ; \ce{N(CH2CH3)3} = N,N-diéthyléthanamine.
\item \textbf{Numéroter} la chaîne principale de façon à donner à l'azote l'indice le plus petit.
\item \textbf{Vérifier} : la somme des groupes cités (chaîne principale + préfixes N) doit correspondre à tous les groupes de la formule.
\end{enumerate}
\end{methode}

\begin{exoresolu}[Nommer des amines — type Bac]
Nommer les amines suivantes :
\textbf{(a)} \ce{CH3-CH2-CH2-NH2} ; \textbf{(b)} \ce{CH3-CH(NH2)-CH2-CH3} ; \textbf{(c)} \ce{CH3-CH2-NH-CH3} ; \textbf{(d)} \ce{CH3-N(CH3)-CH2-CH2-CH3}.
\tcblower
\textbf{(a)} \ce{CH3-CH2-CH2-NH2} : chaîne de 3 carbones, amine primaire, \ce{-NH2} en position 1 $\Rightarrow$ \textbf{propan-1-amine}.

\textbf{(b)} \ce{CH3-CH(NH2)-CH2-CH3} : chaîne de 4 carbones (butane), \ce{-NH2} sur le carbone 2 (numérotation qui donne l'indice le plus petit) $\Rightarrow$ \textbf{butan-2-amine}.

\textbf{(c)} \ce{CH3-CH2-NH-CH3} : amine secondaire ; le groupe le plus long est l'éthyle (chaîne principale : éthanamine), le méthyle est porté par l'azote $\Rightarrow$ \textbf{N-méthyléthanamine}.

\textbf{(d)} \ce{CH3-N(CH3)-CH2-CH2-CH3} : amine tertiaire ; groupe le plus long = propyle (propan-1-amine), deux méthyles sur l'azote $\Rightarrow$ \textbf{N,N-diméthylpropan-1-amine}.

\textbf{Conclusion} : (a) propan-1-amine ; (b) butan-2-amine ; (c) N-méthyléthanamine ; (d) N,N-diméthylpropan-1-amine.
\end{exoresolu}

\courssub{Fiche méthode 3 : Expliquer les propriétés physiques par les liaisons hydrogène}

\begin{methode}[Liaisons hydrogène et propriétés physiques]
\begin{enumerate}
\item \textbf{Repérer les liaisons N--H} dans la molécule : une amine primaire en possède 2, une secondaire 1, une tertiaire \emph{aucune}.
\item \textbf{Conclure sur les liaisons hydrogène intermoléculaires} : elles n'existent qu'entre molécules portant N--H (donneur) et un doublet de N (accepteur). Les amines tertiaires ne peuvent pas en former entre elles.
\item \textbf{Comparer les températures d'ébullition} : amine primaire $>$ amine secondaire $>$ amine tertiaire (à masse molaire comparable), car il faut fournir de l'énergie pour rompre les liaisons hydrogène. Mais une amine bout \emph{moins haut} que l'alcool de masse voisine, car la liaison N--H$\cdots$N est plus faible que O--H$\cdots$O (N moins électronégatif que O).
\item \textbf{Solubilité dans l'eau} : les amines à courte chaîne (moins de 5--6 C) sont solubles car elles forment des liaisons hydrogène avec l'eau (le doublet de N accepte, N--H donne) ; la solubilité \emph{diminue} quand la chaîne carbonée (partie hydrophobe) s'allonge.
\item \textbf{Rédiger la justification} en citant explicitement « liaisons hydrogène intermoléculaires » : c'est le mot-clé attendu par le correcteur.
\end{enumerate}
\end{methode}

\begin{exoresolu}[Températures d'ébullition — type Bac]
On donne les températures d'ébullition de trois composés de masse molaire voisine ($M \approx \SI{59}{g.mol^{-1}}$) : propan-1-ol : \SI{97}{\celsius} ; éthylméthylamine \ce{CH3CH2NHCH3} : \SI{36}{\celsius} ; triméthylamine \ce{N(CH3)3} : \SI{3}{\celsius}.
\begin{enumerate}
\item Justifier l'écart entre l'éthylméthylamine et la triméthylamine.
\item Pourquoi le propan-1-ol bout-il plus haut que l'éthylméthylamine ?
\item L'éthylméthylamine est-elle soluble dans l'eau ? Justifier.
\end{enumerate}
\tcblower
\textbf{1.} L'éthylméthylamine est une amine \emph{secondaire} : elle possède une liaison N--H. Ses molécules s'associent par \cle{liaisons hydrogène} intermoléculaires N--H$\cdots$N, qu'il faut rompre pour vaporiser le liquide. La triméthylamine est \emph{tertiaire} : sans liaison N--H, ses molécules ne sont retenues que par de faibles interactions de van der Waals. D'où $T_{eb}(\text{éthylméthylamine}) = \SI{36}{\celsius} > T_{eb}(\text{triméthylamine}) = \SI{3}{\celsius}$.

\textbf{2.} Le propan-1-ol forme des liaisons hydrogène O--H$\cdots$O. L'oxygène étant plus électronégatif que l'azote, la liaison O--H est plus polarisée que N--H : les liaisons hydrogène entre alcools sont \emph{plus fortes} qu'entre amines. D'où $T_{eb}(\text{propan-1-ol}) = \SI{97}{\celsius} > \SI{36}{\celsius}$.

\textbf{3.} Oui : l'éthylméthylamine, à chaîne carbonée courte, est soluble dans l'eau. Elle établit des liaisons hydrogène avec les molécules d'eau : sa liaison N--H est donneuse, le doublet non liant de son azote est accepteur. Cette hydratation compense la rupture des liaisons hydrogène eau--eau.

\textbf{Conclusion} : l'existence ou non de liaisons hydrogène intermoléculaires explique la hiérarchie des températures d'ébullition, et les liaisons hydrogène avec l'eau expliquent la solubilité des amines légères.
\end{exoresolu}

\courssub{Fiche méthode 4 : Caractère basique — réaction d'une amine avec l'eau}

\begin{methode}[Écrire et exploiter la réaction d'une amine avec l'eau]
\begin{enumerate}
\item \textbf{Identifier le couple acide/base} : l'amine \ce{R-NH2} est une base (doublet non liant sur N) ; l'eau joue le rôle d'acide. Couples : \ce{RNH3+}/\ce{RNH2} et \ce{H2O}/\ce{HO-}.
\item \textbf{Écrire l'équilibre} avec une double flèche (réaction \emph{limitée}) :
\[ \ce{R-NH2 + H2O <=> R-NH3+ + HO-} \]
L'ion \ce{R-NH3+} est l'\cle{ion alkylammonium}.
\item \textbf{Constante de basicité} : $K_B = \dfrac{[\ce{RNH3+}][\ce{HO-}]}{[\ce{RNH2}]}$ ; on utilise aussi $pK_A$ du couple \ce{RNH3+}/\ce{RNH2} : plus le $pK_A$ est grand, plus l'amine est basique.
\item \textbf{Conséquences} : la solution est basique ($pH > 7$), $[\ce{RNH3+}] = [\ce{HO-}]$ si l'eau est la seule autre espèce réactive.
\item \textbf{Comparer les basicités} : l'effet inductif donneur des groupes alkyle enrichit le doublet de N : en général amine secondaire $>$ primaire $>$ \ce{NH3} (la tertiaire est un peu moins basique que la secondaire en solution aqueuse, pour des raisons de solvatation).
\end{enumerate}
\end{methode}

\begin{exoresolu}[Solution aqueuse d'éthylamine — type Bac]
On dissout $n = \SI{2,0e-2}{mol}$ d'éthylamine \ce{C2H5NH2} dans de l'eau distillée pour obtenir $V = \SI{500}{mL}$ de solution. Le $pH$ mesuré vaut $11,9$ à \SI{25}{\celsius}.
\begin{enumerate}
\item Écrire l'équation-bilan de la réaction de l'éthylamine avec l'eau et préciser les couples mis en jeu.
\item Calculer les concentrations des espèces présentes et le taux d'avancement final $\tau$.
\item En déduire le $pK_A$ du couple \ce{C2H5NH3+}/\ce{C2H5NH2}.
\end{enumerate}
On donne : $K_e = \SI{1,0e-14}{}$ à \SI{25}{\celsius}.
\tcblower
\textbf{1. Équation et couples.}
\[ \ce{C2H5NH2 + H2O <=> C2H5NH3+ + HO-} \]
Couples : \ce{C2H5NH3+}/\ce{C2H5NH2} (acide/base) et \ce{H2O}/\ce{HO-}. La double flèche traduit une réaction \emph{limitée} : l'éthylamine est une base faible.

\textbf{2. Concentrations et taux d'avancement.}

Concentration apportée : $C = \dfrac{n}{V} = \dfrac{\SI{2,0e-2}{mol}}{\SI{0,500}{L}} = \SI{4,0e-2}{mol.L^{-1}}$.

$[\ce{H3O+}] = 10^{-pH} = 10^{-11,9} = \SI{1,26e-12}{mol.L^{-1}}$, d'où :
\[ [\ce{HO-}] = \frac{K_e}{[\ce{H3O+}]} = \frac{\SI{1,0e-14}{}}{\SI{1,26e-12}{}} = \SI{7,9e-3}{mol.L^{-1}} \]

Électroneutralité : $[\ce{C2H5NH3+}] = [\ce{HO-}] - [\ce{H3O+}] \approx [\ce{HO-}] = \SI{7,9e-3}{mol.L^{-1}}$.

Conservation : $[\ce{C2H5NH2}] = C - [\ce{C2H5NH3+}] = 4,0\times10^{-2} - 7,9\times10^{-3} = \SI{3,2e-2}{mol.L^{-1}}$.

Taux d'avancement : $\tau = \dfrac{[\ce{C2H5NH3+}]}{C} = \dfrac{7,9\times10^{-3}}{4,0\times10^{-2}} \approx 0,20$, soit environ \SI{20}{\%} : la réaction est bien limitée.

\textbf{3. Calcul du $pK_A$.}
\[ K_A = \frac{[\ce{C2H5NH2}][\ce{H3O+}]}{[\ce{C2H5NH3+}]} = \frac{3,2\times10^{-2}\times 1,26\times10^{-12}}{7,9\times10^{-3}} = \SI{5,1e-12}{} \]
\[ pK_A = -\log K_A = 11,3 \]

\textbf{Conclusion} : l'éthylamine est une base faible ($\tau \approx \SI{20}{\%}$) ; le couple \ce{C2H5NH3+}/\ce{C2H5NH2} a un $pK_A \approx 11,3$, valeur typique des ions alkylammonium.
\end{exoresolu}

\courssub{Fiche méthode 5 : Caractère nucléophile et réaction d'Hofmann}

\begin{methode}[Réaction d'Hofmann : amine + dérivé halogéné]
\begin{enumerate}
\item \textbf{Reconnaître le mécanisme} : le doublet non liant de l'azote (nucléophile) attaque le carbone porteur de l'halogène (électrophile, car C est $\delta+$ dans C--X). C'est une \cle{substitution nucléophile}.
\item \textbf{Écrire l'équation} : l'halogène part avec l'hydrogène de l'azote sous forme d'halogénure d'hydrogène $\mathrm{HX}$ :
\[ \mathrm{R{-}NH_2} + \mathrm{R'{-}X} \rightarrow \mathrm{R{-}NH{-}R'} + \mathrm{HX} \]
\item \textbf{Anticiper les alkylations successives} : l'amine obtenue est encore nucléophile ! La réaction ne s'arrête pas :
\begin{itemize}
\item amine primaire $\to$ amine secondaire $\to$ amine tertiaire $\to$ \cle{sel d'ammonium quaternaire} $[\mathrm{NR_4}]^+\mathrm{X}^-$.
\end{itemize}
On obtient donc un \emph{mélange} ; pour favoriser l'amine primaire, on utilise un excès d'ammoniac (ou d'amine de départ).
\item \textbf{Cas de l'ammoniac} : $\ce{NH3} + \mathrm{R{-}X} \rightarrow \mathrm{R{-}NH_2} + \mathrm{HX}$ : c'est la première étape de la synthèse d'Hofmann des amines.
\item \textbf{Électrophile / nucléophile} : l'amine est nucléophile (donneur de doublet) ; le dérivé halogéné est électrophile (accepteur). Ne jamais inverser.
\end{enumerate}
\end{methode}

\begin{exoresolu}[Synthèse d'Hofmann — type Bac]
Pour synthétiser un intermédiaire pharmaceutique, on fait réagir l'ammoniac \ce{NH3} en excès sur le bromoéthane \ce{C2H5Br}.
\begin{enumerate}
\item Écrire l'équation-bilan de la réaction et nommer le produit organique obtenu.
\item Justifier le caractère nucléophile de l'ammoniac et le caractère électrophile du bromoéthane.
\item Pourquoi utilise-t-on l'ammoniac en excès ? Écrire l'équation de la réaction parasite qui aurait lieu sinon.
\end{enumerate}
\tcblower
\textbf{1. Équation-bilan.}
\[ \ce{NH3 + C2H5Br -> C2H5NH2 + HBr} \]
Le produit organique est l'\textbf{éthanamine} (éthylamine), amine primaire.

\textbf{2. Justification électronique.} L'atome d'azote de \ce{NH3} porte un \cle{doublet non liant} qu'il peut offrir : l'ammoniac est \emph{nucléophile} (« ami des noyaux », donneur de doublet). Dans \ce{C2H5Br}, le brome, plus électronégatif que le carbone, polarise la liaison C--Br : le carbone porte une charge partielle $\delta+$ et constitue un site \emph{électrophile}. Le doublet de N attaque ce carbone, le brome part en \ce{Br-} puis arrache un proton : substitution nucléophile.

\textbf{3. Rôle de l'excès d'ammoniac.} L'éthanamine formée possède encore un doublet sur l'azote : elle est elle-même nucléophile et peut réagir sur une seconde molécule de bromoéthane :
\[ \ce{C2H5NH2 + C2H5Br -> (C2H5)2NH + HBr} \]
(réaction parasite donnant la diéthylamine, amine secondaire ; la suite mènerait à la triéthylamine puis au sel \ce{[(C2H5)4N]+Br-}). En mettant \ce{NH3} en large excès, le bromoéthane rencontre statistiquement bien plus souvent une molécule d'ammoniac qu'une molécule d'éthanamine : on limite les alkylations successives et on améliore le rendement en amine primaire.

\textbf{Conclusion} : \ce{NH3 + C2H5Br -> C2H5NH2 + HBr} ; l'excès d'ammoniac oriente la synthèse vers l'amine primaire en défavorisant la suralkylation.
\end{exoresolu}

\courssub{Fiche méthode 6 : Action d'une amine sur un chlorure d'acyle — synthèse des amides}

\begin{methode}[Préparer un amide à partir d'un chlorure d'acyle]
\begin{enumerate}
\item \textbf{Identifier les réactifs} : chlorure d'acyle $R-COCl$ (groupe \chemfig{C(=[1]O)-[7]Cl}) + amine primaire $R'-NH_2$ ou secondaire $R'-NH-R''$.
\item \textbf{Écrire l'équation} : le chlore et un hydrogène de l'azote s'éliminent sous forme de \ce{HCl} ; la liaison C--N se forme :
\[ R-COCl + R'-NH_2 \rightarrow R-CO-NH-R' + \ce{HCl} \]
Avec une amine secondaire : $R-COCl + R'NHR'' \rightarrow R-CO-NR'R'' + \ce{HCl}$.
\item \textbf{Reconnaître le groupe amide} : \chemfig{R-C(=[1]O)-[7]NH-R'}, liaison \cle{peptidique} si elle relie des acides aminés.
\item \textbf{Préciser les conditions} : réaction \emph{rapide et totale} (contrairement à l'action sur un acide carboxylique, lente et limitée) ; on ajoute souvent une base pour neutraliser le \ce{HCl} formé.
\item \textbf{Amines tertiaires} : sans H sur l'azote, elles ne donnent \emph{pas} d'amide stable par cette voie : réflexe de vérification.
\end{enumerate}
\end{methode}

\begin{exoresolu}[Synthèse d'un amide — type Bac]
Le N-éthyléthanamide est un intermédiaire de synthèse. On le prépare par action du chlorure d'éthanoyle \ce{CH3COCl} sur l'éthylamine \ce{C2H5NH2}.
\begin{enumerate}
\item Écrire l'équation-bilan de la réaction et entourer le groupe fonctionnel formé.
\item Citer deux avantages de cette voie par rapport à l'action de l'acide éthanoïque sur l'éthylamine.
\item On fait réagir $m = \SI{15,7}{g}$ de chlorure d'éthanoyle avec un excès d'éthylamine. Calculer la masse d'amide obtenue, la réaction étant totale.
\end{enumerate}
On donne : $M(\ce{CH3COCl}) = \SI{78,5}{g.mol^{-1}}$ ; $M(\ce{CH3CONHC2H5}) = \SI{87,1}{g.mol^{-1}}$.
\tcblower
\textbf{1. Équation-bilan.}
\[ \ce{CH3COCl + C2H5NH2 -> CH3-CO-NH-C2H5 + HCl} \]
Le produit est le \textbf{N-éthyléthanamide}, qui contient le groupe \cle{amide} \chemfig{CH_3-C(=[1]O)-[7]NH-CH_2CH_3} (la fonction est le groupe \ce{-CO-NH-}).

\textbf{2. Avantages.} La réaction avec le chlorure d'acyle est \emph{rapide} et \emph{totale}, alors que l'action directe de l'acide carboxylique sur l'amine est lente et limitée par un équilibre (elle passe par un sel d'ammonium). Le chlorure d'acyle est donc un réactif bien plus efficace, malgré sa corrosivité.

\textbf{3. Calcul de masse.} La réaction étant totale et l'éthylamine en excès, le chlorure d'éthanoyle est le réactif limitant. D'après l'équation, $n(\text{amide}) = n(\ce{CH3COCl})$.

\begin{center}
\begin{tabularx}{\linewidth}{|l|X|X|X|}
\hline
\rowcolor{popblueL} & \ce{CH3COCl} & \ce{CH3CONHC2H5} & \ce{HCl} \\
\hline
État initial (mol) & $n_0 = 0,200$ & 0 & 0 \\
\hline
État final (mol) & $0,200 - x_{max} = 0$ & $x_{max} = 0,200$ & $0,200$ \\
\hline
\end{tabularx}
\end{center}

\[ n_0 = \frac{m}{M} = \frac{\SI{15,7}{g}}{\SI{78,5}{g.mol^{-1}}} = \SI{0,200}{mol} \]
\[ m(\text{amide}) = n_0 \times M(\text{amide}) = 0,200 \times 87,1 = \SI{17,4}{g} \]

\textbf{Conclusion} : la réaction totale de \SI{15,7}{g} de chlorure d'éthanoyle avec un excès d'éthylamine fournit \SI{17,4}{g} de N-éthyléthanamide.
\end{exoresolu}

\begin{attention}[Les 5 erreurs qui coûtent des points au Bac]
\begin{enumerate}
\item \textbf{Confondre la classe de l'amine avec la classe du carbone} : \ce{(CH3)3C-NH2} est une amine \emph{primaire}. On compte les groupes carbonés portés par l'\emph{azote}, jamais la ramification du reste de la molécule.
\item \textbf{Oublier le « N » dans la nomenclature} : \ce{CH3-NH-C2H5} se nomme N-méthyléthanamine, pas « méthyléthanamine ». Sans le N, le nom est faux et le point est perdu.
\item \textbf{Écrire la réaction amine + eau avec une flèche simple} : c'est un équilibre (base faible) : $RNH_2 + H_2O \rightleftharpoons RNH_3^+ + HO^-$. La double flèche et les ions $RNH_3^+$ et $HO^-$ sont exigés.
\item \textbf{Arrêter la réaction d'Hofmann à la première étape} : le produit d'une alkylation est encore nucléophile ; le jury attend que tu mentionnes les alkylations successives (secondaire, tertiaire, sel d'ammonium quaternaire) et le rôle de l'excès d'ammoniac.
\item \textbf{Former un amide avec une amine tertiaire} ou oublier le \ce{HCl} : la réaction $RCOCl + R'_2NH$ ne donne pas d'amide stable (pas de H sur N). Et dans toute équation avec chlorure d'acyle, le \ce{HCl} doit apparaître pour équilibrer : sans lui, l'équation est fausse.
\end{enumerate}
\end{attention}

\newpage

\coursec{Exercices}

\courssub{Je vérifie mes connaissances}

\begin{exercice}[QCM : les classes d'amines]
Pour chaque question, une seule réponse est exacte. Coche la bonne lettre.
\begin{enumerate}
\item Une amine secondaire possède :
\begin{itemize}
\item[\textbf{a)}] deux atomes d'hydrogène liés à l'azote ;
\item[\textbf{b)}] deux groupes alkyles (ou aryles) liés à l'azote ;
\item[\textbf{c)}] trois groupes alkyles (ou aryles) liés à l'azote.
\end{itemize}
\item La formule générale d'une amine aliphatique saturée et acyclique est :
\begin{itemize}
\item[\textbf{a)}] $C_nH_{2n+1}N$ ;
\item[\textbf{b)}] $C_nH_{2n+3}N$ ;
\item[\textbf{c)}] $C_nH_{2n-1}N$.
\end{itemize}
\item Quel est le nom \textbf{systématique (IUPAC)} de l'amine de formule semi-développée \chemfig{CH_3-NH-CH_2-CH_3} ?
\begin{itemize}
\item[\textbf{a)}] N-méthyléthanamine ;
\item[\textbf{b)}] éthylméthylamine (nom usuel) ;
\item[\textbf{c)}] N-éthylméthanamine.
\end{itemize}
\item Le caractère basique des amines est dû :
\begin{itemize}
\item[\textbf{a)}] à la présence d'un doublet non liant sur l'atome d'azote ;
\item[\textbf{b)}] à la présence d'atomes d'hydrogène sur l'azote ;
\item[\textbf{c)}] à la masse molaire élevée de la molécule.
\end{itemize}
\end{enumerate}
\end{exercice}

\begin{exercice}[Vrai ou faux ? Justifie]
Réponds par \emph{vrai} ou \emph{faux}, puis justifie ta réponse en une ou deux phrases.
\begin{enumerate}
\item Toutes les amines peuvent former des liaisons hydrogène \textbf{entre leurs propres molécules}.
\item Une amine tertiaire ne peut pas réagir avec l'eau car elle ne possède pas d'atome d'hydrogène lié à l'azote.
\item La réaction d'une amine primaire sur un dérivé halogéné permet d'obtenir une amine secondaire.
\item L'action d'une amine sur un chlorure d'acyle conduit à un amide.
\item Les amines de petite masse molaire (jusqu'à 5-6 atomes de carbone) sont très solubles dans l'eau.
\end{enumerate}
\end{exercice}

\begin{exercice}[Texte à trous]
Complète le texte suivant avec les mots ou expressions de la liste ci-dessous (certains peuvent être utilisés plusieurs fois, d'autres pas du tout) :
\medskip
\begin{center}
\emph{doublet non liant ; nucléophile ; électrophile ; pyramide ; plane ; basique ; acide ; amide ; liaisons hydrogène ; liaisons covalentes}
\end{center}
\medskip
La molécule d'ammoniac \ce{NH3} a une géométrie en \ldots\ldots\ldots\ldots\ldots\ldots à base triangulaire ; l'atome d'azote porte un \ldots\ldots\ldots\ldots\ldots\ldots qui confère aux amines un caractère \ldots\ldots\ldots\ldots\ldots\ldots et un caractère \ldots\ldots\ldots\ldots\ldots\ldots. Les amines primaires et secondaires établissent entre leurs molécules des \ldots\ldots\ldots\ldots\ldots\ldots, ce qui explique leurs températures d'ébullition relativement élevées. L'action d'une amine sur un chlorure d'acyle conduit à un \ldots\ldots\ldots\ldots\ldots\ldots.
\end{exercice}

\begin{exercice}[Définitions essentielles]
\begin{enumerate}
\item Définis le terme \cle{amine} et donne sa formule générale pour une amine aliphatique saturée.
\item Définis les trois classes d'amines en précisant le nombre de groupes alkyles (ou aryles) liés à l'atome d'azote.
\item Énonce ce qu'on appelle \cle{réaction d'Hofmann}.
\item Qu'appelle-t-on \cle{caractère nucléophile} d'une amine ? À quoi est-il dû ?
\end{enumerate}
\end{exercice}

\courssub{Je m'entraîne}

\begin{exercice}[Classes et nomenclature]
On considère les cinq amines suivantes, données par leurs formules semi-développées :
\begin{center}
A : \chemfig{CH_3-CH_2-CH_2-NH_2} \qquad
B : \chemfig{CH_3-CH(-[2]NH_2)-CH_3} \qquad
C : \chemfig{CH_3-NH-CH_2-CH_3} \\[6pt]
D : \chemfig{CH_3-N(-[2]CH_3)-CH_2-CH_3} \qquad
E : \chemfig{CH_3-CH_2-NH-CH_2-CH_3}
\end{center}
\begin{enumerate}
\item Précise la classe (primaire, secondaire ou tertiaire) de chacune de ces amines.
\item Nomme chaque amine selon la nomenclature officielle (IUPAC).
\item Parmi ces amines, indique celles qui sont isomères les unes des autres et précise leur formule brute commune.
\end{enumerate}
\end{exercice}

\begin{exercice}[Isomères de formule \ce{C3H9N}]
\begin{enumerate}
\item Écris les formules semi-développées de \textbf{toutes} les amines de formule brute \ce{C3H9N}.
\item Nomme chacune d'elles et indique sa classe.
\item Explique, en t'appuyant sur les liaisons hydrogène, laquelle de ces amines possède la température d'ébullition la plus élevée et laquelle possède la plus faible.
\end{enumerate}
\end{exercice}

\begin{exercice}[Une amine dans l'eau]
On dissout de la propan-1-amine \ce{C3H7NH2} dans l'eau distillée. Le pH de la solution obtenue, mesuré à \SI{25}{\celsius}, vaut $11,8$.
\begin{enumerate}
\item Écris l'équation-bilan de la réaction de la propan-1-amine avec l'eau.
\item Identifie les deux couples acide/base mis en jeu.
\item Justifie, à partir de l'équation-bilan, que le pH de la solution est supérieur à 7.
\item Le p$K_A$ du couple \ce{C3H7NH3+}/\ce{C3H7NH2} vaut $10,6$. 
\begin{enumerate}
\item Trace l'axe des p$K_A$ et indique les domaines de prédominance des deux espèces.
\item Quelle espèce prédomine dans la solution de pH $11,8$ ?
\item Calcule le rapport $\dfrac{[\ce{C3H7NH2}]}{[\ce{C3H7NH3+}]}$ dans cette solution.
\end{enumerate}
\end{enumerate}
\end{exercice}

\begin{popfigure}
\begin{tikzpicture}[scale=0.9, >=Stealth]
% Axe pKa pour propan-1-amine
\draw[->] (0,0) -- (10,0) node[right] {pH};
\foreach \x in {0,2,4,6,7,8,10} \draw (\x,0.1) -- (\x,-0.1) node[below] {\x};
\draw[thick, popblue] (0,0.5) -- (10,0.5) node[right] {Axe p$K_A$};
\draw[popblue, thick] (10.6,0.3) -- (10.6,0.7) node[above] {p$K_A$ = 10,6};
\draw[poporange, thick, <->] (0,0.8) -- (10.6,0.8) node[midway, above] {\ce{C3H7NH3+} prédomine};
\draw[popgreen, thick, <->] (10.6,1.1) -- (10,1.1) node[midway, above] {\ce{C3H7NH2} prédomine};
\draw[popred, dashed, thick] (11.8, -0.3) -- (11.8, 1.3) node[above] {pH = 11,8};
\end{tikzpicture}
\legende{Axe des p$K_A$ pour le couple \ce{C3H7NH3+}/\ce{C3H7NH2} (p$K_A$ = 10,6).}
\end{popfigure}

\begin{exercice}[Solubilité et liaisons hydrogène]
On considère la butan-1-amine \ce{C4H11N} ($M = \SI{73}{g.mol^{-1}}$) et la N,N-diméthyléthanamine \ce{C4H11N} ($M = \SI{73}{g.mol^{-1}}$).
\begin{enumerate}
\item Écris les formules semi-développées de ces deux amines et précise leur classe.
\item À masse molaire égale, laquelle de ces deux amines est la plus soluble dans l'eau ? Laquelle a la température d'ébullition la plus élevée ?
\item Justifie tes réponses en expliquant le rôle des liaisons hydrogène, d'une part entre molécules d'amine (état pur), d'autre part entre l'amine et l'eau (solubilité).
\end{enumerate}
\end{exercice}

\courssub{Je me prépare au Bac}

\begin{exercice}[Type Bac — La réaction d'Hofmann au laboratoire]
Au laboratoire de chimie d'un lycée de Douala, un professeur fait réagir de l'ammoniac \ce{NH3} en excès avec du 1-bromopropane \ce{CH3-CH2-CH2-Br} afin de préparer une amine. Il explique à ses élèves que la réaction ne s'arrête pas forcément au premier stade.

\medskip
\textbf{Données :} $M(\ce{C}) = \SI{12}{g.mol^{-1}}$ ; $M(\ce{H}) = \SI{1}{g.mol^{-1}}$ ; $M(\ce{N}) = \SI{14}{g.mol^{-1}}$ ; $M(\ce{Br}) = \SI{80}{g.mol^{-1}}$.
\begin{enumerate}
\item Rappelle la définition d'une amine et donne la formule générale d'une amine aliphatique saturée.
\item Justifie le caractère nucléophile de l'ammoniac et des amines.
\item Écris l'équation-bilan de la première réaction entre l'ammoniac et le 1-bromopropane. Nomme le produit organique obtenu et précise sa classe.
\item Le produit obtenu peut réagir à son tour avec le 1-bromopropane. Écris les équations-bilans des deux réactions d'Hofmann successives suivantes, en nommant chaque produit organique et en précisant sa classe.
\item Pourquoi le professeur utilise-t-il un excès d'ammoniac s'il souhaite obtenir majoritairement l'amine primaire ?
\item \textbf{Application numérique :} On fait réagir $n = \SI{0,10}{mol}$ d'ammoniac avec $m = \SI{12,3}{g}$ de 1-bromopropane. 
\begin{enumerate}
\item Détermine le réactif limitant pour la formation de l'amine primaire.
\item Calcule la masse maximale d'amine primaire (propan-1-amine) que l'on peut espérer obtenir (rendement théorique 100\,\%).
\end{enumerate}
\end{enumerate}
\end{exercice}

\begin{exercice}[Type Bac — Synthèse du paracétamol]
Le paracétamol est un antalgique très utilisé au Cameroun pour traiter les douleurs et la fièvre. L'une des étapes de sa synthèse industrielle consiste à faire réagir le para-aminophénol \ce{HO-C6H4-NH2} avec le chlorure d'éthanoyle \ce{CH3-COCl}. On obtient le paracétamol \ce{HO-C6H4-NH-CO-CH3} et du chlorure d'hydrogène.

\medskip
\textbf{Données :} $M(\ce{C}) = \SI{12}{g.mol^{-1}}$ ; $M(\ce{H}) = \SI{1}{g.mol^{-1}}$ ; $M(\ce{N}) = \SI{14}{g.mol^{-1}}$ ; $M(\ce{O}) = \SI{16}{g.mol^{-1}}$ ; $M(\ce{Cl}) = \SI{35,5}{g.mol^{-1}}$. Le rendement de la synthèse est de $80\,\%$.
\begin{enumerate}
\item Écris la formule semi-développée du chlorure d'éthanoyle et entoure le groupe caractéristique des chlorures d'acyle.
\item Écris l'équation-bilan de la réaction de synthèse du paracétamol.
\item Sur la formule du paracétamol, entoure et nomme les deux fonctions organiques présentes dans la molécule.
\item À quelle classe appartient la fonction \ce{-NH2} du para-aminophénol ? Pourquoi cette fonction est-elle nucléophile ?
\item Calcule les masses molaires du para-aminophénol et du paracétamol.
\item On fait réagir $m = \SI{10,0}{g}$ de para-aminophénol avec un excès de chlorure d'éthanoyle. Calcule la masse théorique de paracétamol attendue, puis la masse réellement obtenue compte tenu du rendement.
\item Cite une précaution expérimentale importante à prendre lors de la manipulation du chlorure d'éthanoyle (on rappelle qu'il réagit violemment avec l'eau en dégageant \ce{HCl} gazeux).
\end{enumerate}
\end{exercice}

\begin{exercice}[Type Bac — Les amines autour de nous]
Les amines sont présentes dans de nombreux produits de la vie quotidienne : certaines sont responsables de l'odeur du poisson avarié (triméthylamine), d'autres entrent dans la composition de médicaments. On s'intéresse à trois amines : la méthanamine \ce{CH3NH2}, la N-méthylméthanamine \ce{(CH3)2NH} et la N,N-diméthylméthanamine (triméthylamine) \ce{(CH3)3N}.

\medskip
\textbf{Données :} p$K_A$ des couples conjugués : \ce{CH3NH3+}/\ce{CH3NH2} : $10,7$ ; \ce{(CH3)2NH2+}/\ce{(CH3)2NH} : $10,8$ ; \ce{(CH3)3NH+}/\ce{(CH3)3N} : $9,8$. Températures d'ébullition sous \SI{1}{atm} : méthanamine : $-6\,\degres$C ; N-méthylméthanamine : $7\,\degres$C ; triméthylamine : $3\,\degres$C.
\begin{enumerate}
\item Écris les formules semi-développées des trois amines et précise la classe de chacune.
\item Ces trois amines ont des masses molaires faibles et croissantes. Propose une explication, fondée sur les liaisons hydrogène intermoléculaires, à l'évolution des températures d'ébullition observées.
\item Écris l'équation-bilan de la réaction de la triméthylamine avec l'eau et identifie les couples acide/base mis en jeu.
\item Classe les trois amines par basicité croissante en justifiant ta réponse à l'aide des valeurs de p$K_A$.
\item On prépare une solution aqueuse de méthanamine de pH égal à $11,7$. 
\begin{enumerate}
\item À l'aide d'un axe des p$K_A$, détermine l'espèce prédominante du couple \ce{CH3NH3+}/\ce{CH3NH2} dans cette solution.
\item Calcule le rapport $\dfrac{[\ce{CH3NH2}]}{[\ce{CH3NH3+}]}$.
\end{enumerate}
\item La triméthylamine peut être dosée par une solution d'acide chlorhydrique. Écris l'équation-bilan de la réaction de dosage et précise ses caractéristiques (rapide, totale, etc.).
\end{enumerate}
\end{exercice}

\begin{popfigure}
\begin{tikzpicture}[scale=0.9, >=Stealth]
% Axe pKa pour méthanamine
\draw[->] (0,0) -- (10,0) node[right] {pH};
\foreach \x in {0,2,4,6,8,10} \draw (\x,0.1) -- (\x,-0.1) node[below] {\x};
\draw[thick, popblue] (0,0.5) -- (10,0.5) node[right] {Axe p$K_A$};
\draw[popblue, thick] (10.7,0.3) -- (10.7,0.7) node[above] {p$K_A$ = 10,7};
\draw[poporange, thick, <->] (0,0.8) -- (10.7,0.8) node[midway, above] {\ce{CH3NH3+} prédomine};
\draw[popgreen, thick, <->] (10.7,1.1) -- (10,1.1) node[midway, above] {\ce{CH3NH2} prédomine};
\draw[popred, dashed, thick] (11.7, -0.3) -- (11.7, 1.3) node[above] {pH = 11,7};
\end{tikzpicture}
\legende{Axe des p$K_A$ pour le couple \ce{CH3NH3+}/\ce{CH3NH2} (p$K_A$ = 10,7).}
\end{popfigure}

\begin{methode}[Tracer un axe des p$K_A$ et déterminer l'espèce majoritaire]
\begin{enumerate}
\item Dessine un axe horizontal gradué en pH (de 0 à 14).
\item Repère la valeur du p$K_A$ du couple acide/base considéré.
\item Trace une ligne verticale pointillée à cette valeur : elle sépare l'axe en deux domaines.
\item \textbf{Domaine acide (pH < p$K_A$) :} l'espèce acide (forme protonée, \ce{AH} ou \ce{BH+}) est majoritaire.
\item \textbf{Domaine basique (pH > p$K_A$) :} l'espèce basique (forme déprotonée, \ce{A-} ou \ce{B}) est majoritaire.
\item Pour calculer le rapport $\frac{[\text{Base}]}{[\text{Acide}]}$, utilise la relation : $\text{pH} = \text{p}K_A + \log\frac{[\text{Base}]}{[\text{Acide}]}$.
\end{enumerate}
\end{methode}

\begin{attention}[Pièges classiques sur les amines]
\begin{itemize}
\item \textbf{Classe de l'amine :} On compte le nombre de groupes carbonés (alkyles/aryles) \textbf{liés directement à l'atome d'azote}, pas le nombre d'atomes de carbone voisins.
\item \textbf{Liaisons hydrogène :} Seules les amines primaires ($-$NH$_2$) et secondaires ($>$NH) donnent des liaisons H \textbf{donneuses} (H sur N). Les amines tertiaires n'ont pas de H sur N : elles n'acceptent que des liaisons H (donneuses d'eau par ex.), d'où des températures d'ébullition plus basses à masse molaire égale.
\item \textbf{Basicité vs Encombrement :} L'effet inductif +I des groupes alkyles augmente la basicité (1° $<$ 2°), mais l'encombrement stérique et la moins bonne solvatation de l'ion ammonium tétra-substitué font chuter la basicité de l'amine tertiaire (3° $<$ 1° $<$ 2° en général).
\item \textbf{Réaction d'Hofmann :} C'est une substitution nucléophile ($S_\text{N}2$ pour les primaires). L'amine produite est plus nucléophile que l'ammoniac (effet +I) : la réaction continue (polyalkylation) sauf si l'ammoniac est en \textbf{grand excès}.
\item \textbf{Chlorure d'acyle :} Réaction violente, exothermique, dégage \ce{HCl} gazeux (toxique, irritant). Toujours manipuler sous hotte, avec EPI, et verser le chlorure d'acyle \textbf{lentement} dans l'amine (éventuellement en solution dans une base comme \ce{NaOH} ou pyridine pour neutraliser le \ce{HCl} formé).
\end{itemize}
\end{attention}

\newpage

\coursec{Corrigés des exercices}

\begin{corrige}[Exercice 1 — QCM : les classes d'amines]
\begin{enumerate}
\item \textbf{Réponse b)} — Une amine secondaire possède \cle{deux groupes alkyles (ou aryles)} liés à l'atome d'azote : $R-\text{NH}-R'$. La réponse a) décrit une amine primaire $R-\text{NH}_2$, la réponse c) une amine tertiaire $R-\text{N}(R')-R''$.
\item \textbf{Réponse b)} — Une amine aliphatique saturée et acyclique a pour formule générale $C_nH_{2n+3}N$ (on ajoute \ce{NH} à un alcane $C_nH_{2n+2}$). Vérification : \ce{CH3NH2} donne $C_1H_5N$, et $2\times 1+3 = 5$. \checkmark
\item \textbf{Réponse a)} — \chemfig{CH_3-NH-CH_2-CH_3} se nomme \cle{N-méthyléthanamine} : la chaîne la plus longue liée à l'azote (éthyle, 2 carbones) donne le nom de base \emph{éthanamine}, et le groupe méthyle porté par l'azote est cité en préfixe \emph{N-méthyl}. La réponse b) est le nom usuel (non systématique) ; la réponse c) est incorrecte car la chaîne principale doit être la plus longue.
\item \textbf{Réponse a)} — Le caractère basique des amines est dû au \cle{doublet non liant} de l'atome d'azote, capable de capter un proton \ce{H+} pour former l'ion alkylammonium.
\end{enumerate}
\end{corrige}

\begin{corrige}[Exercice 2 — Vrai ou faux ? Justifie]
\begin{enumerate}
\item \textbf{Faux.} Seules les amines \textbf{primaires} (\ce{-NH2}) et \textbf{secondaires} (\ce{R2NH}), qui possèdent au moins un atome d'hydrogène lié à l'azote, peuvent former des liaisons hydrogène entre leurs propres molécules. Les amines tertiaires (\ce{R3N}) n'ont pas de H sur N : elles ne peuvent pas en établir entre elles.
\item \textbf{Faux.} La réaction avec l'eau met en jeu le \emph{doublet non liant} de l'azote, qui capte un proton de l'eau ; elle ne nécessite pas d'atome H sur l'azote. Une amine tertiaire réagit bien avec l'eau : \[ \ce{R3N + H2O <=> R3NH+ + HO-} \]
\item \textbf{Vrai.} Dans la réaction d'Hofmann, l'amine primaire \ce{R-NH2} attaque le dérivé halogéné \ce{R1-X} par son doublet : \[ \ce{R-NH2 + R1-X -> R-NH-R1 + HX} \]
\item \textbf{Vrai.} L'action d'une amine sur un chlorure d'acyle est une substitution nucléophile sur le carbone du groupe \ce{-COCl} ; elle conduit à un \cle{amide} avec dégagement de \ce{HCl} : \[ \ce{R-COCl + R1-NH2 -> R-CO-NH-R1 + HCl} \]
\item \textbf{Vrai.} Les amines de petite taille (jusqu'à 5 ou 6 atomes de carbone) sont très solubles dans l'eau car elles établissent des liaisons hydrogène avec les molécules d'eau. Au-delà, la longue chaîne carbonée hydrophobe l'emporte et la solubilité chute.
\end{enumerate}
\end{corrige}

\begin{corrige}[Exercice 3 — Texte à trous]
La molécule d'ammoniac \ce{NH3} a une géométrie en \textbf{pyramide} à base triangulaire ; l'atome d'azote porte un \textbf{doublet non liant} qui confère aux amines un caractère \textbf{basique} et un caractère \textbf{nucléophile}. Les amines primaires et secondaires établissent entre leurs molécules des \textbf{liaisons hydrogène}, ce qui explique leurs températures d'ébullition relativement élevées. L'action d'une amine sur un chlorure d'acyle conduit à un \textbf{amide}.
\medskip

\emph{Mots non utilisés :} électrophile, plane, acide, liaisons covalentes. \emph{Point de vigilance :} c'est bien le doublet non liant qui est à l'origine des deux caractères (basique \emph{et} nucléophile) ; une amine n'est jamais électrophile (au niveau de l'atome d'azote).
\end{corrige}

\begin{corrige}[Exercice 4 — Définitions essentielles]
\begin{enumerate}
\item Une \cle{amine} est un composé organique qui dérive de l'ammoniac \ce{NH3} par remplacement d'un, deux ou trois atomes d'hydrogène par des groupes alkyles (ou aryles). Pour une amine aliphatique saturée et acyclique, la formule générale est $C_nH_{2n+3}N$.
\item Les trois classes se distinguent par le nombre de groupes carbonés liés à l'azote :
\begin{itemize}
\item \textbf{Amine primaire} : un seul groupe alkyle (ou aryle) lié à N, formule $R-\mathrm{NH}_2$ ;
\item \textbf{Amine secondaire} : deux groupes liés à N, formule $R-\mathrm{NH}-R'$ ;
\item \textbf{Amine tertiaire} : trois groupes liés à N, formule $R-\mathrm{N}(-R')-R''$.
\end{itemize}
\item La \cle{réaction d'Hofmann} est l'action d'une amine (ou de l'ammoniac) sur un dérivé halogéné \ce{R-X} : c'est une substitution nucléophile au cours de laquelle l'amine attaque le carbone porteur de l'halogène, avec départ de l'ion halogénure, pour donner une amine de classe supérieure (ou un sel d'ammonium quaternaire en fin de série).
\item Le \cle{caractère nucléophile} d'une amine est sa tendance à attaquer les sites pauvres en électrons (déficitaires en électrons, dits électrophiles) d'autres molécules. Il est dû au \cle{doublet non liant} de l'atome d'azote, riche en électrons et disponible pour former une nouvelle liaison covalente.
\end{enumerate}
\end{corrige}

\begin{corrige}[Exercice 5 — Classes et nomenclature]
\begin{enumerate}
\item \textbf{Classe de chaque amine} (on compte les groupes carbonés liés à N) :
\begin{center}
\begin{tabularx}{\linewidth}{|l|X|l|}
\hline
\rowcolor{popblueL} \textbf{Amine} & \textbf{Formule semi-développée} & \textbf{Classe} \\
\hline
A & \chemfig{CH_3-CH_2-CH_2-NH_2} & primaire (1 groupe sur N) \\
\hline
B & \chemfig{CH_3-CH(-[2]NH_2)-CH_3} & primaire (1 groupe sur N) \\
\hline
C & \chemfig{CH_3-NH-CH_2-CH_3} & secondaire (2 groupes sur N) \\
\hline
D & \chemfig{CH_3-N(-[2]CH_3)-CH_2-CH_3} & tertiaire (3 groupes sur N) \\
\hline
E & \chemfig{CH_3-CH_2-NH-CH_2-CH_3} & secondaire (2 groupes sur N) \\
\hline
\end{tabularx}
\end{center}
\item \textbf{Nomenclature IUPAC} (la chaîne la plus longue liée à N donne le nom de base ; les autres groupes sur N sont cités en préfixe N-) :
\begin{itemize}
\item A : \textbf{propan-1-amine} ;
\item B : \textbf{propan-2-amine} ;
\item C : \textbf{N-méthyléthanamine} ;
\item D : \textbf{N,N-diméthyléthanamine} ;
\item E : \textbf{N-éthyléthanamine} (nom usuel : diéthylamine).
\end{itemize}
\item \textbf{Isomères :} comptons les atomes de chaque amine.
\begin{itemize}
\item A et B : $C_3H_9N$ — ce sont des isomères de \emph{position} (fonction \ce{-NH2} en position 1 ou 2) ;
\item C : $C_3H_9N$ également ! (\ce{CH3} + \ce{NH} + \ce{C2H5} = 3 C, 9 H, 1 N) ;
\item D : $C_4H_{11}N$ (2 méthyles + 1 éthyle = 4 C) ;
\item E : $C_4H_{11}N$ (2 éthyles = 4 C).
\end{itemize}
Donc \textbf{A, B et C sont isomères} ($\ce{C3H9N}$) et \textbf{D et E sont isomères} ($\ce{C4H11N}$). C'est une belle illustration du fait que des amines de classes différentes peuvent être isomères.
\end{enumerate}
\end{corrige}

\begin{corrige}[Exercice 6 — Isomères de formule \ce{C3H9N}]
\begin{enumerate}
\item \textbf{Formules semi-développées} : on explore les trois classes.
\begin{itemize}
\item Amines primaires (1 groupe $C_3$ sur N, deux chaînes possibles) :

\chemfig{CH_3-CH_2-CH_2-NH_2} \hspace{2cm} \chemfig{CH_3-CH(-[2]NH_2)-CH_3}
\item Amine secondaire (2 groupes sur N : méthyle + éthyle) :

\chemfig{CH_3-NH-CH_2-CH_3}
\item Amine tertiaire (3 méthyles sur N) :

\chemfig{CH_3-N(-[2]CH_3)-CH_3}
\end{itemize}
Il y a donc \textbf{4 isomères} au total.
\item \textbf{Noms et classes :}
\begin{center}
\begin{tabularx}{\linewidth}{|X|l|l|}
\hline
\rowcolor{popblueL} \textbf{Formule} & \textbf{Nom IUPAC} & \textbf{Classe} \\
\hline
\chemfig{CH_3-CH_2-CH_2-NH_2} & propan-1-amine & primaire \\
\hline
\chemfig{CH_3-CH(-[2]NH_2)-CH_3} & propan-2-amine & primaire \\
\hline
\chemfig{CH_3-NH-CH_2-CH_3} & N-méthyléthanamine & secondaire \\
\hline
\chemfig{CH_3-N(-[2]CH_3)-CH_3} & N,N-diméthylméthanamine & tertiaire \\
\hline
\end{tabularx}
\end{center}
\item \textbf{Températures d'ébullition :} à masse molaire égale ($M = \SI{59}{g.mol^{-1}}$ pour tous), ce sont les liaisons hydrogène intermoléculaires qui font la différence.
\begin{itemize}
\item Les amines \textbf{primaires} (\ce{-NH2}) possèdent \emph{deux} atomes H sur N : elles forment le plus de liaisons hydrogène entre leurs molécules $\Rightarrow$ températures d'ébullition \textbf{les plus élevées} (propan-1-amine : $\approx \SI{48}{\celsius}$).
\item L'amine \textbf{secondaire} (\ce{-NH-}) n'a qu'\emph{un} H sur N : liaisons hydrogène moins nombreuses $\Rightarrow$ ébullition intermédiaire ($\approx \SI{36}{\celsius}$).
\item L'amine \textbf{tertiaire} (\ce{N(CH3)3}) n'a \emph{aucun} H sur N : pas de liaison hydrogène entre ses propres molécules $\Rightarrow$ température d'ébullition \textbf{la plus faible} ($\approx \SI{3}{\celsius}$).
\end{itemize}
\end{enumerate}
\end{corrige}

\begin{corrige}[Exercice 7 — Une amine dans l'eau]
\begin{enumerate}
\item \textbf{Équation-bilan} (réaction limitée, d'où la double flèche) : \[ \ce{C3H7NH2 + H2O <=> C3H7NH3+ + HO-} \]
\item \textbf{Couples acide/base mis en jeu :}
\begin{itemize}
\item \ce{C3H7NH3+}/\ce{C3H7NH2} (l'amine joue le rôle de base, elle capte \ce{H+}) ;
\item \ce{H2O}/\ce{HO-} (l'eau joue ici le rôle d'acide, elle cède \ce{H+}) .
\end{itemize}
\item La réaction produit des ions \textbf{hydroxyde \ce{HO-}}. La concentration $[\ce{HO-}]$ augmente, donc $[\ce{H3O+}]$ diminue (produit ionique de l'eau) : la solution est basique, d'où un pH supérieur à 7 (ici $11,8$). Cohérent !
\item \begin{enumerate}
\item \textbf{Axe des p$K_A$} (voir la figure de l'énoncé) : on place p$K_A = 10,6$.
\begin{itemize}
\item Si pH $< 10,6$ : l'espèce acide \ce{C3H7NH3+} prédomine ;
\item Si pH $> 10,6$ : l'espèce basique \ce{C3H7NH2} prédomine.
\end{itemize}
\item Ici pH $= 11,8 > 10,6$ : c'est donc l'amine libre \textbf{\ce{C3H7NH2}} qui prédomine.
\item On applique la relation de Henderson-Hasselbalch :
\[ \text{pH} = \text{p}K_A + \log\frac{[\ce{C3H7NH2}]}{[\ce{C3H7NH3+}]} \]
\[ \log\frac{[\ce{C3H7NH2}]}{[\ce{C3H7NH3+}]} = 11,8 - 10,6 = 1,2 \]
\[ \frac{[\ce{C3H7NH2}]}{[\ce{C3H7NH3+}]} = 10^{1,2} \approx 15,8 \]
Il y a environ \textbf{16 fois plus} d'amine moléculaire que d'ion propylammonium : la prédominance de \ce{C3H7NH2} est confirmée quantitativement.
\end{enumerate}
\end{enumerate}
\end{corrige}

\begin{corrige}[Exercice 8 — Solubilité et liaisons hydrogène]
\begin{enumerate}
\item \textbf{Formules et classes :}
\begin{itemize}
\item Butan-1-amine : \chemfig{CH_3-CH_2-CH_2-CH_2-NH_2} — amine \textbf{primaire} ;
\item N,N-diméthyléthanamine : \chemfig{CH_3-CH_2-N(-[2]CH_3)-CH_3} — amine \textbf{tertiaire}.
\end{itemize}
Ce sont deux isomères de formule brute \ce{C4H11N}.
\item À masse molaire égale :
\begin{itemize}
\item la \textbf{butan-1-amine} (primaire) a la température d'ébullition \textbf{la plus élevée} ($\approx \SI{78}{\celsius}$ contre $\approx \SI{37}{\celsius}$ pour la tertiaire) ;
\item les deux sont solubles dans l'eau (4 C seulement), mais la \textbf{butan-1-amine} est \textbf{la plus soluble}.
\end{itemize}
\item \textbf{Justification :}
\begin{itemize}
\item \emph{État pur (ébullition) :} la butan-1-amine possède deux H sur N : ses molécules s'associent par de nombreuses liaisons hydrogène $N-H\cdots N$, qu'il faut rompre pour vaporiser le liquide $\Rightarrow$ ébullition élevée. L'amine tertiaire n'a aucun H sur N : pas de liaison hydrogène entre ses molécules, seules des interactions de Van der Waals, plus faibles $\Rightarrow$ ébullition basse.
\item \emph{Solubilité dans l'eau :} les deux amines peuvent \emph{accepter} des liaisons hydrogène de l'eau (via le doublet de N), mais la butan-1-amine peut en plus en \emph{donner} (via ses H sur N) : elle forme davantage de liaisons hydrogène $N-H\cdots O$ et $O-H\cdots N$ avec l'eau, d'où sa meilleure solubilité.
\end{itemize}
\end{enumerate}
\end{corrige}

\begin{corrige}[Exercice 9 — Type Bac : la réaction d'Hofmann au laboratoire]
\begin{enumerate}
\item Une \textbf{amine} est un composé dérivé de l'ammoniac \ce{NH3} par remplacement d'un, deux ou trois atomes d'hydrogène par des groupes alkyles (ou aryles). Formule générale d'une amine aliphatique saturée acyclique : $C_nH_{2n+3}N$.
\item L'atome d'azote de \ce{NH3} et des amines porte un \textbf{doublet non liant}, riche en électrons : il peut attaquer un atome de carbone déficitaire en électrons (comme le carbone porteur du brome dans \ce{R-Br}, rendu $\delta+$ par la polarité de la liaison \ce{C-Br}). C'est le \textbf{caractère nucléophile}.
\item \textbf{Première réaction} (substitution nucléophile) : \[ \ce{NH3 + CH3-CH2-CH2-Br -> CH3-CH2-CH2-NH2 + HBr} \] Le produit organique est la \textbf{propan-1-amine}, amine \textbf{primaire}. (Remarque : en présence d'excès de \ce{NH3}, le \ce{HBr} formé est neutralisé : \ce{HBr + NH3 -> NH4+ + Br-}.)
\item \textbf{Réactions successives :} la propan-1-amine, plus nucléophile que \ce{NH3} (effet inductif donneur du groupe propyle), attaque à son tour le 1-bromopropane :
\[ \ce{CH3CH2CH2-NH2 + CH3CH2CH2-Br -> (CH3CH2CH2)2NH + HBr} \]
Produit : la \textbf{N-propylpropan-1-amine} (dipropylamine), amine \textbf{secondaire}. Puis :
\[ \ce{(CH3CH2CH2)2NH + CH3CH2CH2-Br -> (CH3CH2CH2)3N + HBr} \]
Produit : la \textbf{N,N-dipropylpropan-1-amine} (tripropylamine), amine \textbf{tertiaire}. (Une quatrième étape peut même conduire au sel d'ammonium quaternaire \ce{(C3H7)4N+ Br-}.)
\item L'ammoniac en \textbf{grand excès} augmente la probabilité que le 1-bromopropane rencontre une molécule de \ce{NH3} plutôt qu'une molécule de propan-1-amine déjà formée : on favorise ainsi la première étape et on limite la polyalkylation. La propan-1-amine est alors le produit \textbf{majoritaire}.
\item \textbf{Application numérique :}
\begin{enumerate}
\item Masse molaire du 1-bromopropane \ce{C3H7Br} :
\[ M = 3\times 12 + 7\times 1 + 80 = 36 + 7 + 80 = \SI{123}{g.mol^{-1}} \]
Quantité de 1-bromopropane :
\[ n(\ce{C3H7Br}) = \frac{m}{M} = \frac{12,3}{123} = \SI{0,10}{mol} \]
L'équation-bilan se fait mole à mole (stœchiométrie 1:1). On a $n(\ce{NH3}) = \SI{0,10}{mol}$ et $n(\ce{C3H7Br}) = \SI{0,10}{mol}$ : les deux réactifs sont en \textbf{proportions stœchiométriques} ; il n'y a pas de réactif limitant (ou, si l'on veut, les deux s'épuisent en même temps).
\medskip

\emph{Remarque :} dans la pratique, on mettrait \ce{NH3} en excès (question 5) ; ici l'énoncé impose les quantités égales.
\item D'après la stœchiométrie, $n_{\text{max}}(\ce{C3H7NH2}) = \SI{0,10}{mol}$. Masse molaire de la propan-1-amine \ce{C3H9N} :
\[ M = 3\times 12 + 9\times 1 + 14 = 36 + 9 + 14 = \SI{59}{g.mol^{-1}} \]
\[ m_{\text{max}} = n \times M = 0,10 \times 59 = \boxed{\SI{5,9}{g}} \]
\end{enumerate}
\end{enumerate}
\medskip
\textbf{Barème indicatif (sur 10) :} définition + formule générale (1 pt) ; caractère nucléophile justifié (1 pt) ; équation 1 + nom + classe (2 pts) ; deux équations successives + noms + classes (2 pts) ; rôle de l'excès (1 pt) ; calcul de $M$ et $n$ (1,5 pt) ; masse théorique avec unité (1,5 pt).
\medskip

\textbf{Points de vigilance :} ne pas oublier \ce{HBr} dans les équations ; bien compter les groupes sur N pour la classe ; vérifier la stœchiométrie 1:1 avant de conclure sur le réactif limitant ; toujours donner la masse avec son unité.
\end{corrige}

\begin{corrige}[Exercice 10 — Type Bac : synthèse du paracétamol]
\begin{enumerate}
\item \textbf{Chlorure d'éthanoyle} : \chemfig{CH_3-C(=[1]O)-[7]Cl}. Le groupe caractéristique des chlorures d'acyle est \ce{-C(=O)-Cl} (un carbonyle portant un chlore) : c'est lui qu'il faut entourer.
\item \textbf{Équation-bilan :} \[ \ce{HO-C6H4-NH2 + CH3-COCl -> HO-C6H4-NH-CO-CH3 + HCl} \]
\item Sur la formule du paracétamol \ce{HO-C6H4-NH-CO-CH3}, on entoure :
\begin{itemize}
\item le groupe \textbf{\ce{-OH}} fixé sur le cycle benzénique : fonction \textbf{phénol} ;
\item le groupe \textbf{\ce{-NH-CO-}} : fonction \textbf{amide} (plus précisément amide substitué, ou anilide).
\end{itemize}
\item Le groupe \ce{-NH2} du para-aminophénol appartient à la fonction \textbf{amine primaire} (un seul groupe carboné — le cycle — lié à N). Elle est nucléophile grâce au \textbf{doublet non liant} de l'azote, qui attaque le carbone électrophile ($\delta+$) du groupe \ce{-COCl}.
\item \textbf{Masses molaires :}
\begin{itemize}
\item Para-aminophénol \ce{HO-C6H4-NH2} = \ce{C6H7NO} :
\[ M = 6\times 12 + 7\times 1 + 14 + 16 = 72 + 7 + 14 + 16 = \SI{109}{g.mol^{-1}} \]
\item Paracétamol \ce{HO-C6H4-NH-CO-CH3} = \ce{C8H9NO2} :
\[ M = 8\times 12 + 9\times 1 + 14 + 2\times 16 = 96 + 9 + 14 + 32 = \SI{151}{g.mol^{-1}} \]
\end{itemize}
\item \textbf{Calcul des masses :} le chlorure d'éthanoyle est en excès, donc le para-aminophénol est le réactif limitant.
\[ n(\text{para-aminophénol}) = \frac{m}{M} = \frac{10,0}{109} \approx \SI{9,17e-2}{mol} \]
Stœchiométrie 1:1 : $n_{\text{th}}(\text{paracétamol}) = \SI{9,17e-2}{mol}$.
\[ m_{\text{th}} = n \times M = \frac{10,0}{109} \times 151 \approx \SI{13,9}{g} \]
Avec un rendement de $80\,\%$ :
\[ m_{\text{réelle}} = 0,80 \times 13,9 \approx \boxed{\SI{11,1}{g}} \]
\item \textbf{Précaution :} le chlorure d'éthanoyle réagit violemment avec l'eau en dégageant \ce{HCl} gazeux (toxique et irritant). Il faut donc \textbf{manipuler sous hotte ventilée}, avec blouse, gants et lunettes, utiliser de la \textbf{verrerie bien sèche}, et verser le chlorure d'acyle \textbf{lentement} (goutte à goutte) dans le milieu réactionnel refroidi si nécessaire.
\end{enumerate}
\medskip
\textbf{Barème indicatif (sur 10) :} formule du chlorure d'acyle + groupe entouré (1 pt) ; équation-bilan équilibrée (1,5 pt) ; deux fonctions identifiées et nommées (1,5 pt) ; classe + nucléophilie (1 pt) ; masses molaires (1 pt) ; quantité de matière et masse théorique (2 pts) ; masse réelle avec rendement (1 pt) ; précaution (1 pt).
\medskip

\textbf{Points de vigilance :} ne pas oublier \ce{HCl} dans l'équation ; bien compter les atomes pour les formules brutes (\ce{C6H7NO} et \ce{C8H9NO2}) ; le rendement s'applique à la masse \emph{théorique}, jamais l'inverse ; la fonction \ce{-OH} sur un cycle benzénique est un phénol, pas un alcool.
\end{corrige}

\begin{corrige}[Exercice 11 — Type Bac : les amines autour de nous]
\begin{enumerate}
\item \textbf{Formules et classes :}
\begin{itemize}
\item Méthanamine : \chemfig{CH_3-NH_2} — amine \textbf{primaire} ;
\item N-méthylméthanamine : \chemfig{CH_3-NH-CH_3} — amine \textbf{secondaire} ;
\item Triméthylamine : \chemfig{CH_3-N(-[2]CH_3)-CH_3} — amine \textbf{tertiaire}.
\end{itemize}
\item \textbf{Évolution des températures d'ébullition} ($-6\,^{\circ}$C ; $+7\,^{\circ}$C ; $+3\,^{\circ}$C) :
\begin{itemize}
\item La méthanamine (primaire, \ce{-NH2}) possède deux H sur N : nombreuses liaisons hydrogène intermoléculaires. Mais sa masse molaire ($M = \SI{31}{g.mol^{-1}}$) est la plus faible, donc les interactions de Van der Waals sont faibles : ébullition la plus basse ($-6\,^{\circ}$C).
\item La N-méthylméthanamine (secondaire, $>NH$, $M = \SI{45}{g.mol^{-1}}$) garde un H donneur de liaison hydrogène \emph{et} a une masse plus grande : le cumul liaisons H + Van der Waals donne l'ébullition \textbf{la plus élevée} ($7\,^{\circ}$C).
\item La triméthylamine (tertiaire, $M = \SI{59}{g.mol^{-1}}$) n'a \textbf{aucun H sur N} : pas de liaison hydrogène entre ses molécules ; malgré sa masse plus élevée, son ébullition ($3\,^{\circ}$C) retombe sous celle de l'amine secondaire.
\end{itemize}
Conclusion : les liaisons hydrogène intermoléculaires (présentes pour 1° et 2°, absentes pour 3°) expliquent que l'amine secondaire bout plus haut que la tertiaire, alors même qu'elle est plus légère.
\item \textbf{Réaction de la triméthylamine avec l'eau :} \[ \ce{(CH3)3N + H2O <=> (CH3)3NH+ + HO-} \] Couples mis en jeu : \ce{(CH3)3NH+}/\ce{(CH3)3N} et \ce{H2O}/\ce{HO-}.
\item \textbf{Classement par basicité croissante :} plus le p$K_A$ du couple conjugué est élevé, plus la base est forte.
\[ \underbrace{\ce{(CH3)3N}}_{\text{p}K_A = 9,8} \;<\; \underbrace{\ce{CH3NH2}}_{\text{p}K_A = 10,7} \;<\; \underbrace{\ce{(CH3)2NH}}_{\text{p}K_A = 10,8} \]
L'effet inductif donneur des groupes méthyles enrichit le doublet de N (la secondaire est la plus basique), mais pour la tertiaire, l'encombrement stérique et la moins bonne solvatation de l'ion \ce{(CH3)3NH+} font chuter la basicité.
\item \begin{enumerate}
\item pH $= 11,7 > \text{p}K_A = 10,7$ : on est dans le domaine basique de l'axe (voir figure) : c'est l'espèce basique \textbf{\ce{CH3NH2}} qui prédomine.
\item Relation de Henderson-Hasselbalch :
\[ \log\frac{[\ce{CH3NH2}]}{[\ce{CH3NH3+}]} = \text{pH} - \text{p}K_A = 11,7 - 10,7 = 1,0 \]
\[ \frac{[\ce{CH3NH2}]}{[\ce{CH3NH3+}]} = 10^{1,0} = \boxed{10} \]
Il y a 10 fois plus de méthanamine moléculaire que d'ion méthylammonium.
\end{enumerate}
\item \textbf{Dosage de la triméthylamine par \ce{HCl} :} \[ \ce{(CH3)3N + H3O+ -> (CH3)3NH+ + H2O} \] Cette réaction acide-base est \textbf{rapide}, \textbf{totale} (constante d'équilibre $K = 10^{\text{p}K_A} = 10^{9,8} \gg 10^4$) et \textbf{unique} : elle convient parfaitement à un dosage direct. À l'équivalence, la solution contient le chlorure de triméthylammonium \ce{(CH3)3NH+ Cl-}, de pH légèrement acide (acide faible de p$K_A = 9,8$).
\end{enumerate}
\medskip
\textbf{Barème indicatif (sur 10) :} formules + classes (1,5 pt) ; explication des ébullitions par les liaisons H (2 pts) ; équation avec l'eau + couples (1,5 pt) ; classement de basicité justifié (1,5 pt) ; espèce prédominante (0,5 pt) ; calcul du rapport (1 pt) ; équation de dosage + caractéristiques (2 pts).
\medskip

\textbf{Points de vigilance :} une amine tertiaire réagit bien avec l'eau (c'est le doublet qui compte, pas le H sur N) ; ne pas conclure « plus lourd = bout plus haut » sans invoquer les liaisons hydrogène ; dans la relation de Henderson, vérifier que la base est au numérateur ; une réaction de dosage doit être rapide, totale et unique — cite les trois adjectifs.
\end{corrige}

\newpage

\coursec{Fiche bilan}

\begin{popfigure}
\begin{center}\tcbox[colback=popdark,colframe=popdark,coltext=white,arc=9pt,boxsep=4pt,left=6pt,right=6pt]{\bfseries\small Les amines}\end{center}
\vspace{-2pt}
\begin{tcbraster}[raster columns=3,raster equal height=rows,raster column skip=6pt,raster row skip=6pt,enhanced,blankest]
\begin{tcolorbox}[enhanced,colback=white,colframe=popblue,boxrule=0.9pt,arc=3pt,title={Structure et classes},fonttitle=\bfseries\scriptsize,coltitle=white,colbacktitle=popblue,left=4pt,right=4pt,top=2pt,bottom=2pt,boxsep=1pt,toptitle=1.5pt,bottomtitle=1.5pt,halign=left]
\begin{itemize}[nosep,leftmargin=*,itemsep=1pt,font=\scriptsize]
\item primaire, secondaire, tertiaire
\item doublet sur N
\end{itemize}
\end{tcolorbox}
\begin{tcolorbox}[enhanced,colback=white,colframe=poporange,boxrule=0.9pt,arc=3pt,title={Nomenclature},fonttitle=\bfseries\scriptsize,coltitle=white,colbacktitle=poporange,left=4pt,right=4pt,top=2pt,bottom=2pt,boxsep=1pt,toptitle=1.5pt,bottomtitle=1.5pt,halign=left]
\begin{itemize}[nosep,leftmargin=*,itemsep=1pt,font=\scriptsize]
\item alcanamine
\item N-alkyl
\end{itemize}
\end{tcolorbox}
\begin{tcolorbox}[enhanced,colback=white,colframe=popgreen,boxrule=0.9pt,arc=3pt,title={Propri\'et\'es physiques},fonttitle=\bfseries\scriptsize,coltitle=white,colbacktitle=popgreen,left=4pt,right=4pt,top=2pt,bottom=2pt,boxsep=1pt,toptitle=1.5pt,bottomtitle=1.5pt,halign=left]
\begin{itemize}[nosep,leftmargin=*,itemsep=1pt,font=\scriptsize]
\item liaisons H
\item solubilit\'e
\end{itemize}
\end{tcolorbox}
\begin{tcolorbox}[enhanced,colback=white,colframe=poppurple,boxrule=0.9pt,arc=3pt,title={Basicit\'e},fonttitle=\bfseries\scriptsize,coltitle=white,colbacktitle=poppurple,left=4pt,right=4pt,top=2pt,bottom=2pt,boxsep=1pt,toptitle=1.5pt,bottomtitle=1.5pt,halign=left]
\begin{itemize}[nosep,leftmargin=*,itemsep=1pt,font=\scriptsize]
\item \ce{RNH2 + H2O <=> RNH3+ + HO-}
\end{itemize}
\end{tcolorbox}
\begin{tcolorbox}[enhanced,colback=white,colframe=poppink,boxrule=0.9pt,arc=3pt,title={Hofmann},fonttitle=\bfseries\scriptsize,coltitle=white,colbacktitle=poppink,left=4pt,right=4pt,top=2pt,bottom=2pt,boxsep=1pt,toptitle=1.5pt,bottomtitle=1.5pt,halign=left]
\begin{itemize}[nosep,leftmargin=*,itemsep=1pt,font=\scriptsize]
\item amine + \ce{R-X}
\item alkylations successives
\end{itemize}
\end{tcolorbox}
\begin{tcolorbox}[enhanced,colback=white,colframe=popgold,boxrule=0.9pt,arc=3pt,title={Amides},fonttitle=\bfseries\scriptsize,coltitle=white,colbacktitle=popgold,left=4pt,right=4pt,top=2pt,bottom=2pt,boxsep=1pt,toptitle=1.5pt,bottomtitle=1.5pt,halign=left]
\begin{itemize}[nosep,leftmargin=*,itemsep=1pt,font=\scriptsize]
\item amine + \ce{R-COCl}
\end{itemize}
\end{tcolorbox}
\end{tcbraster}

\legende{Carte mentale de la le\c{c}on : les amines, de l'ammoniac aux amides.}
\end{popfigure}

\begin{bilan}
\textbf{D\'efinitions cl\'es}
\begin{itemize}
  \item Une \cle{amine} est un compos\'e organique d\'eriv\'e de l'ammoniac \ce{NH3} par remplacement de un, deux ou trois atomes d'hydrog\`ene par des groupes alkyles (ou aryles).
  \item Formule g\'en\'erale : $C_nH_{2n+3}N$ pour une amine satur\'ee acyclique.
  \item \cle{Classe} d'une amine : nombre de groupes carbon\'es li\'es \`a l'atome d'azote : \textbf{primaire} $R-NH_2$, \textbf{secondaire} $R-NH-R'$, \textbf{tertiaire} $R-N(R')-R''$.
  \item L'atome d'azote porte un \cle{doublet non liant} : c'est l'origine du caract\`ere \textbf{basique} (accepteur de proton) et \textbf{nucl\'eophile} (donneur de doublet) des amines.
\end{itemize}

\textbf{Formules et \'equations-types \`a conna\^itre par c\oe ur}
\begin{center}
\begin{tabularx}{\linewidth}{|l|X|}
\hline
\rowcolor{popblueL} \textbf{Notion} & \textbf{Formule / \'equation} \\
\hline
Amine primaire & $R-NH_2$ \quad (ex. \ce{CH3CH2NH2} : \'ethanamine) \\
\hline
Amine secondaire & $R-NH-R'$ \quad (ex. \ce{CH3NHCH3} : N-m\'ethylm\'ethanamine) \\
\hline
Amine tertiaire & $R-N(-R')-R''$ \quad (ex. \ce{N(CH3)3} : N,N-dim\'ethylm\'ethanamine) \\
\hline
R\'eaction avec l'eau & $R-NH_2 + \ce{H2O} \rightleftharpoons R-NH_3^+ + HO^-$ \\
\hline
R\'eaction d'Hofmann (1\iere{} \'etape) & $R-NH_2 + R'-X \rightarrow R-NH_2^+-R' + X^-$ puis $R-NH-R'$ \\
\hline
Alkylation totale & $R-NH_2 + 3 R'-X \rightarrow R-N(R')_3 + \dots$ jusqu'\`a l'ion \cle{ammonium quaternaire} $R-N(R')_3^+$ \\
\hline
Synth\`ese d'un amide & $R-COCl + R'-NH_2 \rightarrow R-CO-NH-R' + \ce{HCl}$ \\
\hline
\end{tabularx}
\end{center}

\textbf{Propri\'et\'es physiques}
\begin{itemize}
  \item Les amines primaires et secondaires forment des \cle{liaisons hydrog\`ene} intermol\'eculaires (liaison \ce{N-H}) : temp\'eratures d'\'ebullition plus \'elev\'ees que les alcanes de masse voisine, mais plus faibles que les alcools (liaison \ce{N-H} moins polaire que \ce{O-H}).
  \item Les amines tertiaires (pas de \ce{N-H}) ne s'associent pas entre elles : elles bouillent plus bas que leurs isom\`eres primaires et secondaires.
  \item Les amines \`a courte cha\^ine (moins de 6 atomes de C environ) sont \textbf{solubles dans l'eau} gr\^ace aux liaisons hydrog\`ene avec les mol\'ecules d'eau ; la solubilit\'e diminue quand la cha\^ine carbon\'ee s'allonge.
\end{itemize}

\textbf{M\'ethodes express}
\begin{itemize}
  \item \textbf{Nommer une amine} : cha\^ine la plus longue contenant N $\to$ \emph{alcan} + suffixe \emph{-amine}, avec l'indice de position de N ; les groupes port\'es par N sont cit\'es en pr\'efixe avec la lettre \emph{N}.
  \item \textbf{Identifier la classe} : compter le nombre de liaisons \ce{C-N} autour de l'azote (1, 2 ou 3).
  \item \textbf{R\'eaction d'Hofmann} : l'amine attaque le carbone porteur de l'halog\`ene (substitution nucl\'eophile) ; l'alkylation se poursuit jusqu'\`a l'ammonium quaternaire si le d\'eriv\'e halog\'en\'e est en exc\`es.
  \item \textbf{Chlorure d'acyle} : rupture de la liaison \ce{C-Cl}, d\'egagement de \ce{HCl}, formation de la liaison amide \ce{-CO-NH-} (r\'eaction rapide et totale).
\end{itemize}
\end{bilan}

\begin{aretenir}[Checklist avant le Bac]
\begin{itemize}
  \item[$\square$] Je sais d\'efinir une amine et donner sa formule g\'en\'erale $C_nH_{2n+3}N$.
  \item[$\square$] Je distingue sans h\'esiter les trois classes d'amines (primaire, secondaire, tertiaire).
  \item[$\square$] Je sais nommer une amine en nomenclature syst\'ematique (suffixe \emph{-amine}, pr\'efixe \emph{N-}).
  \item[$\square$] Je sais dessiner les formules semi-d\'evelopp\'ees d'amines isom\`eres de classes diff\'erentes.
  \item[$\square$] J'explique la solubilit\'e et les temp\'eratures d'\'ebullition par les liaisons hydrog\`ene.
  \item[$\square$] Je justifie la basicit\'e des amines par le doublet non liant de l'azote.
  \item[$\square$] J'\'ecris l'\'equation \ce{R-NH2 + H2O <=> R-NH3+ + HO-} et j'identifie le couple \ce{RNH3+}/\ce{RNH2}.
  \item[$\square$] Je sais que l'amine est nucl\'eophile et qu'elle attaque le carbone porteur de l'halog\`ene.
  \item[$\square$] J'\'ecris les \'etapes de la r\'eaction d'Hofmann jusqu'\`a l'ion ammonium quaternaire.
  \item[$\square$] J'\'ecris l'\'equation de formation d'un amide \`a partir d'un chlorure d'acyle, avec d\'egagement de \ce{HCl}.
  \item[$\square$] Je v\'erifie toujours la conservation des atomes et des charges dans mes \'equations-bilans.
  \item[$\square$] Je connais une application concr\`ete : synth\`ese de m\'edicaments (parac\'etamol), colorants, savons.
\end{itemize}
\end{aretenir}

\findecours

\end{document}
